% Options for packages loaded elsewhere
% Options for packages loaded elsewhere
\PassOptionsToPackage{unicode}{hyperref}
\PassOptionsToPackage{hyphens}{url}
\PassOptionsToPackage{dvipsnames,svgnames,x11names}{xcolor}
%
\documentclass[
  ngerman,
  a4paper,
  DIV=11,
  oneside]{scrreprt}
\usepackage{xcolor}
\usepackage[margin=2.5cm]{geometry}
\usepackage{amsmath,amssymb}
\setcounter{secnumdepth}{5}
\usepackage{iftex}
\ifPDFTeX
  \usepackage[T1]{fontenc}
  \usepackage[utf8]{inputenc}
  \usepackage{textcomp} % provide euro and other symbols
\else % if luatex or xetex
  \usepackage{unicode-math} % this also loads fontspec
  \defaultfontfeatures{Scale=MatchLowercase}
  \defaultfontfeatures[\rmfamily]{Ligatures=TeX,Scale=1}
\fi
\usepackage{lmodern}
\ifPDFTeX\else
  % xetex/luatex font selection
\fi
% Use upquote if available, for straight quotes in verbatim environments
\IfFileExists{upquote.sty}{\usepackage{upquote}}{}
\IfFileExists{microtype.sty}{% use microtype if available
  \usepackage[]{microtype}
  \UseMicrotypeSet[protrusion]{basicmath} % disable protrusion for tt fonts
}{}
\makeatletter
\@ifundefined{KOMAClassName}{% if non-KOMA class
  \IfFileExists{parskip.sty}{%
    \usepackage{parskip}
  }{% else
    \setlength{\parindent}{0pt}
    \setlength{\parskip}{6pt plus 2pt minus 1pt}}
}{% if KOMA class
  \KOMAoptions{parskip=half}}
\makeatother
% Make \paragraph and \subparagraph free-standing
\makeatletter
\ifx\paragraph\undefined\else
  \let\oldparagraph\paragraph
  \renewcommand{\paragraph}{
    \@ifstar
      \xxxParagraphStar
      \xxxParagraphNoStar
  }
  \newcommand{\xxxParagraphStar}[1]{\oldparagraph*{#1}\mbox{}}
  \newcommand{\xxxParagraphNoStar}[1]{\oldparagraph{#1}\mbox{}}
\fi
\ifx\subparagraph\undefined\else
  \let\oldsubparagraph\subparagraph
  \renewcommand{\subparagraph}{
    \@ifstar
      \xxxSubParagraphStar
      \xxxSubParagraphNoStar
  }
  \newcommand{\xxxSubParagraphStar}[1]{\oldsubparagraph*{#1}\mbox{}}
  \newcommand{\xxxSubParagraphNoStar}[1]{\oldsubparagraph{#1}\mbox{}}
\fi
\makeatother

\usepackage{color}
\usepackage{fancyvrb}
\newcommand{\VerbBar}{|}
\newcommand{\VERB}{\Verb[commandchars=\\\{\}]}
\DefineVerbatimEnvironment{Highlighting}{Verbatim}{commandchars=\\\{\}}
% Add ',fontsize=\small' for more characters per line
\usepackage{framed}
\definecolor{shadecolor}{RGB}{241,243,245}
\newenvironment{Shaded}{\begin{snugshade}}{\end{snugshade}}
\newcommand{\AlertTok}[1]{\textcolor[rgb]{0.68,0.00,0.00}{#1}}
\newcommand{\AnnotationTok}[1]{\textcolor[rgb]{0.37,0.37,0.37}{#1}}
\newcommand{\AttributeTok}[1]{\textcolor[rgb]{0.40,0.45,0.13}{#1}}
\newcommand{\BaseNTok}[1]{\textcolor[rgb]{0.68,0.00,0.00}{#1}}
\newcommand{\BuiltInTok}[1]{\textcolor[rgb]{0.00,0.23,0.31}{#1}}
\newcommand{\CharTok}[1]{\textcolor[rgb]{0.13,0.47,0.30}{#1}}
\newcommand{\CommentTok}[1]{\textcolor[rgb]{0.37,0.37,0.37}{#1}}
\newcommand{\CommentVarTok}[1]{\textcolor[rgb]{0.37,0.37,0.37}{\textit{#1}}}
\newcommand{\ConstantTok}[1]{\textcolor[rgb]{0.56,0.35,0.01}{#1}}
\newcommand{\ControlFlowTok}[1]{\textcolor[rgb]{0.00,0.23,0.31}{\textbf{#1}}}
\newcommand{\DataTypeTok}[1]{\textcolor[rgb]{0.68,0.00,0.00}{#1}}
\newcommand{\DecValTok}[1]{\textcolor[rgb]{0.68,0.00,0.00}{#1}}
\newcommand{\DocumentationTok}[1]{\textcolor[rgb]{0.37,0.37,0.37}{\textit{#1}}}
\newcommand{\ErrorTok}[1]{\textcolor[rgb]{0.68,0.00,0.00}{#1}}
\newcommand{\ExtensionTok}[1]{\textcolor[rgb]{0.00,0.23,0.31}{#1}}
\newcommand{\FloatTok}[1]{\textcolor[rgb]{0.68,0.00,0.00}{#1}}
\newcommand{\FunctionTok}[1]{\textcolor[rgb]{0.28,0.35,0.67}{#1}}
\newcommand{\ImportTok}[1]{\textcolor[rgb]{0.00,0.46,0.62}{#1}}
\newcommand{\InformationTok}[1]{\textcolor[rgb]{0.37,0.37,0.37}{#1}}
\newcommand{\KeywordTok}[1]{\textcolor[rgb]{0.00,0.23,0.31}{\textbf{#1}}}
\newcommand{\NormalTok}[1]{\textcolor[rgb]{0.00,0.23,0.31}{#1}}
\newcommand{\OperatorTok}[1]{\textcolor[rgb]{0.37,0.37,0.37}{#1}}
\newcommand{\OtherTok}[1]{\textcolor[rgb]{0.00,0.23,0.31}{#1}}
\newcommand{\PreprocessorTok}[1]{\textcolor[rgb]{0.68,0.00,0.00}{#1}}
\newcommand{\RegionMarkerTok}[1]{\textcolor[rgb]{0.00,0.23,0.31}{#1}}
\newcommand{\SpecialCharTok}[1]{\textcolor[rgb]{0.37,0.37,0.37}{#1}}
\newcommand{\SpecialStringTok}[1]{\textcolor[rgb]{0.13,0.47,0.30}{#1}}
\newcommand{\StringTok}[1]{\textcolor[rgb]{0.13,0.47,0.30}{#1}}
\newcommand{\VariableTok}[1]{\textcolor[rgb]{0.07,0.07,0.07}{#1}}
\newcommand{\VerbatimStringTok}[1]{\textcolor[rgb]{0.13,0.47,0.30}{#1}}
\newcommand{\WarningTok}[1]{\textcolor[rgb]{0.37,0.37,0.37}{\textit{#1}}}

\usepackage{longtable,booktabs,array}
\newcounter{none} % for unnumbered tables
\usepackage{calc} % for calculating minipage widths
% Correct order of tables after \paragraph or \subparagraph
\usepackage{etoolbox}
\makeatletter
\patchcmd\longtable{\par}{\if@noskipsec\mbox{}\fi\par}{}{}
\makeatother
% Allow footnotes in longtable head/foot
\IfFileExists{footnotehyper.sty}{\usepackage{footnotehyper}}{\usepackage{footnote}}
\makesavenoteenv{longtable}
\usepackage{graphicx}
\makeatletter
\newsavebox\pandoc@box
\newcommand*\pandocbounded[1]{% scales image to fit in text height/width
  \sbox\pandoc@box{#1}%
  \Gscale@div\@tempa{\textheight}{\dimexpr\ht\pandoc@box+\dp\pandoc@box\relax}%
  \Gscale@div\@tempb{\linewidth}{\wd\pandoc@box}%
  \ifdim\@tempb\p@<\@tempa\p@\let\@tempa\@tempb\fi% select the smaller of both
  \ifdim\@tempa\p@<\p@\scalebox{\@tempa}{\usebox\pandoc@box}%
  \else\usebox{\pandoc@box}%
  \fi%
}
% Set default figure placement to htbp
\def\fps@figure{htbp}
\makeatother


% definitions for citeproc citations
\NewDocumentCommand\citeproctext{}{}
\NewDocumentCommand\citeproc{mm}{%
  \begingroup\def\citeproctext{#2}\cite{#1}\endgroup}
\makeatletter
 % allow citations to break across lines
 \let\@cite@ofmt\@firstofone
 % avoid brackets around text for \cite:
 \def\@biblabel#1{}
 \def\@cite#1#2{{#1\if@tempswa , #2\fi}}
\makeatother
\newlength{\cslhangindent}
\setlength{\cslhangindent}{1.5em}
\newlength{\csllabelwidth}
\setlength{\csllabelwidth}{3em}
\newenvironment{CSLReferences}[2] % #1 hanging-indent, #2 entry-spacing
 {\begin{list}{}{%
  \setlength{\itemindent}{0pt}
  \setlength{\leftmargin}{0pt}
  \setlength{\parsep}{0pt}
  % turn on hanging indent if param 1 is 1
  \ifodd #1
   \setlength{\leftmargin}{\cslhangindent}
   \setlength{\itemindent}{-1\cslhangindent}
  \fi
  % set entry spacing
  \setlength{\itemsep}{#2\baselineskip}}}
 {\end{list}}
\usepackage{calc}
\newcommand{\CSLBlock}[1]{\hfill\break\parbox[t]{\linewidth}{\strut\ignorespaces#1\strut}}
\newcommand{\CSLLeftMargin}[1]{\parbox[t]{\csllabelwidth}{\strut#1\strut}}
\newcommand{\CSLRightInline}[1]{\parbox[t]{\linewidth - \csllabelwidth}{\strut#1\strut}}
\newcommand{\CSLIndent}[1]{\hspace{\cslhangindent}#1}

\ifLuaTeX
\usepackage[bidi=basic,shorthands=off]{babel}
\else
\usepackage[bidi=default,shorthands=off]{babel}
\fi
\ifLuaTeX
  \usepackage{selnolig} % disable illegal ligatures
\fi


\setlength{\emergencystretch}{3em} % prevent overfull lines

\providecommand{\tightlist}{%
  \setlength{\itemsep}{0pt}\setlength{\parskip}{0pt}}



 


\usepackage{booktabs}
\usepackage{float}
\usepackage{fontawesome5}
\usepackage{fvextra}
\usepackage{xurl}
\floatplacement{figure}{H}
\DefineVerbatimEnvironment{Highlighting}{Verbatim}{commandchars=\\\{\},breaklines=true,breakanywhere=true,fontsize=\footnotesize}
\KOMAoption{captions}{tableheading}
\makeatletter
\@ifpackageloaded{bookmark}{}{\usepackage{bookmark}}
\makeatother
\makeatletter
\@ifpackageloaded{caption}{}{\usepackage{caption}}
\AtBeginDocument{%
\ifdefined\contentsname
  \renewcommand*\contentsname{Inhaltsverzeichnis}
\else
  \newcommand\contentsname{Inhaltsverzeichnis}
\fi
\ifdefined\listfigurename
  \renewcommand*\listfigurename{Abbildungsverzeichnis}
\else
  \newcommand\listfigurename{Abbildungsverzeichnis}
\fi
\ifdefined\listtablename
  \renewcommand*\listtablename{Tabellenverzeichnis}
\else
  \newcommand\listtablename{Tabellenverzeichnis}
\fi
\ifdefined\figurename
  \renewcommand*\figurename{Abbildung}
\else
  \newcommand\figurename{Abbildung}
\fi
\ifdefined\tablename
  \renewcommand*\tablename{Tabelle}
\else
  \newcommand\tablename{Tabelle}
\fi
}
\@ifpackageloaded{float}{}{\usepackage{float}}
\floatstyle{ruled}
\@ifundefined{c@chapter}{\newfloat{codelisting}{h}{lop}}{\newfloat{codelisting}{h}{lop}[chapter]}
\floatname{codelisting}{Listing}
\newcommand*\listoflistings{\listof{codelisting}{Listingverzeichnis}}
\makeatother
\makeatletter
\makeatother
\makeatletter
\@ifpackageloaded{caption}{}{\usepackage{caption}}
\@ifpackageloaded{subcaption}{}{\usepackage{subcaption}}
\makeatother
\usepackage{bookmark}
\IfFileExists{xurl.sty}{\usepackage{xurl}}{} % add URL line breaks if available
\urlstyle{same}
\hypersetup{
  pdftitle={Langfristige Anlagestrategien anhand historischer Finanzdaten},
  pdfauthor={Riccardo Ambühl},
  pdflang={de},
  colorlinks=true,
  linkcolor={blue},
  filecolor={Maroon},
  citecolor={Blue},
  urlcolor={Blue},
  pdfcreator={LaTeX via pandoc}}


\title{Langfristige Anlagestrategien anhand historischer Finanzdaten}
\author{Riccardo Ambühl}
\date{2026-10-07}
\begin{document}
\maketitle

\renewcommand*\contentsname{Inhaltsverzeichnis}
{
\setcounter{tocdepth}{2}
\tableofcontents
}

\bookmarksetup{startatroot}

\chapter{Überblick}\label{uxfcberblick}

\textbf{Titel:} Langfristige Anlagestrategien anhand historischer
Finanzdaten\\
\textbf{Autor:} Riccardo Ambühl

Dieses Repository enthält ein Quarto-Book-Setup für die Maturarbeit. Die
Arbeit ist in einzelne Kapitel aufgeteilt, Python-Experimente liegen in
\texttt{notebooks/}, Quellen in \texttt{references.bib}, Styles in
\texttt{styles/} und spätere Ergebnisse in \texttt{outputs/}.

\section{Forschungsfrage}\label{forschungsfrage}

Wie können historische Finanzdaten mithilfe eines Python-basierten
Modells verwendet werden, um langfristige Investmentstrategien
hinsichtlich Risiko und Rendite zu vergleichen?

\section{Installationsempfehlungen}\label{installationsempfehlungen}

Installiere lokal:

\begin{itemize}
\tightlist
\item
  \href{https://quarto.org/}{Quarto}
\item
  \href{https://code.visualstudio.com/}{VS Code}
\item
  VS-Code-Erweiterungen: Quarto, Python, Jupyter
\item
  Python 3.11 oder neuer
\item
  Python-Pakete: \texttt{jupyter}, \texttt{pandas}, \texttt{matplotlib}
\item
  Zotero mit Better BibTeX
\end{itemize}

Beispiel für eine lokale Python-Umgebung:

\begin{Shaded}
\begin{Highlighting}[]
\NormalTok{python }\OperatorTok{{-}}\NormalTok{m venv }\OperatorTok{.}\FunctionTok{venv}
\OperatorTok{.}\NormalTok{\textbackslash{}}\OperatorTok{.}\FunctionTok{venv}\NormalTok{\textbackslash{}Scripts\textbackslash{}Activate}\OperatorTok{.}\FunctionTok{ps1}
\NormalTok{python }\OperatorTok{{-}}\NormalTok{m pip install }\OperatorTok{{-}{-}}\NormalTok{upgrade pip}
\NormalTok{python }\OperatorTok{{-}}\NormalTok{m pip install jupyter pandas matplotlib}
\end{Highlighting}
\end{Shaded}

\section{Quarto-Befehle}\label{quarto-befehle}

Vorschau im Browser:

\begin{Shaded}
\begin{Highlighting}[]
\NormalTok{quarto preview}
\end{Highlighting}
\end{Shaded}

Alle Formate rendern:

\begin{Shaded}
\begin{Highlighting}[]
\NormalTok{quarto render}
\end{Highlighting}
\end{Shaded}

Nur HTML rendern:

\begin{Shaded}
\begin{Highlighting}[]
\NormalTok{quarto render }\OperatorTok{{-}{-}}\NormalTok{to html}
\end{Highlighting}
\end{Shaded}

Nur PDF rendern:

\begin{Shaded}
\begin{Highlighting}[]
\NormalTok{quarto render }\OperatorTok{{-}{-}}\NormalTok{to pdf}
\end{Highlighting}
\end{Shaded}

\section{Zotero und Better BibTeX
Workflow}\label{zotero-und-better-bibtex-workflow}

\begin{enumerate}
\def\labelenumi{\arabic{enumi}.}
\tightlist
\item
  Quellen in Zotero sammeln.
\item
  Better BibTeX installieren.
\item
  Eine Sammlung für die Maturarbeit anlegen.
\item
  Die Sammlung als BibLaTeX oder BibTeX nach \texttt{references.bib}
  exportieren.
\item
  Wenn möglich automatische Aktualisierung aktivieren.
\item
  Kurze, stabile Citation Keys verwenden, zum Beispiel
  \texttt{markowitz1952}.
\end{enumerate}

\section{Quellen verwenden}\label{quellen-verwenden}

Quellen werden in den \texttt{.qmd}-Kapiteln mit Quarto-Zitationssyntax
eingebunden:

\begin{Shaded}
\begin{Highlighting}[]
\NormalTok{Die moderne Portfoliotheorie beschreibt Diversifikation systematisch }\CommentTok{[}\OtherTok{@markowitz1952}\CommentTok{]}\NormalTok{.}
\end{Highlighting}
\end{Shaded}

Die globale Bibliographie ist in \texttt{\_quarto.yml} über
\texttt{references.bib} eingebunden.

\section{Neue Kapitel hinzufügen}\label{neue-kapitel-hinzufuxfcgen}

\begin{enumerate}
\def\labelenumi{\arabic{enumi}.}
\tightlist
\item
  Neue Datei erstellen, zum Beispiel \texttt{resultate.qmd}.
\item
  Kapitelüberschrift mit \texttt{\#\ Resultate} beginnen.
\item
  Datei in \texttt{\_quarto.yml} unter \texttt{book:\ chapters:}
  ergänzen.
\item
  Mit \texttt{quarto\ preview} prüfen.
\end{enumerate}

\section{Notebooks sinnvoll
auslagern}\label{notebooks-sinnvoll-auslagern}

Nutze \texttt{notebooks/} für explorative Arbeit, Zwischenschritte und
Tests. Übertrage nur die relevanten, geprüften Ergebnisse in die
\texttt{.qmd}-Kapitel. Dadurch bleibt das Book lesbar, während der
Analyseprozess trotzdem nachvollziehbar bleibt.

\section{Projektstruktur}\label{projektstruktur}

\begin{Shaded}
\begin{Highlighting}[]
\NormalTok{.}
\NormalTok{|{-}{-} \_quarto.yml}
\NormalTok{|{-}{-} index.qmd}
\NormalTok{|{-}{-} einleitung.qmd}
\NormalTok{|{-}{-} theorie.qmd}
\NormalTok{|{-}{-} methodik.qmd}
\NormalTok{|{-}{-} fazit.qmd}
\NormalTok{|{-}{-} references.bib}
\NormalTok{|{-}{-} styles/}
\NormalTok{|{-}{-} src/}
\NormalTok{|{-}{-} notebooks/}
\NormalTok{|{-}{-} data/}
\NormalTok{|{-}{-} images/}
\NormalTok{|{-}{-} outputs/}
\NormalTok{\textasciigrave{}{-}{-} documentation/ai{-}usage/}
\end{Highlighting}
\end{Shaded}

\section{KI-Nachweise}\label{ki-nachweise}

Die Nutzung von KI-Tools wird unter \texttt{documentation/ai-usage/}
dokumentiert. Dort wird festgehalten, wann KI eingesetzt wurde, wofür
sie genutzt wurde, welche Inhalte übernommen wurden und was manuell
geprüft werden muss.

\section{Pendenzenliste}\label{pendenzenliste}

\begin{itemize}
\tightlist
\item[$\boxtimes$]
  \faIcon{calendar-alt} Spätestens KW 27 Disposition abgeben
\item[$\square$]
  \faIcon{calendar-alt} KW 27-29 Theorie
\item[$\square$]
  \faIcon{calendar-alt} KW 30 Daten importieren
\item[$\square$]
  \faIcon{calendar-alt} KW 31 Daten verarbeiten
\item[$\square$]
  \faIcon{calendar-alt} KW 32-38 Analyse

  \begin{itemize}
  \tightlist
  \item[$\square$]
    \faIcon{calendar-alt} KW 32-33 Funktionen für die Formeln erstellen
  \item[$\square$]
    \faIcon{calendar-alt} KW 34 Alle Funktionen bündeln und eine
    Simulations Funktion erstellen
  \item[$\square$]
    \faIcon{calendar-alt} KW 35-38 Simulation verschiedener Strategien
  \end{itemize}
\item[$\square$]
  \faIcon{calendar-alt} 25.09.2026 Probekapitel abgeben
\item[$\square$]
  \faIcon{calendar-alt} KW 39-40 Auswertung

  \begin{itemize}
  \tightlist
  \item[$\square$]
    \faIcon{calendar-alt} KW 39 Auswertungen bündeln
  \item[$\square$]
    \faIcon{calendar-alt} KW 40 Auswertungen Grafisch darstellen (SVG)
  \end{itemize}
\item[$\square$]
  \faIcon{calendar-alt} KW 41 Einleitung
\item[$\square$]
  \faIcon{calendar-alt} KW 42 Fazit
\item[$\square$]
  \faIcon{calendar-alt} 19.--30.10.2026 / KW 43--44 Definitiven Titel im
  Intranet eingeben
\item[$\square$]
  \faIcon{calendar-alt} KW 43-44, spätestens 30.10.2026 Erste Fassung
\item[$\square$]
  \faIcon{calendar-alt} KW 45 Gegenlesen
\item[$\square$]
  \faIcon{calendar-alt} KW 46 Überarbeitung
\item[$\square$]
  \faIcon{calendar-alt} KW 47 Schlusskontrolle und Abgabe vorbereiten

  \begin{itemize}
  \tightlist
  \item[$\square$]
    Abstract fertigstellen und elektronisch abgeben
  \item[$\square$]
    Titelblatt kontrollieren
  \item[$\square$]
    Inhaltsverzeichnis und Seitenzahlen kontrollieren
  \item[$\square$]
    Abbildungs- und Tabellenverzeichnis kontrollieren
  \item[$\square$]
    Literatur- und Quellenverzeichnis kontrollieren
  \item[$\square$]
    Ehrlichkeitserklärung ausfüllen, unterschreiben und als letzte Seite
    einfügen
  \item[$\square$]
    PDF exportieren und vollständig kontrollieren
  \item[$\square$]
    Zwei gebundene Exemplare drucken lassen
  \item[$\square$]
    Eigenes drittes Exemplar für die Präsentation vorbereiten
  \end{itemize}
\item[$\square$]
  \faIcon{calendar-alt} Mittwoch, 25.11.2026, mittags Maturitätsarbeit
  abgeben

  \begin{itemize}
  \tightlist
  \item[$\square$]
    Zwei gebundene Exemplare abgeben
  \item[$\square$]
    Elektronische PDF-Version abgeben
  \item[$\square$]
    Webseite hochladen
  \end{itemize}
\item[$\square$]
  \faIcon{calendar-alt} KW 2--3 2027 Präsentation erstellen
\item[$\square$]
  \faIcon{calendar-alt} KW 4 2027: Präsentation üben und auf Fragen
  vorbereiten
\item[$\square$]
  \faIcon{calendar-alt} Samstag, 30.01.2027 Präsentation
\end{itemize}

\textbf{Optional:}

\begin{itemize}
\tightlist
\item[$\square$]
  \faIcon{calendar-alt} Simulation Tool basierend auf der Simulation
  Funktion
\item[$\square$]
  \faIcon{calendar-alt} Automatisiertes testen von weiteren Strategien
  auf dem Server
\item[$\square$]
  \faIcon{calendar-alt} Verschiedene anschauliche Grafik Typen (SVG)
\end{itemize}

\section{Quellen}\label{quellen}

\subsection{Yahoo Finance}\label{yahoo-finance}

Yahoo Finance bietet historische Kursdaten für eine Vielzahl von
Finanzinstrumenten, darunter Aktien, ETFs und Anleihen. Diese Daten
können für die Berechnung von Renditen, Volatilität und anderen
Kennzahlen verwendet werden.

\subsection{FRED (Federal Reserve Economic
Data)}\label{fred-federal-reserve-economic-data}

FRED ist eine umfangreiche Datenbank der Federal Reserve Bank of
St.~Louis, die wirtschaftliche und finanzielle Daten bereitstellt. Hier
können makroökonomische Indikatoren wie das Bruttoinlandsprodukt (BIP),
Arbeitslosenzahlen und Zinssätze abgerufen werden, die für die Analyse
von Anlagestrategien relevant sein können.

\subsection{Vanguard}\label{vanguard}

Vanguard ist einer der grössten Anbieter von ETFs und bietet
umfangreiche Informationen zu seinen Produkten, einschliesslich
historischer Renditen, Kostenquoten und Fondsinformationen. Diese Daten
können für die Analyse von ETF-basierten Anlagestrategien verwendet
werden.

\subsection{Fidelity Investments}\label{fidelity-investments}

Fidelity ist ein internationaler Online-Anbieter im Finanzsektor, der
neben Finanzdienstleistungen auch fundiertes Fachwissen über
verschiedene Anlagestrategien vermittelt. Für die vorliegende Arbeit
dient die Plattform als wertvolle Quelle, um tiefere Einblicke in die
praktische Nutzung und Anwendung dieser Finanztheorien zu gewinnen.

\subsection{Souverän investieren mit Indexfonds und ETFs (Autor: Gerd
Kommer, ISBN:
978-3-593-52064-3)}\label{souveruxe4n-investieren-mit-indexfonds-und-etfs-autor-gerd-kommer-isbn-978-3-593-52064-3}

Ich erhoffe mir von diesem Buch das theoretische Fundament für meine
passiven Anlagestrategien. Es soll mir dabei helfen, das breit gestreute
Investieren mit ETFs, das 60/40-Portfolio und die BIP-Gewichtung zu
erklären.

\subsection{Der leichte Einstieg in die Welt der ETFs (Autor: Gerd
Kommer, ISBN:
978-3-95972-543-9)}\label{der-leichte-einstieg-in-die-welt-der-etfs-autor-gerd-kommer-isbn-978-3-95972-543-9}

Ich erhoffe mir von diesem Buch ganz einfache Definitionen für die
Grundlagen meiner Arbeit. Ich will damit Begriffe wie ``ETF'' oder die
Reinvestition von Dividenden erklären.

\subsection{Python for Data Analysis (Autor: Wes McKinney, ISBN:
978-1-09-810403-0)}\label{python-for-data-analysis-autor-wes-mckinney-isbn-978-1-09-810403-0}

Ich erhoffe mir von diesem Buch das Wissen, wie ich die historischen
Kurse von Yahoo Finance und FRED richtig verarbeite. Es soll mir zeigen,
wie ich die Daten filtere und anpasse.

\subsection{Python for Finance (Autor: Yves Hilpisch, ISBN:
978-1-4920-2433-0)}\label{python-for-finance-autor-yves-hilpisch-isbn-978-1-4920-2433-0}

Ich erhoffe mir von diesem Buch die mathematische Erklärung für meine
Trendfolgestrategie. Ich will daraus lernen, wie Indikatoren wie der MA
und der RSI genau funktionieren und wie man Kennzahlen wie die Sharpe
Ratio und den maximalen Drawdown berechnet.

\textbf{Python Libraries und wie ich sie anwende:}

\begin{itemize}
\tightlist
\item
  yfinance (Yahoo Finance Kursdaten)
\item
  fredapi (Zugriff auf FRED)
\item
  pandas (Datenverarbeitung, zentrales Werkzeug für Zeitreihen,
  Renditen, Resampling)
\item
  numpy (numerische Basis, Volatilität, kumulierte Renditen)
\item
  quantstats (Performance Reports: Sharpe, Drawdown, CAGR)
\item
  empyrical (Risiko-/Renditkennzahlen)
\item
  vectorbt oder backtesting.py (Backtesting)
\item
  pandas-ta (Technische Indikatoren)
\item
  matplotlib (Visualisierung)
\item
  plotly (interaktive Kursverläufe)
\item
  bt (Portfolio-Backtesting und Rebalancing)
\end{itemize}

Das ist eine Auswahl an Libraries, die ich testen möchte, jedoch nicht
unbedingt alle schlussendlich verwenden werde. Es ist sinnvoll,
möglichst viele geeignete Libraries zu nutzen, um den Aufwand gering zu
halten. Das Ziel ist es, etwas Neues zu schaffen und nicht bereits
Vorhandenes erneut zu programmieren.

\bookmarksetup{startatroot}

\chapter{Einleitung}\label{einleitung}

\bookmarksetup{startatroot}

\chapter{Theorie}\label{theorie}

3.1 Grundlagen historischer Finanzdaten

3.1.1 Indexfonds und ETFs

Ein Index bildet die Wertentwicklung eines festgelegten Marktbereichs
als rechnerische Kennzahl ab. Der S\&P 500 umfasst beispielsweise die
500 grössten Unternehmen der USA. Weil sich ein marktbreiter Index auf
zahlreiche Titel stützt, hängt seine Entwicklung nicht von einem
einzelnen Unternehmen ab. Ein Indexfonds versucht, seinen Referenzindex
möglichst genau nachzubilden, anstatt ihn zu
übertreffen.\footnote{Kommer, \emph{Der leichte Einstieg in die Welt der
  ETFs}, 2. vollständig überarbeitete und aktualisierte Neuausgabe
  (München: FBV, FinanzBuch Verlag, 2025), 464--65.}Welche Wertpapiere
im Index enthalten sind und wie stark sie gewichtet werden, bestimmt
dessen Regelwerk. Es regelt ebenfalls, unter welchen Voraussetzungen
Titel hinzukommen oder entfernt werden. Die Zusammensetzung wird
üblicherweise in einem festen Rhythmus, etwa halb- oder vierteljährlich,
überprüft.\footnote{Kommer, 465.}

Bei einem ETF werden Fondsanteile während der Börsenzeiten an einem
Handelsplatz gehandelt. Der Zugang kann über eine Bank oder einen Broker
erfolgen; ausgeführt wird der Kauf an der Börse. Kommer zählt Indexfonds
und ETFs zu den bedeutendsten Finanzinnovationen seit 1971.\footnote{Kommer,
  58.} Für Privatanleger sind sie unter anderem interessant, weil ihr
Aufbau vergleichsweise leicht nachvollziehbar ist, die Zusammensetzung
offengelegt wird und die Verwaltung mit wenig eigenem Aufwand verbunden
ist. Hinzu kommen in der Regel geringe Gebühren. Ihre Konstruktion
beruht zudem auf Erkenntnissen der Kapitalmarktforschung.\footnote{Kommer,
  59--60.}

3.1.2 Eignung für die Untersuchung

Marktbreite ETFs sind für meine Untersuchung gut geeignet, da sie ganze
Anlageklassen mit geringem Aufwand abbilden. ETFs sind günstig, bequem
und transparent. Für meine Untersuchung ist der Total Return besser
geeignet als der Kursindex. Der Kursindex berücksichtigt keine
Dividenden, die wichtig sind, um die wahre Performance nachzuweisen.

3.2 Historische Finanzdaten und Zeitreihen

3.2.1 Zeitreihen

Historische Finanzdaten bilden die Grundlage für den Vergleich der
Anlagestrategien. Eine finanzielle Zeitreihe entsteht, wenn der Kurs
eines Indexfonds über einen längeren Zeitraum erfasst wird. Die
Zeitreihe besteht also aus Beobachtungen, die einem bestimmten Datum
zugeordnet werden. Wichtig ist dabei die Reihenfolge, denn sie zeigt die
Veränderung von einer Beobachtung zur nächsten. Für diese Zeitreihen
wird pandas in Python sehr wichtig sein. Wes McKinney, der Hauptautor
von pandas, arbeitete als Analyst bei AQR Capital Management, als er mit
der Entwicklung der Library begann. Man kann also sicher sagen, dass
pandas darauf ausgelegt ist, mit Finanzdaten zu arbeiten.\footnote{Hilpisch,
  \emph{Python for Finance}, 304.}

3.2.2 Frequenzen

Zeitreihen können ausserdem in verschiedenen Frequenzen vorliegen:
täglich, wöchentlich, monatlich und jährlich. Diese Frequenzen werden
für meine Arbeit relevant sein. Da meine Untersuchung langfristig
angesetzt ist, sind kleinere Beobachtungsintervalle beziehungsweise
kleinere ``Timeframes'', auf Deutsch Zeitrahmen, nicht von Bedeutung.

{\def\LTcaptype{none} % do not increment counter
\begin{longtable}[]{@{}ll@{}}
\toprule\noalign{}
Alias & Description \\
\midrule\noalign{}
\endhead
\bottomrule\noalign{}
\endlastfoot
\texttt{B} & Business day frequency \\
\texttt{C} & Custom business day frequency (experimental) \\
\texttt{D} & Calendar day frequency \\
\texttt{W} & Weekly frequency \\
\texttt{M} & Month end frequency \\
\texttt{BM} & Business month end frequency \\
\texttt{MS} & Month start frequency \\
\texttt{BMS} & Business month start frequency \\
\texttt{Q} & Quarter end frequency \\
\texttt{BQ} & Business quarter end frequency \\
\texttt{QS} & Quarter start frequency \\
\texttt{BQS} & Business quarter start frequency \\
\texttt{A} & Year end frequency \\
\texttt{BA} & Business year end frequency \\
\texttt{AS} & Year start frequency \\
\texttt{BAS} & Business year start frequency \\
\texttt{H} & Hourly frequency \\
\texttt{T} & Minutely frequency \\
\texttt{S} & Secondly frequency \\
\texttt{L} & Milliseconds \\
\texttt{U} & Microseconds \\
\end{longtable}
}

Ausführliche Liste mit Beispielfrequenzen.\footnote{Hilpisch, 191--92.}

3.2.3 Resampling

Wenn ein Datensatz nicht im gewünschten Zeitrahmen vorliegt, kann man
ihn mithilfe des ``Resamplings'' umwandeln. Normalerweise geschieht das
von einem kleineren Zeitrahmen zu einem grösseren. Das nennt man
``Downsampling''. Es wird später vor allem wichtig sein, um die Daten in
ein einheitliches Format zu bringen und sie vergleichen zu können.

Für das ``Resampling'' gibt es in pandas eine eingebaute Funktion, die
oben stark vereinfacht dargestellt ist. Nun folgt ein konkretes Beispiel
für die Nutzung dieser Funktion.

\begin{Shaded}
\begin{Highlighting}[]
\CommentTok{\# Importiert die Pandas{-}Bibliothek}
\ImportTok{import}\NormalTok{ pandas }\ImportTok{as}\NormalTok{ pd}

\CommentTok{\# Die Funktion resample\_dataframe() nimmt ein DataFrame, eine Frequenz und eine Aggregationsmethode als Eingabe und gibt ein resampletes DataFrame zurück}
\KeywordTok{def}\NormalTok{ resample\_dataframe(df: pd.DataFrame, frequenz: }\BuiltInTok{str}\NormalTok{, aggregation: }\BuiltInTok{dict}\NormalTok{):}
    \ControlFlowTok{return}\NormalTok{ df.resample(frequenz).agg(aggregation)}

\CommentTok{\# Beispiel zur Nutzung der Funktion}
\NormalTok{df }\OperatorTok{=}\NormalTok{ pd.DataFrame(\{}
    \StringTok{"Datum"}\NormalTok{: [}\StringTok{"2026{-}07{-}27"}\NormalTok{, }\StringTok{"2026{-}07{-}28"}\NormalTok{, }\StringTok{"2026{-}07{-}29"}\NormalTok{, }\StringTok{"2026{-}07{-}30"}\NormalTok{, }\StringTok{"2026{-}07{-}31"}\NormalTok{, }\StringTok{"2026{-}08{-}01"}\NormalTok{, }\StringTok{"2026{-}08{-}02"}\NormalTok{, }\StringTok{"2026{-}08{-}03"}\NormalTok{, }\StringTok{"2026{-}08{-}04"}\NormalTok{, }\StringTok{"2026{-}08{-}05"}\NormalTok{, }\StringTok{"2026{-}08{-}06"}\NormalTok{, }\StringTok{"2026{-}08{-}07"}\NormalTok{, }\StringTok{"2026{-}08{-}08"}\NormalTok{, }\StringTok{"2026{-}08{-}09"}\NormalTok{],}
    \StringTok{"Kurs"}\NormalTok{: [}\DecValTok{101}\NormalTok{, }\DecValTok{104}\NormalTok{, }\DecValTok{105}\NormalTok{, }\DecValTok{106}\NormalTok{, }\DecValTok{108}\NormalTok{, }\DecValTok{107}\NormalTok{, }\DecValTok{110}\NormalTok{, }\DecValTok{112}\NormalTok{, }\DecValTok{111}\NormalTok{, }\DecValTok{115}\NormalTok{, }\DecValTok{116}\NormalTok{, }\DecValTok{118}\NormalTok{, }\DecValTok{117}\NormalTok{, }\DecValTok{120}\NormalTok{]}
\NormalTok{\})}

\CommentTok{\# Konvertiert die Spalte "Datum" in das Datetime{-}Format, um Zeitreihenoperationen zu ermöglichen.}
\NormalTok{df[}\StringTok{"Datum"}\NormalTok{] }\OperatorTok{=}\NormalTok{ pd.to\_datetime(df[}\StringTok{"Datum"}\NormalTok{])}
\CommentTok{\# Setzt die Spalte "Datum" als Index des DataFrames}
\NormalTok{df }\OperatorTok{=}\NormalTok{ df.set\_index(}\StringTok{"Datum"}\NormalTok{)}

\NormalTok{weekly }\OperatorTok{=}\NormalTok{ resample\_dataframe(}
\NormalTok{    df,}
    \StringTok{"W"}\NormalTok{, }\CommentTok{\# Resampling{-}Frequenz: "W" steht für wöchentlich}
\NormalTok{    \{}\StringTok{"Kurs"}\NormalTok{: }\StringTok{"last"}\NormalTok{\} }\CommentTok{\# Aggregationsmethode: "last" bedeutet, dass der letzte Kurswert jeder Woche genommen wird}
\NormalTok{)}

\BuiltInTok{print}\NormalTok{(}\SpecialStringTok{f"Wöchentliche Kurse:}\CharTok{\textbackslash{}n}\SpecialCharTok{\{}\NormalTok{weekly}\SpecialCharTok{\}}\SpecialStringTok{"}\NormalTok{)}
\end{Highlighting}
\end{Shaded}

\begin{verbatim}
Wöchentliche Kurse:
            Kurs
Datum           
2026-08-02   110
2026-08-09   120
\end{verbatim}

Quelle zur Nutzung von to\_datetime().\footnote{McKinney, \emph{Python
  for Data Analysis}, 511.}

Um Vorausschauverzerrung zu vermeiden, sollte man grundsätzlich immer
den letzten verfügbaren Wert des Intervalls verwenden.

3.3 Rendite

3.3.1 Einfache Rendite

Hilpisch unterscheidet absolute Veränderungen und prozentuale
Veränderungen beziehungsweise einfache Renditen.\footnote{Hilpisch,
  \emph{Python for Finance}, 315--16.}

Absolute Änderung: \[
\Delta P_t=P_t-P_{t-1}
\] Einfache Rendite: \[
r_t=\frac{P_t}{P_{t-1}}-1
\]

Quelle der Formeln und der entsprechenden Pandas-Umsetzung:.\footnote{Hilpisch,
  315--16; {„Pandas.{Series}.Pct\_change --- Pandas 3.0.5
  Documentation``}.}

\begin{Shaded}
\begin{Highlighting}[]
\CommentTok{\# Importiert die Pandas{-}Bibliothek}
\ImportTok{import}\NormalTok{ pandas }\ImportTok{as}\NormalTok{ pd}

\CommentTok{\# Die Funktion absolute\_aenderung() berechnet die absolute Änderung der Werte in einer Pandas{-}Series}
\KeywordTok{def}\NormalTok{ absolute\_aenderung(df: pd.Series):}
    \CommentTok{\# df.diff() berechnet den absoluten Unterschied zum vorherigen Wert}
    \ControlFlowTok{return}\NormalTok{ df.diff()}

\CommentTok{\# Die Funktion prozentuale\_aenderung() berechnet die prozentuale Änderung der Werte in einer Pandas{-}Series}
\KeywordTok{def}\NormalTok{ prozentuale\_aenderung(df: pd.Series):}
    \CommentTok{\# df.pct\_change() berechnet die prozentuale Änderung zum vorherigen Wert}
    \ControlFlowTok{return}\NormalTok{ df.pct\_change()}

\CommentTok{\# Beispiel zur Nutzung der Funktionen}
\NormalTok{df }\OperatorTok{=}\NormalTok{ pd.Series([}\DecValTok{100}\NormalTok{, }\DecValTok{110}\NormalTok{, }\FloatTok{126.5}\NormalTok{])}

\NormalTok{aenderung\_absolut }\OperatorTok{=}\NormalTok{ absolute\_aenderung(df)}

\NormalTok{aenderung\_prozentuale }\OperatorTok{=}\NormalTok{ prozentuale\_aenderung(df)}

\CommentTok{\# Der erste Wert ist NaN, weil es keinen vorherigen Wert gibt}
\BuiltInTok{print}\NormalTok{(}\SpecialStringTok{f"Absolute Änderung:}\CharTok{\textbackslash{}n}\SpecialCharTok{\{}\NormalTok{aenderung\_absolut}\SpecialCharTok{\}}\SpecialStringTok{"}\NormalTok{)}

\CommentTok{\# Die prozentuale Änderung wird als Dezimalzahl dargestellt. Für Prozent mit 100 multiplizieren}
\BuiltInTok{print}\NormalTok{(}\SpecialStringTok{f"Prozentuale Änderung:}\CharTok{\textbackslash{}n}\SpecialCharTok{\{}\NormalTok{aenderung\_prozentuale }\OperatorTok{*} \DecValTok{100}\SpecialCharTok{\}}\SpecialStringTok{"}\NormalTok{)}
\end{Highlighting}
\end{Shaded}

\begin{verbatim}
Absolute Änderung:
0     NaN
1    10.0
2    16.5
dtype: float64
Prozentuale Änderung:
0     NaN
1    10.0
2    15.0
dtype: float64
\end{verbatim}

3.3.2 Logarithmische Rendite und kumulierte Rendite

Logarithmische Renditen lassen sich über die Zeit addieren. Durch die
Exponentialfunktion erhält man aus ihrer Summe wieder den
Wachstumsfaktor.\footnote{Hilpisch, \emph{Python for Finance}, 317--18.}

Logarithmische Rendite: \[
\ell_t=\ln\left(\frac{P_t}{P_{t-1}}\right)
\] Kumulierte logarithmische Rendite und Wachstumsfaktor: \[
L_t=\sum_{i=1}^{t}\ell_i,\qquad G_t=e^{L_t}
\]

Quelle der Formeln:.\footnote{Hilpisch, 317--18.}

\begin{Shaded}
\begin{Highlighting}[]
\CommentTok{\# Importiert die Numpy{-}Bibliothek}
\ImportTok{import}\NormalTok{ numpy }\ImportTok{as}\NormalTok{ np}
\CommentTok{\# Importiert die Pandas{-}Bibliothek}
\ImportTok{import}\NormalTok{ pandas }\ImportTok{as}\NormalTok{ pd}

\CommentTok{\# Die Funktion kumulierte\_rendite() berechnet die kumulierte logarithmische Rendite}
\KeywordTok{def}\NormalTok{ kumulierte\_rendite(preise: pd.Series):}
    \CommentTok{\# np.log() berechnet den natürlichen Logarithmus}
    \CommentTok{\# und shift(1) verschiebt die Serie um einen Tag zurück, um die Rendite zu berechnen}
\NormalTok{    log\_rendite }\OperatorTok{=}\NormalTok{ np.log(preise }\OperatorTok{/}\NormalTok{ preise.shift(}\DecValTok{1}\NormalTok{))}
    \CommentTok{\# cumsum() summiert die logarithmischen Renditen über die Zeit}
\NormalTok{    kumulierte\_log\_rendite }\OperatorTok{=}\NormalTok{ log\_rendite.cumsum()}
    \ControlFlowTok{return}\NormalTok{ kumulierte\_log\_rendite}

\CommentTok{\# Die Funktion wachstumsfaktor() berechnet den Gesamtwachstumsfaktor aus der kumulierten logarithmischen Rendite}
\KeywordTok{def}\NormalTok{ wachstumsfaktor(preise: pd.Series):}
\NormalTok{    kumulierte\_log\_rendite }\OperatorTok{=}\NormalTok{ kumulierte\_rendite(preise)}
    \CommentTok{\# np.exp() macht den Logarithmus rückgängig und zeigt den Gesamtwachstumsfaktor}
\NormalTok{    wachstumsfaktoren }\OperatorTok{=}\NormalTok{ np.exp(kumulierte\_log\_rendite)}
    \ControlFlowTok{return}\NormalTok{ wachstumsfaktoren}


\CommentTok{\# Beispiel zur Nutzung der Funktionen}
\NormalTok{df }\OperatorTok{=}\NormalTok{ pd.Series([}\DecValTok{100}\NormalTok{, }\DecValTok{110}\NormalTok{, }\FloatTok{126.5}\NormalTok{])}

\BuiltInTok{print}\NormalTok{(}\SpecialStringTok{f"Kumulierte logarithmische Rendite:}\CharTok{\textbackslash{}n}\SpecialCharTok{\{}\NormalTok{kumulierte\_rendite(df)}\SpecialCharTok{\}}\CharTok{\textbackslash{}n}\SpecialStringTok{"}\NormalTok{)}
\BuiltInTok{print}\NormalTok{(}\SpecialStringTok{f"Wachstumsfaktor:}\CharTok{\textbackslash{}n}\SpecialCharTok{\{}\NormalTok{wachstumsfaktor(df)}\SpecialCharTok{\}}\SpecialStringTok{"}\NormalTok{)}
\end{Highlighting}
\end{Shaded}

\begin{verbatim}
Kumulierte logarithmische Rendite:
0         NaN
1    0.095310
2    0.235072
dtype: float64

Wachstumsfaktor:
0      NaN
1    1.100
2    1.265
dtype: float64
\end{verbatim}

3.3.3 Geometrische Rendite

Die geometrische Rendite beschreibt das tatsächlich realisierte
durchschnittliche Wachstum pro Periode über mehrere Zeiträume hinweg.
Der arithmetische Durchschnitt ist der normale mathematische Mittelwert.
Dieser wird zum Beispiel auch bei der Sharpe Ratio genutzt.\footnote{Pham,
  Cui, und Ruthbah, \emph{The {Performance} of the 60/40 {Portfolio}},
  50--51.} Eine Form der geometrischen Rendite ist die annualisierte
Rendite. Dabei wird die Rendite auf ein Jahr normiert.

Geometrischer Mittelwert (Formel): \[
\operatorname{Mean}(R)
=
\sqrt[N]{\prod_{i=1}^{N} R_i}
\] Dabei bezeichnet \textbf{\(\operatorname{Mean}(R)\)} den
geometrischen Mittelwert, \textbf{\(N\)} die Anzahl der Perioden,
\textbf{\(R_i\)} den Renditefaktor einer einzelnen Periode,
\textbf{\(\prod\)} das Produkt aller Renditefaktoren und
\textbf{\(\sqrt[N]{\cdot}\)} die \emph{N}-te Wurzel.\footnote{Pham, Cui,
  und Ruthbah, 51.} Annualisierte Rendite bei m Perioden pro Jahr: \[
r_{\mathrm{ann}}=\operatorname{Mean}(R)^m-1
\]

Die Annualisierungsformel ist die direkte Umrechnung des zuvor
definierten geometrischen Periodenmittels auf m Perioden pro Jahr.

\begin{Shaded}
\begin{Highlighting}[]
\CommentTok{\# Importiert die Numpy{-}Bibliothek}
\ImportTok{import}\NormalTok{ numpy }\ImportTok{as}\NormalTok{ np}

\CommentTok{\# Die Funktion geometrisches\_mittel() berechnet das geometrische Mittel einer Liste von Werten}
\KeywordTok{def}\NormalTok{ geometrisches\_mittel(werte: np.ndarray):}
    \CommentTok{\# np.log() berechnet den natürlichen Logarithmus jedes Wertes,}
    \CommentTok{\# np.mean() berechnet den Durchschnitt der Logarithmen}
    \CommentTok{\# und np.exp() macht den Logarithmus rückgängig}
    \ControlFlowTok{return}\NormalTok{ np.exp(np.mean(np.log(werte)))}

\CommentTok{\# Beispiel zur Nutzung der Funktion}
\NormalTok{werte }\OperatorTok{=}\NormalTok{ np.array([}\FloatTok{1.1}\NormalTok{, }\FloatTok{1.2}\NormalTok{, }\FloatTok{0.9}\NormalTok{, }\FloatTok{1.05}\NormalTok{])}
\BuiltInTok{print}\NormalTok{(}\SpecialStringTok{f"Geometrisches Mittel: }\SpecialCharTok{\{}\NormalTok{geometrisches\_mittel(werte)}\SpecialCharTok{\}}\SpecialStringTok{"}\NormalTok{)}
\end{Highlighting}
\end{Shaded}

\begin{verbatim}
Geometrisches Mittel: 1.0568210009926833
\end{verbatim}

\begin{Shaded}
\begin{Highlighting}[]
\CommentTok{\# Die Funktion annualisierte\_rendite() berechnet die annualisierte Rendite aus einer Liste von Werten und der Anzahl der Perioden pro Jahr}
\KeywordTok{def}\NormalTok{ annualisierte\_rendite(werte: np.ndarray, perioden\_pro\_jahr: }\BuiltInTok{int}\NormalTok{):}
    \ControlFlowTok{return}\NormalTok{ geometrisches\_mittel(werte) }\OperatorTok{**}\NormalTok{ perioden\_pro\_jahr }\OperatorTok{{-}} \DecValTok{1}

\CommentTok{\# Beispiel zur Nutzung der Funktion}
\NormalTok{werte }\OperatorTok{=}\NormalTok{ np.array([}\FloatTok{1.1}\NormalTok{, }\FloatTok{1.2}\NormalTok{, }\FloatTok{0.9}\NormalTok{, }\FloatTok{1.05}\NormalTok{])}
\NormalTok{perioden\_pro\_jahr }\OperatorTok{=} \DecValTok{12}  \CommentTok{\# zum Beispiel monatliche Renditen}
\BuiltInTok{print}\NormalTok{(}\SpecialStringTok{f"Annualisierte Rendite: }\SpecialCharTok{\{}\NormalTok{annualisierte\_rendite(werte, perioden\_pro\_jahr)}\SpecialCharTok{\}}\SpecialStringTok{"}\NormalTok{)}
\end{Highlighting}
\end{Shaded}

\begin{verbatim}
Annualisierte Rendite: 0.9409628324239998
\end{verbatim}

3.3.4 Nominal und real

Die reale Rendite ist die inflationsbereinigte Rendite. Vor allem über
längere Zeiträume ist es sehr wichtig, das zu berücksichtigen. Aus der
Grafik des Bundesamts für Statistik kann man für die letzten 10 Jahre
eine durchschnittliche Inflation von 0.8\% in der Schweiz
ablesen.\footnote{{„Preise``}.} Das heisst, ein Investment muss jährlich
mindestens 0.8\% erzielen, damit die Kaufkraft erhalten bleibt.

3.4 Volatilität

3.4.1 Bedeutung der Volatilität

Zwei Anlagestrategien können die gleiche Rendite erzielen, aber
unterschiedlich stark schwanken. Die Volatilität beschreibt die Stärke
dieser Schwankungen. Je höher die Volatilität ist, desto höher wird das
Schwankungsrisiko eingeschätzt. Kommer bezeichnet die Volatilität als
das verbreitetste Risikomass der Finanzökonomie. Ausserdem erklärt er,
dass die Begriffe Volatilität und Standardabweichung im Anlagekontext
häufig gleichbedeutend verwendet werden.\footnote{Kommer, \emph{Souverän
  investieren mit Indexfonds und ETFs}, 7. Auflage (Weinheim: Campus
  Verlag, 2025), 143--44.} Da die Volatilität in die Berechnung der
Sharpe Ratio einfliesst, beeinflusst sie auch die Beurteilung einer
Rendite im Verhältnis zum übernommenen Risiko.

3.4.2 Berechnung der Volatilität

Die historische Volatilität wird anhand der Standardabweichung der
periodischen Renditen bestimmt. Dazu wird zuerst der arithmetische
Mittelwert der betrachteten Renditen berechnet. Anschliessend wird
untersucht, wie stark die Rendite jeder einzelnen Periode von diesem
Durchschnitt abweicht. Die Standardabweichung fasst diese Abweichungen
in einer Kennzahl zusammen.\footnote{Kommer, 144.}

Formel für die Standardabweichung: \[
\sigma = \sqrt{\frac{1}{n-1}\sum_{i=1}^{n}(r_i-\bar{r})^2}
\]

Die Formel entspricht der im Code verwendeten
Stichproben-Standardabweichung mit ddof=1.\footnote{{„Pandas.{Series}.Std
  --- Pandas 3.0.5 Documentation``}.} Kommer erläutert Bedeutung und
Interpretation der Standardabweichung im Anlagekontext.\footnote{Kommer,
  \emph{Souverän investieren mit Indexfonds und ETFs}, 7. Auflage
  (Weinheim: Campus Verlag, 2025), 144.}

Dabei bezeichnet \textbf{\(\sigma\)} die Standardabweichung,
\textbf{\(n\)} die Anzahl der Renditen, \textbf{\(r_i\)} die Rendite
einer einzelnen Periode und \textbf{\(\bar r\)} den arithmetischen
Mittelwert aller Renditen. Die Abweichungen vom Mittelwert werden
quadriert. Deshalb fliessen sowohl negative als auch positive
Abweichungen in die Berechnung ein.\footnote{Kommer, 144.}

Wichtig ist, dass man prozentuale Renditen und nicht absolute
Kursänderungen verwendet. Sonst kann man verschiedene Kurse nicht
miteinander vergleichen.

\begin{Shaded}
\begin{Highlighting}[]
\CommentTok{\# Importiert die Numpy{-}Bibliothek}
\ImportTok{import}\NormalTok{ numpy }\ImportTok{as}\NormalTok{ np}
\CommentTok{\# Importiert die Pandas{-}Bibliothek}
\ImportTok{import}\NormalTok{ pandas }\ImportTok{as}\NormalTok{ pd}

\CommentTok{\# Die Funktion standardabweichung() berechnet die Standardabweichung einer Liste von Werten}
\KeywordTok{def}\NormalTok{ standardabweichung(werte: np.ndarray):}
    \CommentTok{\# Werte in eine Pandas{-}Series umwandeln}
\NormalTok{    werte }\OperatorTok{=}\NormalTok{ pd.Series(werte, dtype}\OperatorTok{=}\StringTok{"float64"}\NormalTok{)}

    \CommentTok{\# Standardabweichung mithilfe der Pandas{-}Funktion .std() berechnen}
    \ControlFlowTok{return}\NormalTok{ werte.std(ddof}\OperatorTok{=}\DecValTok{1}\NormalTok{)}


\CommentTok{\# Beispiel zur Nutzung der Funktion}
\NormalTok{werte }\OperatorTok{=}\NormalTok{ np.array([}\FloatTok{1.1}\NormalTok{, }\FloatTok{1.2}\NormalTok{, }\FloatTok{0.9}\NormalTok{, }\FloatTok{1.05}\NormalTok{])}
\BuiltInTok{print}\NormalTok{(}\SpecialStringTok{f"Standardabweichung: }\SpecialCharTok{\{}\NormalTok{standardabweichung(werte)}\SpecialCharTok{\}}\SpecialStringTok{"}\NormalTok{)}
\end{Highlighting}
\end{Shaded}

\begin{verbatim}
Standardabweichung: 0.12499999999999997
\end{verbatim}

3.4.3 Annualisierung

Die berechnete Standardabweichung bezieht sich zunächst auf die
verwendete Frequenz. Wählt man beispielsweise für einen Zeitraum von
einem Monat eine Periodenlänge von einem Tag, erhält man eine tägliche
Volatilität. Damit man verschiedene Datensätze vergleichen kann, rechnet
man sie normalerweise auf ein Jahr hoch. Das nennt man Annualisierung.

Formel zur Berechnung der annualisierten Volatilität:

\[
\sigma_{\mathrm{ann}} = \sigma_{\mathrm{Periode}}\sqrt{m}
\]

Quelle der Formel:.\footnote{Hilpisch, \emph{Python for Finance}, 694.}

Dabei bezeichnet \textbf{\(m\)} die Anzahl der betrachteten Perioden pro
Jahr. Die tägliche Standardabweichung
\textbf{\(\sigma_{\text{Periode}}\)} wird deshalb mit der Quadratwurzel
aus \textbf{\(m\)} multipliziert. Bei täglichen Börsendaten wird häufig
mit 252 Handelstagen gerechnet. Daraus resultiert die annualisierte
Volatilität \textbf{\(\sigma_{\text{ann}}\)}.\footnote{Hilpisch, 694.}
Die Quadratwurzel-Skalierung setzt insbesondere voraus, dass die
periodischen Renditen nicht seriell korreliert sind.\footnote{Sharpe,
  {„The {Sharpe Ratio} ({Fall} 1994)``}.}

\begin{Shaded}
\begin{Highlighting}[]
\CommentTok{\# Importiert die Numpy{-}Bibliothek}
\ImportTok{import}\NormalTok{ numpy }\ImportTok{as}\NormalTok{ np}

\CommentTok{\# Die Funktion annualisierte\_volatilitaet() berechnet die annualisierte Volatilität aus einer Liste von Werten und der Anzahl der Perioden pro Jahr}
\KeywordTok{def}\NormalTok{ annualisierte\_volatilitaet(werte: np.ndarray, perioden\_pro\_jahr: }\BuiltInTok{int}\NormalTok{):}
    \CommentTok{\# Die Standardabweichung der Werte wird mit der Quadratwurzel der Anzahl der Perioden pro Jahr multipliziert}
\NormalTok{    annualisierte\_volatilitaet }\OperatorTok{=}\NormalTok{ standardabweichung(werte) }\OperatorTok{*}\NormalTok{ perioden\_pro\_jahr }\OperatorTok{**}\NormalTok{ (}\DecValTok{1}\OperatorTok{/}\DecValTok{2}\NormalTok{)}
    \ControlFlowTok{return}\NormalTok{ annualisierte\_volatilitaet}

\CommentTok{\# Beispiel zur Nutzung der Funktion}
\NormalTok{werte }\OperatorTok{=}\NormalTok{ np.array([}\FloatTok{1.1}\NormalTok{, }\FloatTok{1.2}\NormalTok{, }\FloatTok{0.9}\NormalTok{, }\FloatTok{1.05}\NormalTok{])}
\NormalTok{perioden\_pro\_jahr }\OperatorTok{=} \DecValTok{12}  \CommentTok{\# zum Beispiel monatliche Renditen}
\BuiltInTok{print}\NormalTok{(}\SpecialStringTok{f"Annualisierte Volatilität: }\SpecialCharTok{\{}\NormalTok{annualisierte\_volatilitaet(werte, perioden\_pro\_jahr)}\SpecialCharTok{\}}\SpecialStringTok{"}\NormalTok{)}
\end{Highlighting}
\end{Shaded}

\begin{verbatim}
Annualisierte Volatilität: 0.4330127018922192
\end{verbatim}

3.4.4 Interpretation

Eine tiefe Volatilität bedeutet, dass die periodischen Renditen nahe bei
ihrem Mittelwert liegen. Die Wertentwicklung verläuft somit
vergleichsweise gleichmässig. Eine hohe Volatilität zeigt dagegen, dass
es stärkere positive und negative Ausschläge gab. Eine hohe Volatilität
bedeutet nicht automatisch, dass eine Anlage eine schlechte Rendite
erzielt. Es deutet jedoch darauf hin, dass der Preis während dieser Zeit
stärkere Schwankungen durchgemacht hat.\footnote{Kommer, \emph{Souverän
  investieren mit Indexfonds und ETFs}, 7. Auflage (Weinheim: Campus
  Verlag, 2025), 147--49.}

3.4.5 Diversifikation

Diversifikation kann die Volatilität eines Gesamtportfolios senken. Bei
Marktschwankungen können sich zwei Kurse gegenseitig ausgleichen. Ein
breit diversifiziertes Portfolio zeigt meist eine tiefere Volatilität,
sofern die Renditen der Wertpapiere nicht untereinander
zusammenhängen.\footnote{Kommer, 148.}

Die Volatilität kennt jedoch auch Grenzen. Da sowohl positive als auch
negative Ausschläge gemessen werden, wird nicht zwischen unerwünschten
Verlusten und erwünschten Gewinnen unterschieden. Ausserdem sagt sie
nichts über die Stärke der Verluste und die Länge der Verlustphase aus.
Sie misst nur einen Risikotyp und deckt nicht das ganze Risikospektrum
ab. Deswegen kann man die Volatilität nicht als einzigen Faktor nutzen,
um das Gesamtrisiko zu beurteilen.\footnote{Kommer, 151.}

3.4.6 Beispiel

\begin{Shaded}
\begin{Highlighting}[]
\CommentTok{\# Importiert die Numpy{-}Bibliothek}
\ImportTok{import}\NormalTok{ numpy }\ImportTok{as}\NormalTok{ np}
\CommentTok{\# Importiert die Pandas{-}Bibliothek}
\ImportTok{import}\NormalTok{ pandas }\ImportTok{as}\NormalTok{ pd}
\CommentTok{\# Importiert die Matplotlib{-}Bibliothek für die Visualisierung}
\ImportTok{import}\NormalTok{ matplotlib.pyplot }\ImportTok{as}\NormalTok{ plt}


\CommentTok{\# Monatliche Renditen der weniger volatilen Strategie}
\NormalTok{strategie\_a }\OperatorTok{=}\NormalTok{ pd.Series([}\FloatTok{0.008}\NormalTok{, }\FloatTok{0.012}\NormalTok{, }\FloatTok{0.009}\NormalTok{, }\FloatTok{0.011}\NormalTok{, }\FloatTok{0.007}\NormalTok{, }\FloatTok{0.013}\NormalTok{, }\FloatTok{0.010}\NormalTok{, }\FloatTok{0.010}\NormalTok{, }\FloatTok{0.006}\NormalTok{, }\FloatTok{0.014}\NormalTok{, }\FloatTok{0.009}\NormalTok{, }\FloatTok{0.011}\NormalTok{])}


\CommentTok{\# Monatliche Renditen der volatilen Strategie}
\NormalTok{strategie\_b }\OperatorTok{=}\NormalTok{ pd.Series([}\FloatTok{0.05}\NormalTok{, }\OperatorTok{{-}}\FloatTok{0.03}\NormalTok{, }\FloatTok{0.04}\NormalTok{, }\OperatorTok{{-}}\FloatTok{0.02}\NormalTok{, }\FloatTok{0.06}\NormalTok{, }\OperatorTok{{-}}\FloatTok{0.04}\NormalTok{, }\FloatTok{0.03}\NormalTok{, }\OperatorTok{{-}}\FloatTok{0.01}\NormalTok{, }\FloatTok{0.07}\NormalTok{, }\OperatorTok{{-}}\FloatTok{0.05}\NormalTok{, }\FloatTok{0.02}\NormalTok{, }\FloatTok{0.00}\NormalTok{])}


\CommentTok{\# Durchschnittliche Monatsrenditen berechnen}
\NormalTok{mittelwert\_a }\OperatorTok{=}\NormalTok{ strategie\_a.mean()}
\NormalTok{mittelwert\_b }\OperatorTok{=}\NormalTok{ strategie\_b.mean()}


\CommentTok{\# Annualisierte Volatilitäten berechnen}
\NormalTok{volatilitaet\_a }\OperatorTok{=}\NormalTok{ annualisierte\_volatilitaet(strategie\_a, }\DecValTok{12}\NormalTok{)}
\NormalTok{volatilitaet\_b }\OperatorTok{=}\NormalTok{ annualisierte\_volatilitaet(strategie\_b, }\DecValTok{12}\NormalTok{)}


\CommentTok{\# Ergebnisse ausgeben}
\BuiltInTok{print}\NormalTok{(}\StringTok{"Strategie A"}\NormalTok{)}
\BuiltInTok{print}\NormalTok{(}\SpecialStringTok{f"Durchschnittliche Monatsrendite: }\SpecialCharTok{\{}\NormalTok{mittelwert\_a}\SpecialCharTok{:.2\%\}}\SpecialStringTok{"}\NormalTok{)}
\BuiltInTok{print}\NormalTok{(}\SpecialStringTok{f"Annualisierte Volatilität: }\SpecialCharTok{\{}\NormalTok{volatilitaet\_a}\SpecialCharTok{:.2\%\}}\SpecialStringTok{"}\NormalTok{)}

\BuiltInTok{print}\NormalTok{()}

\BuiltInTok{print}\NormalTok{(}\StringTok{"Strategie B"}\NormalTok{)}
\BuiltInTok{print}\NormalTok{(}\SpecialStringTok{f"Durchschnittliche Monatsrendite: }\SpecialCharTok{\{}\NormalTok{mittelwert\_b}\SpecialCharTok{:.2\%\}}\SpecialStringTok{"}\NormalTok{)}
\BuiltInTok{print}\NormalTok{(}\SpecialStringTok{f"Annualisierte Volatilität: }\SpecialCharTok{\{}\NormalTok{volatilitaet\_b}\SpecialCharTok{:.2\%\}}\SpecialStringTok{"}\NormalTok{)}


\CommentTok{\# Monatsnummern für die Abbildung}
\NormalTok{monate }\OperatorTok{=} \BuiltInTok{range}\NormalTok{(}\DecValTok{1}\NormalTok{, }\DecValTok{13}\NormalTok{)}


\CommentTok{\# Renditen als Prozentwerte darstellen}
\NormalTok{plt.figure(figsize}\OperatorTok{=}\NormalTok{(}\DecValTok{10}\NormalTok{, }\DecValTok{5}\NormalTok{))}

\NormalTok{plt.plot(}
\NormalTok{    monate,}
\NormalTok{    strategie\_a }\OperatorTok{*} \DecValTok{100}\NormalTok{,}
\NormalTok{    marker}\OperatorTok{=}\StringTok{"o"}\NormalTok{,}
\NormalTok{    label}\OperatorTok{=}\StringTok{"Strategie A"}
\NormalTok{)}

\NormalTok{plt.plot(}
\NormalTok{    monate,}
\NormalTok{    strategie\_b }\OperatorTok{*} \DecValTok{100}\NormalTok{,}
\NormalTok{    marker}\OperatorTok{=}\StringTok{"o"}\NormalTok{,}
\NormalTok{    label}\OperatorTok{=}\StringTok{"Strategie B"}
\NormalTok{)}

\CommentTok{\# Linie bei null Prozent}
\NormalTok{plt.axhline(}
\NormalTok{    y}\OperatorTok{=}\DecValTok{0}\NormalTok{,}
\NormalTok{    linewidth}\OperatorTok{=}\DecValTok{1}
\NormalTok{)}

\NormalTok{plt.xlabel(}\StringTok{"Monat"}\NormalTok{)}
\NormalTok{plt.ylabel(}\StringTok{"Monatliche Rendite in \%"}\NormalTok{)}
\NormalTok{plt.title(}
    \StringTok{"Ähnliche Durchschnittsrendite bei unterschiedlicher Volatilität"}
\NormalTok{)}

\NormalTok{plt.xticks(monate)}
\NormalTok{plt.legend()}
\NormalTok{plt.grid(alpha}\OperatorTok{=}\FloatTok{0.3}\NormalTok{)}
\NormalTok{plt.show()}
\end{Highlighting}
\end{Shaded}

\begin{verbatim}
Strategie A
Durchschnittliche Monatsrendite: 1.00%
Annualisierte Volatilität: 0.82%

Strategie B
Durchschnittliche Monatsrendite: 1.00%
Annualisierte Volatilität: 14.09%
\end{verbatim}

\begin{center}
\pandocbounded{\includegraphics[keepaspectratio]{notebooks/Theorie_files/figure-pdf/cell-9-output-2.png}}
\end{center}

Das Beispiel zeigt übertrieben, wie ein Kurs bei gleicher Rendite eine
ganz unterschiedliche Volatilität aufweisen kann.

3.5 Drawdown

3.5.1 Bedeutung des Drawdowns

Die Volatilität misst die allgemeine Streuung der Renditen, zeigt aber
nicht direkt, wie tief eine Anlage während einer Verlustphase unter
ihren bisherigen Höchststand fällt. Ein Drawdown vergleicht deshalb den
aktuellen Vermögenswert mit dem höchsten Wert, der bis zu diesem
Zeitpunkt erreicht wurde. An einem neuen Höchststand beträgt der
Drawdown 0\%. Liegt der Vermögenswert unter dem bisherigen Höchststand,
ist der Drawdown negativ.\footnote{Kommer, 151--52.}

Der Maximum Drawdown ist der grösste kumulierte historische Verlust
innerhalb des untersuchten Zeitraums. Er zeigt den schlechtesten
beobachteten Rückgang von einem früheren Höchststand bis zum
darauffolgenden Tiefststand.\footnote{Kommer, 151--53.}

3.5.2 Berechnung

Formel für den laufenden Drawdown:

\[
DD_t = \frac{V_t}{\max(V_0,\ldots,V_t)}-1
\]

Quelle der relativen Drawdown-Formel:.\footnote{Pham, Cui, und Ruthbah,
  \emph{The {Performance} of the 60/40 {Portfolio}}, 51.}

Dabei bezeichnet Vt den Vermögenswert zum Zeitpunkt t. Im Nenner steht
der höchste Vermögenswert, der bis zu diesem Zeitpunkt beobachtet wurde.
DDt steht für den Drawdown zum Zeitpunkt t. Der Maximum Drawdown ist der
kleinste und somit stärkste negative Wert der gesamten
Drawdown-Reihe.\footnote{Pham, Cui, und Ruthbah, 51.}

Formel für den Maximum Drawdown:

\[
MDD = \min_t(DD_t)
\]

Quelle der Formel:.\footnote{Pham, Cui, und Ruthbah, 51.}

\begin{Shaded}
\begin{Highlighting}[]
\CommentTok{\# Importiert die Pandas{-}Bibliothek}
\ImportTok{import}\NormalTok{ pandas }\ImportTok{as}\NormalTok{ pd}

\CommentTok{\# Die Funktion drawdown() berechnet den laufenden Drawdown aus periodischen Renditen}
\KeywordTok{def}\NormalTok{ drawdown(renditen: pd.Series):}
    \CommentTok{\# Aus den Renditen wird zunächst mit + 1 der Wachstumsfaktor gebildet}
    \CommentTok{\# Anschliessend werden mit .cumprod() alle Renditen nacheinander multipliziert, damit man den Zinseszinseffekt erhält}
\NormalTok{    vermoegen }\OperatorTok{=}\NormalTok{ (}\DecValTok{1} \OperatorTok{+}\NormalTok{ renditen).cumprod()}
    \CommentTok{\# cummax() bestimmt den bis dahin höchsten Vermögenswert}
\NormalTok{    bisheriger\_hoechststand }\OperatorTok{=}\NormalTok{ vermoegen.cummax()}
    \CommentTok{\# vermoegen / bisheriger\_hoechststand berechnet den Wachstumsfaktor im Verhältnis zum höchsten Vermögenswert}
    \CommentTok{\# Durch {-} 1 erhält man den prozentualen Unterschied zum höchsten Vermögenswert als Dezimalzahl}
    \ControlFlowTok{return}\NormalTok{ vermoegen }\OperatorTok{/}\NormalTok{ bisheriger\_hoechststand }\OperatorTok{{-}} \DecValTok{1}

\CommentTok{\# Die Funktion maximum\_drawdown() gibt den tiefsten Wert der Drawdown{-}Reihe zurück}
\KeywordTok{def}\NormalTok{ maximum\_drawdown(renditen: pd.Series):}
    \ControlFlowTok{return}\NormalTok{ drawdown(renditen).}\BuiltInTok{min}\NormalTok{()}

\CommentTok{\# Beispiel zur Nutzung der Funktionen}
\NormalTok{beispiel\_renditen }\OperatorTok{=}\NormalTok{ pd.Series([}\FloatTok{0.10}\NormalTok{, }\OperatorTok{{-}}\FloatTok{0.05}\NormalTok{, }\OperatorTok{{-}}\FloatTok{0.20}\NormalTok{, }\FloatTok{0.08}\NormalTok{, }\FloatTok{0.15}\NormalTok{, }\FloatTok{0.04}\NormalTok{])}

\BuiltInTok{print}\NormalTok{(}\SpecialStringTok{f"Drawdown{-}Reihe:}\CharTok{\textbackslash{}n}\SpecialCharTok{\{}\NormalTok{drawdown(beispiel\_renditen)}\SpecialCharTok{\}}\CharTok{\textbackslash{}n}\SpecialStringTok{"}\NormalTok{)}
\BuiltInTok{print}\NormalTok{(}\SpecialStringTok{f"Maximum Drawdown: }\SpecialCharTok{\{}\NormalTok{maximum\_drawdown(beispiel\_renditen)}\SpecialCharTok{:.2\%\}}\SpecialStringTok{"}\NormalTok{) }\CommentTok{\# Von einer Dezimalzahl zu Prozent}
\end{Highlighting}
\end{Shaded}

\begin{verbatim}
Drawdown-Reihe:
0    0.000000
1   -0.050000
2   -0.240000
3   -0.179200
4   -0.056080
5   -0.018323
dtype: float64

Maximum Drawdown: -24.00%
\end{verbatim}

3.5.3 Erholungsdauer und Grenzen

Die Erholungsdauer beschreibt die Zeit vom früheren Höchststand, bis
dieser Wert vollständig wieder erreicht wird. Die Verlusttiefe und die
Verlustdauer sind zwei verschiedene Eigenschaften: Ein Drawdown kann
tief, aber kurz oder weniger tief, aber sehr lang sein.\footnote{Kommer,
  \emph{Souverän investieren mit Indexfonds und ETFs}, 7. Auflage
  (Weinheim: Campus Verlag, 2025), 163--64.}

Für die Erholungsdauer und alle wichtigen Daten dazu nutzen wir zwei
Funktionen aus QuantStats: .to\_drawdown\_series() und
.drawdown\_details().\footnote{Aroussi, \emph{Ranaroussi/Quantstats}.}

\begin{Shaded}
\begin{Highlighting}[]
\CommentTok{\# Importiert die QuantStats{-}Bibliothek}
\ImportTok{import}\NormalTok{ quantstats }\ImportTok{as}\NormalTok{ qs}
\CommentTok{\# Importiert die Pandas{-}Bibliothek}
\ImportTok{import}\NormalTok{ pandas }\ImportTok{as}\NormalTok{ pd}

\NormalTok{beispiel\_renditen }\OperatorTok{=}\NormalTok{ pd.Series(\{}
    \StringTok{"2024{-}01{-}31"}\NormalTok{: }\FloatTok{0.10}\NormalTok{,}
    \StringTok{"2024{-}02{-}29"}\NormalTok{: }\OperatorTok{{-}}\FloatTok{0.05}\NormalTok{,}
    \StringTok{"2024{-}03{-}31"}\NormalTok{: }\OperatorTok{{-}}\FloatTok{0.05}\NormalTok{,}
    \StringTok{"2024{-}04{-}30"}\NormalTok{: }\FloatTok{0.12}\NormalTok{,}
    \StringTok{"2024{-}05{-}31"}\NormalTok{: }\FloatTok{0.03}\NormalTok{,}
    \StringTok{"2024{-}06{-}30"}\NormalTok{: }\OperatorTok{{-}}\FloatTok{0.10}\NormalTok{,}
    \StringTok{"2024{-}07{-}31"}\NormalTok{: }\FloatTok{0.04}\NormalTok{,}
    \StringTok{"2024{-}08{-}31"}\NormalTok{: }\FloatTok{0.08}\NormalTok{,}
    \StringTok{"2024{-}09{-}30"}\NormalTok{: }\OperatorTok{{-}}\FloatTok{0.06}\NormalTok{,}
    \StringTok{"2024{-}10{-}31"}\NormalTok{: }\OperatorTok{{-}}\FloatTok{0.04}\NormalTok{,}
    \StringTok{"2024{-}11{-}30"}\NormalTok{: }\FloatTok{0.05}\NormalTok{,}
    \StringTok{"2024{-}12{-}31"}\NormalTok{: }\FloatTok{0.03}
\NormalTok{\})}

\NormalTok{drawdown\_reihe }\OperatorTok{=}\NormalTok{ qs.stats.to\_drawdown\_series(beispiel\_renditen)}
\NormalTok{erholungsphasen }\OperatorTok{=}\NormalTok{ qs.stats.drawdown\_details(drawdown\_reihe)}
\BuiltInTok{print}\NormalTok{(}\SpecialStringTok{f"Drawdown{-}Reihe:}\CharTok{\textbackslash{}n}\SpecialCharTok{\{}\NormalTok{drawdown\_reihe}\SpecialCharTok{\}}\CharTok{\textbackslash{}n}\SpecialStringTok{"}\NormalTok{)}
\BuiltInTok{print}\NormalTok{(}\SpecialStringTok{f"Erholungsphasen:}\CharTok{\textbackslash{}n}\SpecialCharTok{\{}\NormalTok{erholungsphasen}\SpecialCharTok{\}}\CharTok{\textbackslash{}n}\SpecialStringTok{"}\NormalTok{)}
\end{Highlighting}
\end{Shaded}

\begin{verbatim}
Drawdown-Reihe:
2024-01-31    0.000000
2024-02-29   -0.050000
2024-03-31   -0.097500
2024-04-30    0.000000
2024-05-31    0.000000
2024-06-30   -0.100000
2024-07-31   -0.064000
2024-08-31    0.000000
2024-09-30   -0.060000
2024-10-31   -0.097600
2024-11-30   -0.052480
2024-12-31   -0.024054
dtype: float64

Erholungsphasen:
        start      valley         end  days  max drawdown  99% max drawdown
0  2024-02-29  2024-03-31  2024-03-31    32         -9.75              -5.0
1  2024-06-30  2024-06-30  2024-07-31    32        -10.00              -6.4
2  2024-09-30  2024-10-31  2024-12-31    93         -9.76              -6.0
\end{verbatim}

Der berechnete Maximum Drawdown hängt vom gewählten Zeitraum und von der
Frequenz der Daten ab. Werden nur Jahresendwerte verwendet, können
zwischenzeitliche Verluste übersehen und der tatsächliche Drawdown
deutlich unterschätzt werden. Kommer empfiehlt deshalb einen langen
Beobachtungszeitraum und weist auf die genauere Erfassung mit
Monatsdaten hin. Der historische Maximum Drawdown beschreibt nur den
schlimmsten Verlust in der untersuchten Stichprobe und stellt keine
garantierte Verlustgrenze für die Zukunft dar.\footnote{Kommer,
  \emph{Souverän investieren mit Indexfonds und ETFs}, 7. Auflage
  (Weinheim: Campus Verlag, 2025), 152--53.}

3.6 Sharpe Ratio

3.6.1 Bedeutung der Sharpe Ratio

Die Rendite allein zeigt nicht, mit welchen Schwankungen ein
Anlageergebnis erreicht wurde. Für einen risikobereinigten Vergleich
wird deshalb zunächst die risikofreie Rendite von der Anlagenrendite
abgezogen. Die Sharpe Ratio stellt den Durchschnitt dieser
Überschussrenditen der Volatilität gegenüber.\footnote{Kommer, 165--66.}

Das Ergebnis lässt sich als historische Entschädigung für eine Einheit
des gemessenen Schwankungsrisikos interpretieren. Unter ansonsten
gleichen Berechnungsbedingungen weist ein höherer Wert auf ein
vorteilhafteres Verhältnis zwischen Überschussrendite und Volatilität
hin.\footnote{Kommer, 166.}

3.6.2 Berechnung

Formel für die annualisierte Sharpe Ratio:

\[
SR_{\mathrm{ann}} = \frac{\overline{r-r_f}}{\sigma(r-r_f)}\sqrt{m}
\]

Die periodische Sharpe Ratio ist bei Pham et al.~angegeben.\footnote{Pham,
  Cui, und Ruthbah, \emph{The {Performance} of the 60/40 {Portfolio}},
  51.} Die Annualisierung mit der Quadratwurzel der Periodenzahl setzt
insbesondere unkorrelierte periodische Überschussrenditen
voraus.\footnote{Sharpe, {„The {Sharpe Ratio} ({Fall} 1994)``}.}

Dabei bezeichnet \textbf{\(r\)} die periodische Anlagenrendite,
\textbf{\(r_f\)} die risikofreie Rendite derselben Periode,
\textbf{\(\sigma(r-r_f)\)} die Standardabweichung der periodischen
Überschussrenditen und \textbf{\(m\)} die Anzahl der Perioden pro Jahr.

\begin{Shaded}
\begin{Highlighting}[]
\CommentTok{\# Importiert die Pandas{-}Bibliothek}
\ImportTok{import}\NormalTok{ pandas }\ImportTok{as}\NormalTok{ pd}
\CommentTok{\# Importiert die QuantStats{-}Bibliothek}
\ImportTok{import}\NormalTok{ quantstats }\ImportTok{as}\NormalTok{ qs}


\CommentTok{\# Die Funktion sharpe\_ratio() berechnet die annualisierte Sharpe Ratio}
\KeywordTok{def}\NormalTok{ sharpe\_ratio(renditen: pd.Series, risikofreie\_renditen: pd.Series, perioden\_pro\_jahr: }\BuiltInTok{int}\NormalTok{):}
\NormalTok{    daten }\OperatorTok{=}\NormalTok{ pd.concat(}
\NormalTok{        [renditen, risikofreie\_renditen],}
\NormalTok{        axis}\OperatorTok{=}\DecValTok{1}
\NormalTok{    )}

\NormalTok{    daten.columns }\OperatorTok{=}\NormalTok{ [}\StringTok{"Anlage"}\NormalTok{, }\StringTok{"Risikofrei"}\NormalTok{]}

    \CommentTok{\# Berechnet die periodischen Überschussrenditen}
\NormalTok{    ueberschussrenditen }\OperatorTok{=}\NormalTok{ (}
\NormalTok{        daten[}\StringTok{"Anlage"}\NormalTok{]}
        \OperatorTok{{-}}\NormalTok{ daten[}\StringTok{"Risikofrei"}\NormalTok{]}
\NormalTok{    )}

    \CommentTok{\# Berechnet und annualisiert die Sharpe Ratio mit QuantStats}
    \ControlFlowTok{return}\NormalTok{ qs.stats.sharpe(}
\NormalTok{        ueberschussrenditen,}
\NormalTok{        rf}\OperatorTok{=}\DecValTok{0}\NormalTok{,}
\NormalTok{        periods}\OperatorTok{=}\NormalTok{perioden\_pro\_jahr}
\NormalTok{    )}


\NormalTok{anlage }\OperatorTok{=}\NormalTok{ pd.Series([}\FloatTok{0.02}\NormalTok{, }\OperatorTok{{-}}\FloatTok{0.01}\NormalTok{, }\FloatTok{0.03}\NormalTok{, }\FloatTok{0.01}\NormalTok{, }\OperatorTok{{-}}\FloatTok{0.02}\NormalTok{, }\FloatTok{0.025}\NormalTok{])}

\NormalTok{risikofrei }\OperatorTok{=}\NormalTok{ pd.Series([}\FloatTok{0.002}\NormalTok{, }\FloatTok{0.002}\NormalTok{, }\FloatTok{0.002}\NormalTok{, }\FloatTok{0.002}\NormalTok{, }\FloatTok{0.002}\NormalTok{, }\FloatTok{0.002}\NormalTok{])}

\NormalTok{ergebnis }\OperatorTok{=}\NormalTok{ sharpe\_ratio(}
\NormalTok{    anlage,}
\NormalTok{    risikofrei,}
\NormalTok{    perioden\_pro\_jahr}\OperatorTok{=}\DecValTok{12}
\NormalTok{)}

\BuiltInTok{print}\NormalTok{(}\SpecialStringTok{f"Annualisierte Sharpe Ratio: }\SpecialCharTok{\{}\NormalTok{ergebnis}\SpecialCharTok{:.2f\}}\SpecialStringTok{"}\NormalTok{)}
\end{Highlighting}
\end{Shaded}

\begin{verbatim}
Annualisierte Sharpe Ratio: 1.23
\end{verbatim}

3.6.3 Interpretation und Grenzen

Sharpe Ratios können nur sinnvoll verglichen werden, wenn Zeitraum,
Datenfrequenz und Berechnungsweise identisch sind. Eine Sharpe Ratio für
einen Zeitraum darf nicht direkt mit einer Kennzahl aus einem anderen
Zeitraum verglichen werden.\footnote{Kommer, \emph{Souverän investieren
  mit Indexfonds und ETFs}, 7. Auflage (Weinheim: Campus Verlag, 2025),
  166.}

Neben der vollständigen Variante existiert eine vereinfachte Raw Sharpe
Ratio, bei welcher die risikofreie Rendite nicht abgezogen wird. Beide
Varianten dürfen nicht vermischt werden.\footnote{Kommer, 165--66.}
Ausserdem bildet die Sharpe Ratio das Risiko nur über die Volatilität
ab. Da die Volatilität nicht das gesamte Anlagerisiko beschreibt, werden
Rendite und Drawdown in dieser Arbeit zusätzlich als getrennte
Kennzahlen betrachtet.\footnote{Kommer, 166.}

3.7 Diversifikation und Korrelation

3.7.1 Diversifikation

Bei der Diversifikation wird das Kapital nicht auf eine einzelne Anlage
konzentriert. Eine Streuung ist sowohl innerhalb derselben Anlageklasse,
etwa über viele verschiedene Aktien, als auch zwischen Anlageklassen wie
Aktien und Anleihen möglich.\footnote{Kommer, 82.}

Auf diese Weise lassen sich Risiken, die nur einzelne Unternehmen oder
Branchen betreffen, erheblich verringern. Marktweite Kursschwankungen
wirken dagegen auf viele Titel gleichzeitig und verschwinden auch in
einem breit gestreuten Aktienportfolio nicht vollständig.\footnote{Kommer,
  85--86.} Für die Beurteilung ist deshalb entscheidend, wie die
Bestandteile im Gesamtportfolio zusammenwirken.\footnote{Kommer, 85.}

3.7.2 Korrelation

Mit der Korrelation lässt sich untersuchen, wie eng die Renditen zweier
Anlagen linear zusammenhängen. Ihr Koeffizient reicht von -1 bis +1: Am
positiven Grenzwert bewegen sich beide Reihen vollständig in dieselbe
Richtung, am negativen Grenzwert genau entgegengesetzt. Ein Wert von 0
zeigt, dass kein linearer Zusammenhang nachweisbar ist.\footnote{Kommer,
  82--83.}

Ein Portfolio kann bereits von einer Korrelation unter +1 profitieren,
da sich die Wertveränderungen seiner Bestandteile dann nicht vollkommen
parallel entwickeln. Eine negative Korrelation ist somit keine
Voraussetzung. Tendenziell nimmt die risikomindernde Wirkung der
Kombination zu, wenn der Korrelationskoeffizient sinkt.\footnote{Kommer,
  82--83.}

Formel für die Korrelation:

\[
\rho_{X,Y}=\frac{\operatorname{Cov}(X,Y)}{\sigma_X\sigma_Y}
\] Dabei bezeichnet \textbf{\(\rho_{X,Y}\)} die Korrelation zwischen den
Variablen \textbf{\(X\)} und \textbf{\(Y\)},
\textbf{\(\operatorname{Cov}(X,Y)\)} ihre Kovarianz,
\textbf{\(\sigma_X\)} die Standardabweichung von \textbf{\(X\)} und
\textbf{\(\sigma_Y\)} die Standardabweichung von
\textbf{\(Y\)}.\footnote{Markowitz, {„Portfolio {Selection}``}, 80.}
Portfoliovarianz bei zwei Anlagen:

\[
\sigma_p^2=w_1^2\sigma_1^2+w_2^2\sigma_2^2+2w_1w_2\sigma_1\sigma_2\rho_{1,2}
\] Dabei bezeichnet \textbf{\(\sigma_p^2\)} die Varianz des Portfolios,
\textbf{\(w_1\)} und \textbf{\(w_2\)} die Gewichtungen der beiden
Anlagen, \textbf{\(\sigma_1\)} und \textbf{\(\sigma_2\)} deren
Standardabweichungen und \textbf{\(\rho_{1,2}\)} die Korrelation
zwischen den Renditen der beiden Anlagen.\footnote{Markowitz, 80--81;
  Pham, Cui, und Ruthbah, \emph{The {Performance} of the 60/40
  {Portfolio}}, 51.}

\begin{Shaded}
\begin{Highlighting}[]
\CommentTok{\# Importiert die Pandas{-}Bibliothek}
\ImportTok{import}\NormalTok{ pandas }\ImportTok{as}\NormalTok{ pd}

\CommentTok{\# Die Funktion korrelationsmatrix() berechnet die Korrelationsmatrix der Renditen}
\KeywordTok{def}\NormalTok{ korrelationsmatrix(renditen: pd.DataFrame):}
    \CommentTok{\# corr() berechnet die Korrelationsmatrix}
    \ControlFlowTok{return}\NormalTok{ renditen.corr()}

\CommentTok{\# Beispiel zur Nutzung der Funktion}
\NormalTok{renditen }\OperatorTok{=}\NormalTok{ pd.DataFrame(\{}
    \StringTok{"Aktien"}\NormalTok{: [}\FloatTok{0.04}\NormalTok{, }\OperatorTok{{-}}\FloatTok{0.03}\NormalTok{, }\FloatTok{0.02}\NormalTok{, }\FloatTok{0.05}\NormalTok{, }\OperatorTok{{-}}\FloatTok{0.01}\NormalTok{, }\FloatTok{0.03}\NormalTok{],}
    \StringTok{"Anleihen"}\NormalTok{: [}\FloatTok{0.01}\NormalTok{, }\FloatTok{0.00}\NormalTok{, }\OperatorTok{{-}}\FloatTok{0.01}\NormalTok{, }\FloatTok{0.01}\NormalTok{, }\FloatTok{0.02}\NormalTok{, }\FloatTok{0.00}\NormalTok{],}
    \StringTok{"ETF"}\NormalTok{: [}\FloatTok{0.03}\NormalTok{, }\OperatorTok{{-}}\FloatTok{0.02}\NormalTok{, }\FloatTok{0.01}\NormalTok{, }\FloatTok{0.04}\NormalTok{, }\FloatTok{0.00}\NormalTok{, }\FloatTok{0.02}\NormalTok{]}
\NormalTok{\})}

\NormalTok{korrelation }\OperatorTok{=}\NormalTok{ korrelationsmatrix(renditen)}
\BuiltInTok{print}\NormalTok{(}\SpecialStringTok{f"Korrelationsmatrix:}\CharTok{\textbackslash{}n}\SpecialCharTok{\{}\NormalTok{korrelation}\SpecialCharTok{\}}\CharTok{\textbackslash{}n}\SpecialStringTok{"}\NormalTok{)}
\end{Highlighting}
\end{Shaded}

\begin{verbatim}
Korrelationsmatrix:
                Aktien      Anleihen       ETF
Aktien    1.000000e+00  1.679908e-17  0.982953
Anleihen  1.679908e-17  1.000000e+00  0.176547
ETF       9.829533e-01  1.765470e-01  1.000000
\end{verbatim}

\begin{Shaded}
\begin{Highlighting}[]
\CommentTok{\# Importiert die Pandas{-}Bibliothek}
\ImportTok{import}\NormalTok{ pandas }\ImportTok{as}\NormalTok{ pd}
\CommentTok{\# Importiert die Numpy{-}Bibliothek}
\ImportTok{import}\NormalTok{ numpy }\ImportTok{as}\NormalTok{ np}

\CommentTok{\# Die Funktion portfolio\_risiko() berechnet die Varianz und Volatilität eines Portfolios anhand der Renditen und Gewichte der einzelnen Anlagen}
\KeywordTok{def}\NormalTok{ portfolio\_risiko(renditen, gewichte):}
\NormalTok{    cov }\OperatorTok{=}\NormalTok{ renditen.cov()}
\NormalTok{    varianz }\OperatorTok{=}\NormalTok{ gewichte }\OperatorTok{@}\NormalTok{ cov }\OperatorTok{@}\NormalTok{ gewichte}
\NormalTok{    volatilitaet }\OperatorTok{=}\NormalTok{ np.sqrt(varianz)}
    \ControlFlowTok{return}\NormalTok{ varianz, volatilitaet}

\CommentTok{\# Beispiel zur Nutzung der Funktion}
\NormalTok{renditen }\OperatorTok{=}\NormalTok{ pd.DataFrame(\{}
    \StringTok{"Aktien"}\NormalTok{: [}\FloatTok{0.04}\NormalTok{, }\OperatorTok{{-}}\FloatTok{0.03}\NormalTok{, }\FloatTok{0.02}\NormalTok{, }\FloatTok{0.05}\NormalTok{, }\OperatorTok{{-}}\FloatTok{0.01}\NormalTok{, }\FloatTok{0.03}\NormalTok{],}
    \StringTok{"Anleihen"}\NormalTok{: [}\FloatTok{0.01}\NormalTok{, }\FloatTok{0.00}\NormalTok{, }\OperatorTok{{-}}\FloatTok{0.01}\NormalTok{, }\FloatTok{0.01}\NormalTok{, }\FloatTok{0.02}\NormalTok{, }\FloatTok{0.00}\NormalTok{],}
    \StringTok{"ETF"}\NormalTok{: [}\FloatTok{0.03}\NormalTok{, }\OperatorTok{{-}}\FloatTok{0.02}\NormalTok{, }\FloatTok{0.01}\NormalTok{, }\FloatTok{0.04}\NormalTok{, }\FloatTok{0.00}\NormalTok{, }\FloatTok{0.02}\NormalTok{]}
\NormalTok{\})}
\NormalTok{gewichte }\OperatorTok{=}\NormalTok{ pd.Series(\{}\StringTok{"Aktien"}\NormalTok{: }\FloatTok{0.20}\NormalTok{, }\StringTok{"Anleihen"}\NormalTok{: }\FloatTok{0.40}\NormalTok{, }\StringTok{"ETF"}\NormalTok{: }\FloatTok{0.20}\NormalTok{\})}
\NormalTok{var, vol }\OperatorTok{=}\NormalTok{ portfolio\_risiko(renditen, gewichte)}
\BuiltInTok{print}\NormalTok{(}\SpecialStringTok{f"Volatilität: }\SpecialCharTok{\{}\NormalTok{vol }\OperatorTok{*} \DecValTok{100}\SpecialCharTok{:.2f\}}\SpecialStringTok{\%"}\NormalTok{)}
\end{Highlighting}
\end{Shaded}

\begin{verbatim}
Volatilität: 1.15%
\end{verbatim}

3.7.3 Veränderliche Korrelationen

Der Zusammenhang zwischen zwei Renditereihen kann sich im Zeitverlauf
verändern. In Phasen fallender Börsenkurse bewegen sich risikobehaftete
Anlagen zeitweise ähnlicher, sodass ihre Kombination ausgerechnet
während einer Krise weniger stabilisierend wirken kann.\footnote{Kommer,
  \emph{Souverän investieren mit Indexfonds und ETFs}, 7. Auflage
  (Weinheim: Campus Verlag, 2025), 84.} Historische Ländervergleiche
zeigen ausserdem, dass Aktien und Anleihen langfristig nicht durchgehend
negativ korrelierten und je nach untersuchtem Markt unterschiedliche
Beziehungen aufwiesen.\footnote{Pham, Cui, und Ruthbah, \emph{The
  {Performance} of the 60/40 {Portfolio}}, 1--7.}

3.8 Asset-Allokation und Rebalancing

3.8.1 Asset-Allokation

Anlagen mit vergleichbaren Rendite-, Risiko- und Liquiditätsmerkmalen
werden zu einer Assetklasse zusammengefasst. Die Asset-Allokation legt
fest, welche dieser Klassen ein Portfolio kombiniert und welchen Anteil
jede davon am Gesamtvermögen erhält.\footnote{Kommer, \emph{Souverän
  investieren mit Indexfonds und ETFs}, 7. Auflage (Weinheim: Campus
  Verlag, 2025), 169.}

Die Portfoliorendite ergibt sich aus den einzelnen Renditebeiträgen, die
jeweils mit dem zugehörigen Portfolioanteil gewichtet werden.\footnote{Pham,
  Cui, und Ruthbah, \emph{The {Performance} of the 60/40 {Portfolio}},
  50.}

\[
r_{p,t}=\sum_{k=1}^{N}w_{k,t}r_{k,t}
\]

Dabei bezeichnet wk,t das Gewicht und rk,t die Rendite der Anlage k zum
Zeitpunkt t.\footnote{Pham, Cui, und Ruthbah, 50.}

3.8.2 Gewichtungsdrift und Rebalancing

Wenn sich die Bestandteile eines Portfolios unterschiedlich entwickeln,
verändern sich ihre Gewichte automatisch. Eine Anlage mit einer höheren
Rendite erhält danach einen grösseren Anteil am Gesamtportfolio. Kommer
zeigt diese Gewichtungsdrift anhand eines Portfolios, dessen
Zusammensetzung bereits nach einem Jahr von den festgelegten
Zielgewichten abweicht.\footnote{Kommer, \emph{Souverän investieren mit
  Indexfonds und ETFs}, 7. Auflage (Weinheim: Campus Verlag, 2025), 523.}

\[
w_{k,t+1}=\frac{w_{k,t}(1+r_{k,t})}{1+r_{p,t}}
\]

Dabei bezeichnet wk,t das Gewicht der Anlage k zu Beginn der Periode,
rk,t ihre Rendite, rp,t die Rendite des Gesamtportfolios und wk,t+1 das
daraus entstehende Gewicht am Ende der Periode.

Diese Formel ist eine eigene algebraische Darstellung der
Gewichtungsdrift, die Kommer anhand eines Zahlenbeispiels
zeigt.\footnote{Kommer, 523--24.}

Rebalancing stellt die ursprünglich festgelegten Zielgewichte wieder
her. Sein Hauptzweck ist die Erhaltung des bewusst gewählten
Rendite-Risiko-Profils. Es handelt sich um regelgebundenes
Risikomanagement und nicht um eine Prognose, welche Anlage als Nächstes
steigen wird.\footnote{Kommer, 524--25.}

Beim kalenderbasierten Rebalancing wird das Portfolio zu festen
Zeitpunkten angepasst. Beim bänderbasierten Rebalancing erfolgt die
Anpassung erst, wenn ein Gewicht eine vorher festgelegte Abweichung
überschreitet.\footnote{Kommer, 525--26.} Für das spätere Modell wird
eine feste jährliche Regel verwendet, damit die Strategie eindeutig
reproduzierbar bleibt.

\begin{Shaded}
\begin{Highlighting}[]
\CommentTok{\# Importiert die Pandas{-}Bibliothek}
\ImportTok{import}\NormalTok{ pandas }\ImportTok{as}\NormalTok{ pd}


\CommentTok{\# Die Funktion neue\_gewichtung() berechnet die Werte und Gewichte nach den Renditen}
\KeywordTok{def}\NormalTok{ neue\_gewichtung(}
\NormalTok{    startwerte: pd.Series,}
\NormalTok{    renditen: pd.Series}
\NormalTok{):}
    \CommentTok{\# Berechnet die neuen Werte nach der Rendite}
\NormalTok{    neue\_werte }\OperatorTok{=}\NormalTok{ startwerte }\OperatorTok{*}\NormalTok{ (}\DecValTok{1} \OperatorTok{+}\NormalTok{ renditen)}
    \CommentTok{\# Berechnet den Gesamtwert des Portfolios}
\NormalTok{    gesamtwert }\OperatorTok{=}\NormalTok{ neue\_werte.}\BuiltInTok{sum}\NormalTok{()}
    \CommentTok{\# Berechnet die neuen Gewichte}
\NormalTok{    neue\_gewichte }\OperatorTok{=}\NormalTok{ neue\_werte }\OperatorTok{/}\NormalTok{ gesamtwert}

    \ControlFlowTok{return}\NormalTok{ neue\_werte, neue\_gewichte}


\CommentTok{\# Beispiel zur Nutzung der Funktion}
\NormalTok{startwerte }\OperatorTok{=}\NormalTok{ pd.Series(\{}
    \StringTok{"Aktien"}\NormalTok{: }\FloatTok{70.0}\NormalTok{,}
    \StringTok{"Anleihen"}\NormalTok{: }\FloatTok{20.0}\NormalTok{,}
    \StringTok{"ETF"}\NormalTok{: }\FloatTok{10.0}
\NormalTok{\})}

\NormalTok{renditen }\OperatorTok{=}\NormalTok{ pd.Series(\{}
    \StringTok{"Aktien"}\NormalTok{: }\FloatTok{0.20}\NormalTok{,}
    \StringTok{"Anleihen"}\NormalTok{: }\FloatTok{0.02}\NormalTok{,}
    \StringTok{"ETF"}\NormalTok{: }\FloatTok{0.05}
\NormalTok{\})}

\NormalTok{start, rend }\OperatorTok{=}\NormalTok{ neue\_gewichtung(startwerte, renditen)}

\BuiltInTok{print}\NormalTok{(pd.DataFrame(\{}
    \StringTok{"Neuer Wert"}\NormalTok{: start,}
    \StringTok{"Neues Gewicht"}\NormalTok{: rend}
\NormalTok{\}))}
\end{Highlighting}
\end{Shaded}

\begin{verbatim}
          Neuer Wert  Neues Gewicht
Aktien          84.0       0.731070
Anleihen        20.4       0.177546
ETF             10.5       0.091384
\end{verbatim}

\begin{Shaded}
\begin{Highlighting}[]
\CommentTok{\# Importiert die Pandas{-}Bibliothek}
\ImportTok{import}\NormalTok{ pandas }\ImportTok{as}\NormalTok{ pd}

\CommentTok{\# Die Funktion rebalancing() berechnet die notwendigen Käufe und Verkäufe}
\KeywordTok{def}\NormalTok{ rebalancing(}
\NormalTok{    aktuelle\_werte: pd.Series,}
\NormalTok{    zielgewichte: pd.Series}
\NormalTok{):}
    \CommentTok{\# Berechnet den aktuellen Gesamtwert des Portfolios}
\NormalTok{    gesamtwert }\OperatorTok{=}\NormalTok{ aktuelle\_werte.}\BuiltInTok{sum}\NormalTok{()}
    \CommentTok{\# Berechnet die gewünschten Werte nach dem Rebalancing}
\NormalTok{    zielwerte }\OperatorTok{=}\NormalTok{ gesamtwert }\OperatorTok{*}\NormalTok{ zielgewichte}
    \CommentTok{\# Berechnet, wie viel gekauft oder verkauft werden muss}
\NormalTok{    transaktionen }\OperatorTok{=}\NormalTok{ zielwerte }\OperatorTok{{-}}\NormalTok{ aktuelle\_werte}
    \CommentTok{\# Berechnet die aktuellen Gewichte}
\NormalTok{    aktuelle\_gewichte }\OperatorTok{=}\NormalTok{ aktuelle\_werte }\OperatorTok{/}\NormalTok{ gesamtwert}
    \ControlFlowTok{return}\NormalTok{ aktuelle\_werte, aktuelle\_gewichte, zielgewichte, zielwerte, transaktionen}


\CommentTok{\# Beispiel zur Nutzung der Funktion}
\NormalTok{aktuelle\_werte }\OperatorTok{=}\NormalTok{ pd.Series(\{}
    \StringTok{"Aktien"}\NormalTok{: }\FloatTok{84.0}\NormalTok{,}
    \StringTok{"Anleihen"}\NormalTok{: }\FloatTok{20.4}\NormalTok{,}
    \StringTok{"ETF"}\NormalTok{: }\FloatTok{10.5}
\NormalTok{\})}

\NormalTok{zielgewichte }\OperatorTok{=}\NormalTok{ pd.Series(\{}
    \StringTok{"Aktien"}\NormalTok{: }\FloatTok{0.70}\NormalTok{,}
    \StringTok{"Anleihen"}\NormalTok{: }\FloatTok{0.20}\NormalTok{,}
    \StringTok{"ETF"}\NormalTok{: }\FloatTok{0.10}
\NormalTok{\})}

\NormalTok{aktuellw, aktuellg, zielg, zielw, tran }\OperatorTok{=}\NormalTok{ rebalancing(}
\NormalTok{    aktuelle\_werte,}
\NormalTok{    zielgewichte}
\NormalTok{)}

\BuiltInTok{print}\NormalTok{(pd.DataFrame(\{}
        \StringTok{"Aktueller Wert"}\NormalTok{: aktuellw,}
        \StringTok{"Aktuelles Gewicht"}\NormalTok{: aktuellg,}
        \StringTok{"Zielgewicht"}\NormalTok{: zielg,}
        \StringTok{"Zielwert"}\NormalTok{: zielw,}
        \StringTok{"Kauf / Verkauf"}\NormalTok{: tran}
\NormalTok{\}))}
\end{Highlighting}
\end{Shaded}

\begin{verbatim}
          Aktueller Wert  Aktuelles Gewicht  Zielgewicht  Zielwert  \
Aktien              84.0           0.731070          0.7     80.43   
Anleihen            20.4           0.177546          0.2     22.98   
ETF                 10.5           0.091384          0.1     11.49   

          Kauf / Verkauf  
Aktien             -3.57  
Anleihen            2.58  
ETF                 0.99  
\end{verbatim}

3.9 Marktbreites Buy-and-Hold

3.9.1 Funktionsweise

Buy-and-Hold bedeutet, ein Portfolio langfristig zu halten und nicht
aufgrund kurzfristiger Marktprognosen laufend zu kaufen und zu
verkaufen. Kommer verbindet diese Strategie mit breit diversifizierten
Indexfonds oder ETFs, festen Regeln, tiefen Kosten und einem
langfristigen Anlagehorizont.\footnote{Kommer, 452--58.}

Marktbreit bedeutet, dass nicht einzelne ausgewählte Unternehmen,
sondern ein grosses Anlageuniversum abgebildet wird. Laufende Erträge
wie Dividenden werden wieder investiert und gehören deshalb zur
Gesamtrendite. Verkäufe können trotzdem für das Rebalancing oder einen
tatsächlichen Kapitalbedarf erfolgen.\footnote{Kommer, 458.}

3.9.2 Vorteile, Risiken und Benchmark

Kommer nennt als wichtige Vorteile von Buy-and-Hold tiefere
Transaktionskosten, einen geringeren Arbeitsaufwand und die Vermeidung
möglicher Timing-Fehler durch häufiges Ein- und Aussteigen.\footnote{Kommer,
  447 und 450.}

Die Strategie verhindert jedoch keine Markteinbrüche. Ein reines
Aktienportfolio nimmt an den Schwankungen, Drawdowns und Erholungsphasen
des Aktienmarkts teil.\footnote{Kommer, 447 und 452-458.} Gerade weil
Buy-and-Hold eine einfache passive Aktienstrategie darstellt, wird es in
dieser Untersuchung als Benchmark für die komplexeren Strategien
verwendet.

\[
V_t=V_0\prod_{i=1}^{t}(1+r_i)
\]

Dabei bezeichnet Vt den Portfoliowert zum Zeitpunkt t, V0 das
Startkapital und ri die Rendite der Periode i. Das Produktzeichen steht
für die fortlaufende Multiplikation aller Wachstumsfaktoren (1 + ri).

Quelle der kumulierten Vermögensformel:.\footnote{Pham, Cui, und
  Ruthbah, \emph{The {Performance} of the 60/40 {Portfolio}}, 51.} Bei
Total-Return-Daten sind die wieder investierten Dividenden bereits in
den periodischen Renditen enthalten.\footnote{Kommer, \emph{Souverän
  investieren mit Indexfonds und ETFs}, 7. Auflage (Weinheim: Campus
  Verlag, 2025), 466--67.}

\begin{Shaded}
\begin{Highlighting}[]
\CommentTok{\# Importiert die Pandas{-}Bibliothek}
\ImportTok{import}\NormalTok{ pandas }\ImportTok{as}\NormalTok{ pd}

\CommentTok{\# Die Funktion buy\_and\_hold() berechnet die Entwicklung eines einmal investierten Startkapitals}
\KeywordTok{def}\NormalTok{ buy\_and\_hold(renditen: pd.Series, startkapital: }\BuiltInTok{float}\NormalTok{):}
    \ControlFlowTok{return}\NormalTok{ startkapital }\OperatorTok{*}\NormalTok{ (}\DecValTok{1} \OperatorTok{+}\NormalTok{ renditen).cumprod()}

\CommentTok{\# Beispiel zur Nutzung der Funktion}
\NormalTok{monatsrenditen }\OperatorTok{=}\NormalTok{ pd.Series([}\FloatTok{0.03}\NormalTok{, }\OperatorTok{{-}}\FloatTok{0.02}\NormalTok{, }\FloatTok{0.01}\NormalTok{, }\FloatTok{0.04}\NormalTok{, }\OperatorTok{{-}}\FloatTok{0.01}\NormalTok{, }\FloatTok{0.02}\NormalTok{])}

\BuiltInTok{print}\NormalTok{(}\SpecialStringTok{f"Vermögensentwicklung:}\CharTok{\textbackslash{}n}\SpecialCharTok{\{}\NormalTok{buy\_and\_hold(monatsrenditen, }\FloatTok{100.0}\NormalTok{)}\SpecialCharTok{\}}\SpecialStringTok{"}\NormalTok{)}
\end{Highlighting}
\end{Shaded}

\begin{verbatim}
Vermögensentwicklung:
0    103.000000
1    100.940000
2    101.949400
3    106.027376
4    104.967102
5    107.066444
dtype: float64
\end{verbatim}

3.10 Das 60/40-Portfolio

3.10.1 Aufbau und Ziel

Das traditionelle 60/40-Portfolio besteht aus 60\% Aktien und 40\%
festverzinslichen Anlagen wie Staatsanleihen. Der Aktienanteil soll
langfristige Renditechancen liefern. Der Anleihenanteil soll das
Portfolio stabilisieren und das Verlustrisiko gegenüber einem reinen
Aktienportfolio senken.\footnote{Pham, Cui, und Ruthbah, \emph{The
  {Performance} of the 60/40 {Portfolio}}, 4.}

\[
r_{60/40,t}=0{,}60r_{\mathrm{Aktien},t}+0{,}40r_{\mathrm{Anleihen},t}
\]

Dabei bezeichnet r60/40,t die Rendite des Gesamtportfolios zum Zeitpunkt
t. rAktien,t und rAnleihen,t sind die Renditen der beiden Anlageklassen;
0,60 und 0,40 sind ihre Zielgewichte.

Die Formel ist der 60/40-Spezialfall der gewichteten
Portfoliorendite.\footnote{Pham, Cui, und Ruthbah, 50.} Durch
unterschiedliche Renditen verschieben sich die Anteile im Zeitablauf.
Das Portfolio muss deshalb regelmässig auf die Zielgewichte von 60\% und
40\% zurückgesetzt werden.\footnote{Kommer, \emph{Souverän investieren
  mit Indexfonds und ETFs}, 7. Auflage (Weinheim: Campus Verlag, 2025),
  523--26.}

3.10.2 Diversifikation und Grenzen

Der Diversifikationseffekt entsteht, weil Aktien und Anleihen nicht
perfekt gleichgerichtet verlaufen. Die Korrelation zwischen beiden
Anlageklassen ist jedoch weder konstant noch immer negativ. Sie
unterschied sich in der langfristigen Untersuchung zwischen Ländern und
veränderte sich über die Zeit.\footnote{Pham, Cui, und Ruthbah,
  \emph{The {Performance} of the 60/40 {Portfolio}}, 1--7.}

Im Jahr 2022 verloren Aktien und Anleihen gleichzeitig an Wert. Der
Anleihenanteil schützt deshalb nicht in jeder einzelnen Marktphase. Über
lange Untersuchungszeiträume stellte das CFA-Paper trotzdem geringere
Volatilitäten und Drawdowns sowie positive reale Renditen des
60/40-Portfolios fest, wobei sich die Ergebnisse zwischen den Ländern
unterschieden.\footnote{Pham, Cui, und Ruthbah, 2--7.}

Auch Kommers historischer Vergleich zeigt den Zielkonflikt: Das
60/40-Portfolio erzielte eine tiefere Rendite als ein reines globales
Aktienportfolio, hatte dafür aber eine deutlich tiefere Volatilität und
einen geringeren Maximum Drawdown.\footnote{Kommer, \emph{Souverän
  investieren mit Indexfonds und ETFs}, 7. Auflage (Weinheim: Campus
  Verlag, 2025), 88.}

3.11 Marktkapitalisierungs- und BIP-Gewichtung

3.11.1 Marktkapitalisierungsgewichtung

Bei der Marktkapitalisierungsgewichtung erhält jedes Unternehmen ein
Gewicht entsprechend seinem Anteil am gesamten Börsenwert aller
enthaltenen Unternehmen. Dadurch entstehen indirekt auch Ländergewichte:
Länder mit einer grossen Marktkapitalisierung ihrer börsennotierten
Unternehmen erhalten ein grösseres Gewicht.\footnote{Kommer, 473--74.}

\[
w_i^{MK}=\frac{MK_i}{\sum_j MK_j}
\] Dabei bezeichnet \textbf{\(w_i^{\mathrm{MK}}\)} das
Marktkapitalisierungsgewicht des Elements \textbf{\(i\)},
\textbf{\(\mathrm{MK}_i\)} dessen Marktkapitalisierung und
\textbf{\(\sum_j \mathrm{MK}_j\)} die gesamte Marktkapitalisierung aller
berücksichtigten Elemente.

Die Gleichung ist eine eigene formale Darstellung der von Kommer
beschriebenen Gewichtung einzelner Aktien und Länder nach
Marktkapitalisierung.\footnote{Kommer, 474.}

3.11.2 BIP-Gewichtung

Bei der von Kommer beschriebenen BIP-Methode werden die Länder nach
ihrem Bruttoinlandsprodukt gewichtet. Innerhalb der einzelnen Länder
werden die Unternehmen weiterhin nach ihrer Marktkapitalisierung
gewichtet. Die Methode verändert somit hauptsächlich die regionale
Zusammensetzung des Portfolios.\footnote{Kommer, 474.}

\[
w_c^{BIP}=\frac{BIP_c}{\sum_d BIP_d}
\] Dabei bezeichnet \textbf{\(w_c^{\mathrm{BIP}}\)} das BIP-Gewicht des
Landes \textbf{\(c\)}, \textbf{\(\mathrm{BIP}_c\)} dessen
Bruttoinlandsprodukt und \textbf{\(\sum_d \mathrm{BIP}_d\)} das gesamte
Bruttoinlandsprodukt aller berücksichtigten Länder.

Die Gleichung ist eine eigene formale Darstellung der von Kommer
beschriebenen Ländergewichtung nach Bruttoinlandsprodukt.\footnote{Kommer,
  474--76.}

Kommers historische Gegenüberstellung zeigt über den Gesamtzeitraum nur
geringe Unterschiede. In einzelnen Teilperioden wechselte die besser
abschneidende Methode. Auch Volatilität und Maximum Drawdown lagen nahe
beieinander.\footnote{Kommer, 474--76.}

Kommer bevorzugt die BIP-Gewichtung, weil sie hohe Länderanteile und
mögliche nationale Klumpenrisiken reduzieren kann. Er hält jedoch selbst
fest, dass dieses Risiko mit den vorhandenen historischen Daten kaum
beweisbar ist. Die BIP-Gewichtung kann zudem komplexer sein und etwas
höhere Transaktionskosten verursachen.\footnote{Kommer, 476--77.} Diese
Position ist deshalb keine allgemein bewiesene Überlegenheit der
BIP-Methode.

3.12 Trendfolge mit gleitenden Durchschnitten

3.12.1 Einfacher gleitender Durchschnitt

Eine Trendfolgestrategie reagiert auf bereits beobachtete
Kursbewegungen. Hilpisch zeigt dazu eine algorithmische Handelsregel mit
zwei einfachen gleitenden Durchschnitten. Ein Simple Moving Average,
kurz SMA, ist der arithmetische Mittelwert der letzten n
Kursbeobachtungen.\footnote{Hilpisch, \emph{Python for Finance}, 326--28
  und 689-692.}

\[
SMA_t^{(n)}=\frac{1}{n}\sum_{i=0}^{n-1}P_{t-i}
\]

Dabei bezeichnet SMAt(n) den einfachen gleitenden Durchschnitt zum
Zeitpunkt t, n die Anzahl berücksichtigter Beobachtungen und Pt-i den
Kurs, der i Perioden vor t beobachtet wurde.

Quelle der Formel:.\footnote{Hilpisch, 326--28 und 689-692.}

Ein kurzer SMA (n=10-20) verwendet weniger Beobachtungen und reagiert
deshalb schneller auf neue Kursbewegungen. Ein langer SMA (n=100-200)
verändert sich langsamer, während sich ein mittelfristiger SMA
(n=20-100) dazwischen bewegt. Wenn sich die beiden Linien kreuzen,
spricht man von einem Crossover.\footnote{Hilpisch, 326--28.}

3.12.2 Handelssignal

Hilpisch verwendet eine Long/Short-Regel: Liegt der kurze SMA über dem
langen SMA, wird eine Long-Position gehalten. Liegt er darunter, wird
eine Short-Position gehalten.\footnote{Hilpisch, 691--92.} Für diese
Untersuchung wird die Regel methodisch zu Long/Cash angepasst. Bei einem
positiven Signal ist das Modell im Markt investiert, ansonsten hält es
Cash. Diese Anpassung ist eine eigene Entscheidung und nicht Hilpischs
ursprüngliche Regel.

\[
s_t=\begin{cases}1,&SMA_t^{kurz}>SMA_t^{lang}\\0,&SMA_t^{kurz}\leq SMA_t^{lang}\end{cases}
\]

Dabei bezeichnet st das Handelssignal zum Zeitpunkt t: Der Wert 1
bedeutet, dass die Strategie investiert ist, der Wert 0 steht für Cash.
SMAtkurz und SMAtlang sind der kurze und der lange gleitende
Durchschnitt.

Grundidee nach;\footnote{Hilpisch, 691--92.} die Long/Cash-Anpassung ist
eine eigene methodische Entscheidung.

Ein am Ende der Periode t berechnetes Signal darf nicht auf die bereits
realisierte Rendite derselben Periode angewendet werden. Die Position
wird deshalb um eine Periode verschoben.\footnote{Hilpisch, 694.}

\[
r_{\mathrm{Strategie},t}=s_{t-1}r_{\mathrm{Markt},t}
\]

Dabei bezeichnet rStrategie,t die Strategierendite zum Zeitpunkt t, st-1
das Signal der vorherigen Periode und rMarkt,t die aktuelle
Marktrendite.

Quelle der zeitlich verschobenen Anwendung des Signals:.\footnote{Hilpisch,
  694.}

\begin{Shaded}
\begin{Highlighting}[]
\CommentTok{\# Importiert die Numpy{-}Bibliothek}
\ImportTok{import}\NormalTok{ numpy }\ImportTok{as}\NormalTok{ np}
\CommentTok{\# Importiert die Pandas{-}Bibliothek}
\ImportTok{import}\NormalTok{ pandas }\ImportTok{as}\NormalTok{ pd}

\CommentTok{\# Die Funktion trendfolge() berechnet eine vereinfachte Long/Cash{-}Strategie}
\KeywordTok{def}\NormalTok{ trendfolge(preise: pd.Series, kurzes\_fenster: }\BuiltInTok{int}\NormalTok{, langes\_fenster: }\BuiltInTok{int}\NormalTok{):}
    \ControlFlowTok{if}\NormalTok{ kurzes\_fenster }\OperatorTok{\textgreater{}=}\NormalTok{ langes\_fenster:}
        \ControlFlowTok{raise} \PreprocessorTok{ValueError}\NormalTok{(}\StringTok{"Das kurze Fenster muss kleiner als das lange Fenster sein."}\NormalTok{)}
    
\NormalTok{    daten }\OperatorTok{=}\NormalTok{ pd.DataFrame(\{}\StringTok{"Kurs"}\NormalTok{: preise\}).dropna().sort\_index()}
\NormalTok{    daten[}\StringTok{"SMA kurz"}\NormalTok{] }\OperatorTok{=}\NormalTok{ daten[}\StringTok{"Kurs"}\NormalTok{].rolling(kurzes\_fenster).mean()}
\NormalTok{    daten[}\StringTok{"SMA lang"}\NormalTok{] }\OperatorTok{=}\NormalTok{ daten[}\StringTok{"Kurs"}\NormalTok{].rolling(langes\_fenster).mean()}
    
    \CommentTok{\# 1 bedeutet investiert und 0 bedeutet Cash}
\NormalTok{    daten[}\StringTok{"Signal"}\NormalTok{] }\OperatorTok{=}\NormalTok{ np.where(daten[}\StringTok{"SMA kurz"}\NormalTok{] }\OperatorTok{\textgreater{}}\NormalTok{ daten[}\StringTok{"SMA lang"}\NormalTok{], }\DecValTok{1}\NormalTok{, }\DecValTok{0}\NormalTok{)}
\NormalTok{    daten[}\StringTok{"Marktrendite"}\NormalTok{] }\OperatorTok{=}\NormalTok{ daten[}\StringTok{"Kurs"}\NormalTok{].pct\_change()}
    
    \CommentTok{\# shift(1) verhindert die Nutzung eines noch nicht verfügbaren Signals}
\NormalTok{    daten[}\StringTok{"Strategierendite"}\NormalTok{] }\OperatorTok{=}\NormalTok{ daten[}\StringTok{"Signal"}\NormalTok{].shift(}\DecValTok{1}\NormalTok{) }\OperatorTok{*}\NormalTok{ daten[}\StringTok{"Marktrendite"}\NormalTok{]}
    \ControlFlowTok{return}\NormalTok{ daten}

\CommentTok{\# Beispiel zur Nutzung der Funktion}
\NormalTok{beispiel\_preise }\OperatorTok{=}\NormalTok{ pd.Series(}
\NormalTok{    [}\DecValTok{100}\NormalTok{, }\DecValTok{101}\NormalTok{, }\DecValTok{103}\NormalTok{, }\DecValTok{102}\NormalTok{, }\DecValTok{105}\NormalTok{, }\DecValTok{107}\NormalTok{, }\DecValTok{104}\NormalTok{, }\DecValTok{102}\NormalTok{, }\DecValTok{100}\NormalTok{, }\DecValTok{103}\NormalTok{, }\DecValTok{106}\NormalTok{],}
\NormalTok{    index}\OperatorTok{=}\NormalTok{pd.date\_range(}\StringTok{"2026{-}01{-}01"}\NormalTok{, periods}\OperatorTok{=}\DecValTok{11}\NormalTok{, freq}\OperatorTok{=}\StringTok{"D"}\NormalTok{)}
\NormalTok{)}

\BuiltInTok{print}\NormalTok{(trendfolge(beispiel\_preise, }\DecValTok{3}\NormalTok{, }\DecValTok{5}\NormalTok{))}
\end{Highlighting}
\end{Shaded}

\begin{verbatim}
            Kurs    SMA kurz  SMA lang  Signal  Marktrendite  Strategierendite
2026-01-01   100         NaN       NaN       0           NaN               NaN
2026-01-02   101         NaN       NaN       0      0.010000          0.000000
2026-01-03   103  101.333333       NaN       0      0.019802          0.000000
2026-01-04   102  102.000000       NaN       0     -0.009709         -0.000000
2026-01-05   105  103.333333     102.2       1      0.029412          0.000000
2026-01-06   107  104.666667     103.6       1      0.019048          0.019048
2026-01-07   104  105.333333     104.2       1     -0.028037         -0.028037
2026-01-08   102  104.333333     104.0       1     -0.019231         -0.019231
2026-01-09   100  102.000000     103.6       0     -0.019608         -0.019608
2026-01-10   103  101.666667     103.2       0      0.030000          0.000000
2026-01-11   106  103.000000     103.0       0      0.029126          0.000000
\end{verbatim}

3.12.3 Grenzen der Trendfolge

Gleitende Durchschnitte sind nachlaufende Indikatoren, weil sie nur
vergangene Kurse verwenden. Ein Trend wird somit erst erkannt, nachdem
die Kursbewegung bereits begonnen hat. In Seitwärtsphasen können häufige
Kreuzungen zu Fehlsignalen und wiederholten Positionswechseln
führen.\footnote{Hilpisch, 326--28 und 689-694.}

Hilpisch zeigt, dass viele Kombinationen von kurzen und langen Fenstern
nachträglich getestet werden können.\footnote{Hilpisch, 697.} Wird nur
die historisch beste Kombination ausgewählt und danach mit denselben
Daten beurteilt, besteht ein Overfitting-Risiko. Die Hauptparameter
müssen deshalb vor dem eigentlichen Vergleich begründet und festgelegt
werden.

Der vereinfachte Beispiel-Backtest von Hilpisch lässt unter anderem
Gebühren, Geld-Brief-Spannen (Differenz zwischen Bid und Ask) und
weitere Marktmechanismen weg.\footnote{Hilpisch, 696.} Solche
Vereinfachungen können eine Strategie mit häufigen Transaktionen im
historischen Vergleich begünstigen.

3.13 Aussagegrenzen historischer Vergleiche und Backtests

3.13.1 Historische Stichprobe und Vergleichbarkeit

Ein Backtest zeigt, wie eine genau definierte Strategie innerhalb einer
gewählten historischen Stichprobe abgeschnitten hätte. Das Ergebnis
garantiert nicht, dass dieselbe Beziehung auch in der Zukunft gilt. Das
CFA-Paper zeigt beispielsweise unterschiedliche Ergebnisse des
60/40-Portfolios je nach Land, Generation und Marktphase.\footnote{Pham,
  Cui, und Ruthbah, \emph{The {Performance} of the 60/40 {Portfolio}},
  1--7.} Auch beim Vergleich von Marktkapitalisierungs- und
BIP-Gewichtung wechselte die besser abschneidende Methode zwischen den
Teilperioden.\footnote{Kommer, \emph{Souverän investieren mit Indexfonds
  und ETFs}, 7. Auflage (Weinheim: Campus Verlag, 2025), 474--76.}

Für einen fairen Vergleich müssen alle Strategien auf derselben
Grundlage untersucht werden. Dazu gehören ein identischer Zeitraum, eine
einheitliche Datenfrequenz je Kennzahl, dieselbe Währung, dieselbe
Behandlung von Dividenden und eine einheitliche nominale oder reale
Berechnung.\footnote{Kommer, 22--23, 166 und 466-467; Hilpisch,
  \emph{Python for Finance}, 319--21.}

3.13.2 Foresight Bias

Foresight Bias oder Look-ahead Bias entsteht, wenn in einem Backtest
Informationen verwendet werden, die zum betrachteten Zeitpunkt noch
nicht verfügbar waren. Beim Resampling kann dies geschehen, wenn ein
Monatsendwert zeitlich am Monatsanfang eingeordnet wird. Hilpisch
fordert deshalb eine zeitlich korrekte Beschriftung und den letzten
verfügbaren Wert des Intervalls.\footnote{Hilpisch, \emph{Python for
  Finance}, 319--21.}

Bei der Trendfolge darf ein am Ende des Tages berechnetes Signal nicht
auf die Rendite desselben Tages angewendet werden. Die Verschiebung des
Signals um eine Periode verhindert, dass das Modell zukünftige
Informationen nutzt.\footnote{Hilpisch, 694.}

3.13.3 Data Snooping und Overfitting

Data Snooping entsteht, wenn historische Daten bewusst oder unbewusst so
lange durchsucht werden, bis ein gewünschtes Ergebnis erscheint. Kommer
verbindet den Begriff unter anderem mit Backtesting with Overfitting und
verweist auf Out-of-Sample-Tests als Gegenmassnahme.\footnote{Kommer,
  \emph{Souverän investieren mit Indexfonds und ETFs}, 7. Auflage
  (Weinheim: Campus Verlag, 2025), 607--8.}

Ein Beispiel ist die nachträgliche Suche nach der besten Kombination
zweier SMA-Fenster. Wenn dieselben Daten sowohl zur Auswahl der
Parameter als auch zur Bewertung der Strategie genutzt werden, kann das
Modell zufällige Eigenschaften der Vergangenheit statt einer stabilen
Regel abbilden.\footnote{Hilpisch, \emph{Python for Finance}, 697.}

3.13.4 Kosten, Inflation und Datenqualität

Ein vereinfachter Backtest kann Transaktionsgebühren, Geld-Brief-Spannen
und weitere Marktmechanismen auslassen. Hilpisch nennt diese Punkte
ausdrücklich als Vereinfachungen seines Beispiels.\footnote{Hilpisch,
  696.} Werden sie nicht berücksichtigt, können häufig handelnde
Strategien im Vergleich zu Buy-and-Hold zu gut erscheinen.

Eine Analyse ohne Inflationsbereinigung zeigt nominales
Vermögenswachstum und nicht die Veränderung der Kaufkraft. Für
langfristige Vergleiche können reale Renditen deshalb eine andere
Aussage liefern als nominale Renditen.\footnote{Kommer, \emph{Souverän
  investieren mit Indexfonds und ETFs}, 7. Auflage (Weinheim: Campus
  Verlag, 2025), 22--23.}

Auch fehlende Werte, unterschiedliche Handelstage, die zeitliche
Ausrichtung mehrerer Reihen und die Verwendung von Kurs- statt
Total-Return-Daten können das Ergebnis beeinflussen.\footnote{Hilpisch,
  \emph{Python for Finance}, 307--21; Kommer, \emph{Souverän investieren
  mit Indexfonds und ETFs}, 7. Auflage (Weinheim: Campus Verlag, 2025),
  466--67.} Diese Annahmen müssen in der Methodik für alle Strategien
eindeutig und gleich festgelegt werden.

\bookmarksetup{startatroot}

\chapter{Methodik}\label{methodik}

\section{Investment Strategien}\label{investment-strategien}

\subsection{Marktbreites investieren}\label{marktbreites-investieren}

Mit ETFs und Index Fonds kann man in die Gesamtwirtschaft investieren
und kann über einen langen Zeitraum meistens Rendite erzielen.

\subsection{60/40 balanced Portfolio}\label{balanced-portfolio}

Es wird in 60\% Aktien für Rendite und 40\% Staatsanleihen für die
Sicherheit investiert. Pham, Cui, und Ruthbah\footnote{\emph{The
  {Performance} of the 60/40 {Portfolio}}.}

\subsection{Trendfolge-Strategie}\label{trendfolge-strategie}

Hinter der Trendfolgestrategie steckt die Annahme, dass es einfacher
ist, bereits bestehende Trends auszunutzen als frühzeitig eine
Trendumkehr zu erkennen.~Man partizipiert dann zwar nicht von
Tiefststand zu Höchststand (oder umgekehrt, wenn man Short geht) von
einem Trend, läuft aber auch nicht Gefahr, die völlig falsche
Entscheidung getroffen haben, wenn das Gegenteil von dem eintritt, was
man angenommen hat. Admirals\footnote{{„Wie Sie mit der
  Trendfolgestrategie bestehende Trends nutzen``}.} Um den Trend zu
erkennen, werden technische Indikatoren wie der gleitende Durchschnitt
(Moving Average) Probst\footnote{{„Der Gleitende Durchschnitt - so
  bestimmen Sie die Trendrichtung``}.} oder der Relative Strength Index
(RSI) Ratgeber\footnote{{„Relative Strength Index - so funktioniert der
  RSI! - finanzen.ch``}.} verwendet.

\subsection{BIP-Gewichtet investieren (Gerd
Kommer)}\label{bip-gewichtet-investieren-gerd-kommer}

Der Gerd-Kommer-ETF ist ein All-Cap-- und All-World-ETF, also ein ETF
auf die „Welt-AG``, der bei der Ländergewichtung sowohl
Markt\-kapitalisierung als auch Wirtschafts\-leistung („BIP``)
berücksichtigt, dabei die wichtigsten fünf Faktor\-prämien („Factor
Investing``) übergewichtet und mit über 4.000 Einzelaktien
Ultradiversifikation bietet. Gerd Kommer\footnote{{„Investieren``}.}

\section{Datenquellen}\label{datenquellen}

\subsection{Yahoo Finance}\label{yahoo-finance-1}

Yahoo Finance bietet historische Kursdaten für eine Vielzahl von
Finanzinstrumenten, darunter Aktien, ETFs und Anleihen. Diese Daten
können für die Berechnung von Renditen, Volatilität und anderen
Kennzahlen verwendet werden. Yahoo\footnote{{„Yahoo FinanzenDeutschland
  \textbar{} Börse live, Kurse, Nachrichten aus Wirtschafts- und
  Finanzwelt``}.}

\subsection{FRED (Federal Reserve Economic
Data)}\label{fred-federal-reserve-economic-data-1}

FRED ist eine umfangreiche Datenbank der Federal Reserve Bank of
St.~Louis, die wirtschaftliche und finanzielle Daten bereitstellt. Hier
können makroökonomische Indikatoren wie das Bruttoinlandsprodukt (BIP),
Arbeitslosenzahlen und Zinssätze abgerufen werden, die für die Analyse
von Anlagestrategien relevant sein können. Federal Reserve Bank of St.
Louis\footnote{{„Federal {Reserve Economic Data} \textbar{} {FRED}
  \textbar{} {St}. {Louis Fed}``}.}

\subsection{Vanguard}\label{vanguard-1}

Vanguard ist einer der grössten Anbieter von ETFs und bietet
umfangreiche Informationen zu seinen Produkten, einschliesslich
historischer Renditen, Kostenquoten und Fondsinformationen. Diese Daten
können für die Analyse von ETF-basierten Anlagestrategien verwendet
werden. Vanguard\footnote{{„Startseite \textbar{} Vanguard Switzerland
  Private Anleger``}.}

\section{Kennzahlen}\label{kennzahlen}

\subsection{Stabilität}\label{stabilituxe4t}

\textbf{Volatilität}: Die Volatilität misst die Schwankungsbreite der
Renditen. Eine hohe Volatilität deutet auf starke Kursschwankungen hin,
während eine niedrige Volatilität auf stabilere Renditen hindeutet.
Buchhaltung-Einfach-Sicher\footnote{{„Volatilität Einfach Erklärt
  \textbar{} {Definition}, {Berechnung} Und {Nutzen}``}.}

\textbf{Verlustphasen}: Verlustphasen sind Zeiträume, in denen die
Renditen negativ sind. Sie können durch verschiedene Faktoren verursacht
werden, wie z.B. wirtschaftliche Abschwünge, geopolitische Ereignisse
oder Unternehmenskrisen. Das Verständnis von Verlustphasen ist wichtig,
um die Risiken einer Anlagestrategie zu bewerten und um angemessene
Massnahmen zur Risikominderung zu ergreifen. Gegenspieler von positiven
Phasen.

\textbf{Positive Phasen}: Positive Phasen sind Zeiträume, in denen die
Renditen positiv sind. Sie können durch wirtschaftliches Wachstum,
technologische Innovationen oder andere günstige Faktoren verursacht
werden. Das Verständnis von positiven Phasen ist wichtig, um die Chancen
einer Anlagestrategie zu bewerten und um angemessene Massnahmen zur
Maximierung der Rendite zu ergreifen. Gegenspieler von Verlustphasen.

\textbf{Maximaler Drawdown}: Der maximale Drawdown misst den grössten
prozentualen Verlust vom Höchststand bis zum Tiefststand eines
Portfolios. Er gibt an, wie viel ein Anleger in einem ungünstigen
Szenario verlieren könnte. Fidelity\footnote{{„Maximum {Drawdown}``}.}

\textbf{Krisenfestigkeit}: Die Krisenfestigkeit eines Portfolios
beschreibt, wie gut es in Zeiten wirtschaftlicher oder finanzieller
Krisen performt. Ein krisenfestes Portfolio sollte in der Lage sein,
Verluste zu minimieren und sich schneller von Rückschlägen zu erholen.
Beispiel Krisen: Dotcom-Blase, Finanzkrise 2008, Corona-Krise

\subsection{Rendite}\label{rendite}

\textbf{Gesamtrendite}: Die Gesamtrendite berücksichtigt sowohl
Kursgewinne als auch Dividendenzahlungen. Sie gibt an, wie viel Prozent
das investierte Kapital über einen bestimmten Zeitraum gewachsen ist.

\textbf{Durchschnittliche jährliche Rendite}: Diese Kennzahl berechnet
die durchschnittliche jährliche Rendite über einen bestimmten Zeitraum,
indem sie die Gesamtrendite auf eine jährliche Basis umrechnet.

\subsection{Sharpe Ratio}\label{sharpe-ratio}

Die Sharpe Ratio misst die risikobereinigte Rendite einer
Anlagestrategie. Sie berechnet sich als Differenz zwischen der Rendite
der Strategie und der risikofreien Rendite, dividiert durch die
Volatilität der Strategie. Eine höhere Sharpe Ratio deutet auf eine
bessere risikobereinigte Performance hin. Fidelity\footnote{{„Sharpe-{Ratio}``}.}

\section{Bedingungen}\label{bedingungen}

\begin{itemize}
\tightlist
\item
  Zeitabschnitt
\item
  Gleiches Startkapital
\item
  Reinvestition der Dividenden
\item
  Gleiche Datenfrequenz (täglich, wöchentlich, monatlich, jährlich)
\item
  jährliches Rebalancing (für 60/40 Portfolio)
\end{itemize}

Nicht berücksichtigt werden: - Transaktionskosten - Steuern - Inflation
- weitere Gebühren

\bookmarksetup{startatroot}

\chapter{Analyse}\label{analyse}

5.1 Ausgangspunkt der Analyse

Die Theorie hat bereits die Bausteine festgelegt, die das spätere Modell
benötigt. Dazu gehören die Verarbeitung von Zeitreihen,
Renditeberechnungen, Volatilität, Drawdown, Sharpe Ratio, Korrelation,
Portfoliorisiko, Rebalancing, Buy-and-Hold und eine Trendfolgestrategie.
Die Methodik verlangt ausserdem einen vergleichbaren Zeitabschnitt,
dasselbe Startkapital, die Reinvestition der Dividenden und eine
einheitliche Datenfrequenz. Beim 60/40-Portfolio ist ein jährliches
Rebalancing vorgesehen.

Damit besteht meine Aufgabe in der Analyse nicht darin, eine völlig neue
Finanztheorie zu entwickeln. Ich muss die bereits beschriebenen
Bestandteile so zusammenführen, dass sie mit realen historischen Daten
reproduzierbar ausgeführt werden können.

Dabei unterscheide ich zwei Ebenen: 1. Das allgemeine System: Es soll
verschiedene geeignete Datensätze und Parameter verarbeiten können. 2.
Die Untersuchung der Maturarbeit: Sie verwendet später eine feste und
begründete Auswahl von Daten und Parametern.

Diese Trennung ist wichtig, weil die spätere Benutzeroberfläche dieselbe
Engine mit anderen Eingaben verwenden können soll. Für die
wissenschaftliche Auswertung werden die Parameter dagegen vor dem
Hauptlauf festgelegt und nicht nachträglich so verändert, dass ein
günstigeres Resultat entsteht. Das entspricht auch der im Theorieteil
behandelten Gefahr von Data Snooping und Overfitting.\footnote{Hilpisch,
  \emph{Python for Finance}, 697; Kommer, \emph{Souverän investieren mit
  Indexfonds und ETFs}, 7. Auflage (Weinheim: Campus Verlag, 2025),
  607--8.}

5.2 Reihenfolge der technischen Entwicklung

Aus den Anforderungen ergibt sich für mich folgende Reihenfolge:

\[
\begin{array}{c}
\boxed{\text{\Large 1. Internes Datenformat festlegen}}
\\[0.4cm]
\downarrow
\\[0.4cm]
\boxed{\text{\Large 2. Datenquellen auf dieses Format abbilden}}
\\[0.4cm]
\downarrow
\\[0.4cm]
\boxed{\text{\Large 3. Import, Normalisierung und Validierung definieren}}
\\[0.4cm]
\downarrow
\\[0.4cm]
\boxed{\text{\Large 4. Veränderbare Einstellungen des Systems festlegen}}
\\[0.4cm]
\downarrow
\\[0.4cm]
\boxed{\text{\Large 5. Einheitlichen Ergebnisvertrag festlegen}}
\\[0.4cm]
\downarrow
\\[0.4cm]
\boxed{\text{\Large 6. Leeres Grundgerüst der Engine definieren}}
\\[0.4cm]
\downarrow
\\[0.4cm]
\boxed{\text{\Large 7. Vorhandene Funktionen integrieren}}
\end{array}
\]

Die Reihenfolge ist bewusst so gewählt. Würde ich zuerst die
Simulationsfunktion programmieren und erst danach entscheiden, welche
Eingaben und Ausgaben sie besitzen soll, würden Datenformat, Strategien
und Benutzeroberfläche unnötig voneinander abhängig.

Die konkrete Reihenfolge ist meine eigene Designentscheidung. Die
Quellen liefern die Grundlagen dafür, welche Datenstrukturen und
Finanzinformationen benötigt werden; sie schreiben nicht diese exakte
Softwarearchitektur vor.

5.3 Entwicklung eines einheitlichen Datenformats

5.3.1 Warum überhaupt ein eigenes internes Format?

Die später verwendeten Daten können aus unterschiedlichen Quellen
stammen. Yahoo Finance liefert Marktzeitreihen anders als FRED eine
Zinsreihe oder eine Makrodatenbank jährliche BIP-Werte. Würde jede
Strategie direkt mit dem Format einer bestimmten Quelle arbeiten, müsste
sich auch die Strategie ändern, sobald der Anbieter gewechselt wird.

Deshalb erhält das Modell einen anbieterunabhängigen Datenvertrag. Ein
Importmodul ist dafür verantwortlich, Quelldaten in diesen Vertrag zu
übersetzen. Die eigentliche Strategie sieht danach nur noch
standardisierte Daten.

Wes McKinney beschreibt Datenanalyse als Abfolge von Einlesen,
Bereinigen, Kombinieren, Normalisieren, Umformen und anschliessender
Berechnung. Er behandelt dabei sowohl tabellarische Daten als auch
Zeitreihen und zeigt, dass pandas gerade für solche strukturierten Daten
entwickelt wurde.\footnote{McKinney, \emph{Python for Data Analysis}.}
Diese Grundidee passt zu meiner Untersuchung: Die Quelle darf
unterschiedlich sein, die Berechnung soll danach aber auf einer
einheitlichen Struktur stattfinden.

Die Entscheidung, dafür einen eigenen Datenvertrag zu definieren, ist
meine technische Schlussfolgerung aus diesem Bedarf.

5.3.2 Warum CSV als Austauschformat?

Für den Austausch zwischen Import, Kontrolle und späterer
Benutzeroberfläche verwende ich in der ersten Version CSV-Dateien.
McKinney behandelt in Kapitel 6 das Einlesen tabellarischer Textdaten
und verwendet pandas.read\_csv() als eines der zentralen Werkzeuge für
solche Daten.\footnote{McKinney.}

CSV ist für mein Projekt nicht deshalb ausgewählt, weil es allen anderen
Formaten technisch überlegen wäre. Die Gründe sind praktischer:

\begin{itemize}
\tightlist
\item
  Die Datei kann direkt mit einem Texteditor, Tabellenprogramm oder
  Python geöffnet werden.
\item
  Beim Entwickeln kann ich leicht kontrollieren, ob ein Adapter
  tatsächlich die erwarteten Spalten erzeugt.
\item
  Eine spätere Weboberfläche kann eine CSV-Datei einfach als
  Benutzereingabe akzeptieren.
\item
  pandas kann das Format ohne zusätzliche Spezialbibliothek einlesen.
\end{itemize}

CSV besitzt gleichzeitig Grenzen. Datentypen und die Bedeutung einer
Spalte werden nicht ausreichend durch die Datei selbst erklärt. Darum
genügt eine CSV allein nicht: Das System benötigt zusätzlich einen klar
definierten Datenvertrag und Metadaten.

Für die interne Berechnung bleibt die Datei nicht als Textstruktur
bestehen. Nach dem Einlesen wird sie in pandas-Objekte überführt.

5.3.3 Warum ein Long-Format?

Für die Marktzeitreihen verwende ich als Austauschformat eine lange
Tabelle. Eine Beobachtung besteht aus einem Datum, einer Anlagenkennung
und einem Wert.

McKinney beschreibt im Abschnitt ``Pivoting `Long' to `Wide' Format''
ausdrücklich, dass mehrere Zeitreihen häufig in einem sogenannten long
beziehungsweise stacked format gespeichert werden. In diesem Format
stellt jede einzelne Beobachtung eine eigene Zeile dar. Er zeigt danach
mit DataFrame.pivot(), wie dieselben Daten für Berechnungen in eine
breite Tabelle mit einer Spalte pro Zeitreihe umgeformt werden
können.\footnote{McKinney.}

Diese Beschreibung passt direkt zu meinem Problem. Für ein allgemeines
Eingabeformat ist Long praktisch, weil eine neue Anlage zusätzliche
Zeilen erzeugt, ohne dass sich die Grundstruktur der Tabelle ändert. Für
bestimmte Berechnungen ist dagegen eine breite Darstellung komfortabler.
Deshalb lege ich nicht fest, dass jede Stufe des Systems Long sein muss:

\[
\begin{array}{c}
\text{\large \textbf{Austausch / Speicherung:}}
\\[4pt]
\boxed{\texttt{date \;|\; asset\_id \;|\; performance\_value}}
\\[10pt]
\downarrow
\\[-2pt]
\text{\texttt{pandas pivot()}}
\\[10pt]
\downarrow
\\[-2pt]
\text{\large \textbf{Interne Berechnung:}}
\\[4pt]
\boxed{\texttt{date \;|\; ASSET\_A \;|\; ASSET\_B \;|\; ASSET\_C \;|\; \ldots}}
\end{array}
\]

McKinney liefert die nachvollziehbare Grundlage dafür, dass Long ein
gebräuchliches Format für mehrere Zeitreihen ist und sich mit pandas
eindeutig in Wide umformen lässt. Meine Entscheidung ist, Long deshalb
als äussere Schnittstelle und Wide dort zu verwenden, wo die Berechnung
davon profitiert.

5.3.4 Marktzeitreihen: minimaler Datenvertrag

Für eine Marktzeitreihe definiere ich zunächst drei Pflichtfelder:

\[
\begin{array}{|l|l|}
\hline
\textbf{Feld} & \textbf{Bedeutung} \\
\hline
\texttt{date} & \text{Datum der Beobachtung} \\
\hline
\texttt{asset\_id} & \text{interne, anbieterunabhängige Kennung der Anlage} \\
\hline
\texttt{performance\_value} & \text{Wertreihe, aus der die Performance berechnet werden darf} \\
\hline
\end{array}
\]

Optional kann eine vierte Spalte vorhanden sein:

\[
\begin{array}{|l|l|}
\hline
\textbf{Feld} & \textbf{Bedeutung} \\
\hline
\texttt{signal\_value} & \text{separate Wertreihe für technische Handelssignale} \\
\hline
\end{array}
\]

Beispiel:

\[
\begin{array}{|c|c|c|c|}
\hline
\textbf{date} & \textbf{asset\_id} & \textbf{performance\_value} & \textbf{signal\_value} \\
\hline
2020\text{-}01\text{-}02 & \texttt{EQ\_DEMO} & 100.0 & 50.25 \\
\hline
2020\text{-}01\text{-}03 & \texttt{EQ\_DEMO} & 102.0 & 51.00 \\
\hline
2020\text{-}01\text{-}02 & \texttt{BD\_DEMO} & 100.0 & 104.80 \\
\hline
2020\text{-}01\text{-}03 & \texttt{BD\_DEMO} & 100.5 & 104.90 \\
\hline
\end{array}
\]

asset\_id ist bewusst nicht dasselbe wie ein Yahoo-Ticker. Der Ticker
gehört zur Quelle; asset\_id gehört zum Modell. Dadurch könnte dieselbe
logische Anlage später aus einer anderen Quelle importiert werden, ohne
dass alle Strategiekonfigurationen angepasst werden müssen.

Auch performance\_value ist absichtlich kein Anbieterbegriff wie Adj
Close. Der Importadapter muss dokumentieren, welche Quellenspalte auf
dieses Feld abgebildet wurde.

\begin{Shaded}
\begin{Highlighting}[]
\ImportTok{from}\NormalTok{ io }\ImportTok{import}\NormalTok{ StringIO}
\ImportTok{import}\NormalTok{ pandas }\ImportTok{as}\NormalTok{ pd}

\NormalTok{beispiel\_csv }\OperatorTok{=} \StringTok{"}\CharTok{\textbackslash{}n}\StringTok{"}\NormalTok{.join([}
    \StringTok{"date,asset\_id,performance\_value"}\NormalTok{,}
    \StringTok{"2020{-}01{-}02,EQ\_DEMO,100.0"}\NormalTok{,}
    \StringTok{"2020{-}01{-}03,EQ\_DEMO,102.0"}\NormalTok{,}
    \StringTok{"2020{-}01{-}06,EQ\_DEMO,101.0"}\NormalTok{,}
    \StringTok{"2020{-}01{-}02,BD\_DEMO,100.0"}\NormalTok{,}
    \StringTok{"2020{-}01{-}03,BD\_DEMO,100.5"}\NormalTok{,}
    \StringTok{"2020{-}01{-}06,BD\_DEMO,101.0"}\NormalTok{,}
\NormalTok{])}

\NormalTok{markt\_long }\OperatorTok{=}\NormalTok{ pd.read\_csv(StringIO(beispiel\_csv))}
\NormalTok{markt\_long[}\StringTok{"date"}\NormalTok{] }\OperatorTok{=}\NormalTok{ pd.to\_datetime(markt\_long[}\StringTok{"date"}\NormalTok{])}
\NormalTok{markt\_long}
\end{Highlighting}
\end{Shaded}

{\def\LTcaptype{none} % do not increment counter
\begin{longtable}[]{@{}llll@{}}
\toprule\noalign{}
& date & asset\_id & performance\_value \\
\midrule\noalign{}
\endhead
\bottomrule\noalign{}
\endlastfoot
0 & 2020-01-02 & EQ\_DEMO & 100.0 \\
1 & 2020-01-03 & EQ\_DEMO & 102.0 \\
2 & 2020-01-06 & EQ\_DEMO & 101.0 \\
3 & 2020-01-02 & BD\_DEMO & 100.0 \\
4 & 2020-01-03 & BD\_DEMO & 100.5 \\
5 & 2020-01-06 & BD\_DEMO & 101.0 \\
\end{longtable}
}

Für die spätere Portfolioberechnung kann dieselbe Tabelle eindeutig
umgeformt werden. Dadurch muss die externe Datei nicht bereits für jede
konkrete mathematische Funktion optimiert sein.

\begin{Shaded}
\begin{Highlighting}[]
\NormalTok{markt\_wide }\OperatorTok{=}\NormalTok{ (}
\NormalTok{    markt\_long}
\NormalTok{    .pivot(index}\OperatorTok{=}\StringTok{"date"}\NormalTok{, columns}\OperatorTok{=}\StringTok{"asset\_id"}\NormalTok{, values}\OperatorTok{=}\StringTok{"performance\_value"}\NormalTok{)}
\NormalTok{    .sort\_index()}
\NormalTok{    .reset\_index()}
\NormalTok{)}
    
\NormalTok{markt\_wide.columns.name }\OperatorTok{=} \VariableTok{None}
\NormalTok{markt\_wide}
\end{Highlighting}
\end{Shaded}

{\def\LTcaptype{none} % do not increment counter
\begin{longtable}[]{@{}llll@{}}
\toprule\noalign{}
& date & BD\_DEMO & EQ\_DEMO \\
\midrule\noalign{}
\endhead
\bottomrule\noalign{}
\endlastfoot
0 & 2020-01-02 & 100.0 & 100.0 \\
1 & 2020-01-03 & 100.5 & 102.0 \\
2 & 2020-01-06 & 101.0 & 101.0 \\
\end{longtable}
}

5.3.5 Was bedeutet performance\_value?

Die Spalte performance\_value ist die Wertreihe, aus der das Modell
später die tatsächliche Wertentwicklung einer Anlage berechnet. Dabei
soll nicht einfach irgendein verfügbarer Kurs übernommen werden. Für
meine Untersuchung ist entscheidend, dass die gewählte Reihe zur
festgelegten Behandlung der Dividenden passt.

Im Theorieteil wurde bereits festgelegt, dass für den langfristigen
Vergleich eine Gesamtrendite beziehungsweise ein Total Return geeigneter
ist als eine reine Kursbetrachtung. Ein Total Return berücksichtigt
neben Kursveränderungen auch laufende Erträge wie Dividenden.\footnote{Kommer,
  \emph{Souverän investieren mit Indexfonds und ETFs}, 7. Auflage
  (Weinheim: Campus Verlag, 2025), 466--67.} Da in meiner Methodik die
Reinvestition der Dividenden vorgesehen ist, muss die verwendete
Performance-Reihe diese Annahme abbilden.

performance\_value bezeichnet deshalb bewusst keinen bestimmten
Anbieterbegriff. Je nach Datenquelle kann dahinter zum Beispiel ein
Total-Return-Index oder eine entsprechend bereinigte Wertreihe stehen.
Der jeweilige Import muss dokumentieren, welche ursprüngliche Reihe
verwendet wurde.

Dadurch bleibt die Berechnung unabhängig von der Quelle. Die Engine
erhält immer dieselbe Spalte, während die Umwandlung davor sicherstellt,
dass ihre Bedeutung vergleichbar ist.

5.3.6 Wozu dient signal\_value?

Für Buy-and-Hold oder ein Rebalancing-Portfolio genügt grundsätzlich die
Performance-Reihe. Bei einer technischen Strategie wie der Trendfolge
kommt jedoch eine zweite Aufgabe hinzu: Aus einer Kursreihe wird ein
Handelssignal berechnet.

Hilpisch verwendet für seine SMA-Strategie zwei gleitende Durchschnitte
und leitet daraus eine Long- beziehungsweise Short-Position
ab.\footnote{Hilpisch, \emph{Python for Finance}, 691--92.} Für meine
Untersuchung wurde diese Regel bereits im Theorieteil zu Long/Cash
angepasst. Ausserdem wird das Signal um eine Periode verschoben, damit
keine bereits bekannte Rendite derselben Periode verwendet
wird.\footnote{Hilpisch, 694.}

Damit das allgemeine System später nicht voraussetzen muss, dass
Performance-Berechnung und Signal immer auf exakt derselben Reihe
beruhen, kann zusätzlich signal\_value angegeben werden. Fehlt diese
Spalte, darf die Strategie performance\_value nur dann verwenden, wenn
dies in der Konfiguration ausdrücklich erlaubt ist.

Für die erste Untersuchung kann sich später herausstellen, dass beide
Felder auf dieselbe Reihe zeigen. Die Trennung bleibt trotzdem sinnvoll,
weil sie die Datenstruktur nicht auf diese eine Untersuchung beschränkt.

5.3.7 Metadaten und zusätzliche Datentypen

Die Marktdatei enthält absichtlich nur die Werte, die sich über die Zeit
verändern. Weitere Informationen zu einer Anlage werden separat
gespeichert. Dazu gehören beispielsweise Name, Anlageklasse, Währung,
Land und ursprüngliche Datenquelle.

Für diese Informationen ist eine Datei assets.csv vorgesehen:

{[}

\begin{array}{c|c}
\textbf{Feld} & \textbf{Bedeutung} \\
\hline
\texttt{asset\_id} & \text{Interne, anbieterunabhängige Kennung der Anlage} \\
\texttt{name} & \text{Name der Anlage} \\
\texttt{asset\_class} & \text{Anlageklasse, z.\,B. Aktien oder Anleihen} \\
\texttt{country} & \text{Zugeordnetes Land bzw. Markt} \\
\texttt{currency} & \text{Währung der Anlage} \\
\texttt{provider} & \text{Datenanbieter} \\
\texttt{provider\_symbol} & \text{Symbol bzw. Kennung beim Datenanbieter}
\end{array}

{]}

Die interne Kennung asset\_id verbindet die Metadaten mit den
Marktwerten. Dadurch kann zum Beispiel US\_EQUITY im Modell gleich
bleiben, auch wenn die Daten später von einem anderen Anbieter bezogen
werden.

Nicht alle benötigten Daten sind Marktpreise. Für die Sharpe Ratio wird
eine risikofreie Rendite benötigt. Sie muss in derselben Periodenlogik
wie die untersuchte Rendite verwendet werden.\footnote{Kommer,
  \emph{Souverän investieren mit Indexfonds und ETFs}, 7. Auflage
  (Weinheim: Campus Verlag, 2025), 165--66.} Dafür ist eine getrennte
Datei vorgesehen:

{[}

\begin{array}{c|c}
\textbf{Feld} & \textbf{Bedeutung} \\
\hline
\texttt{date} & \text{Datum der Beobachtung} \\
\texttt{series\_id} & \text{Kennung der risikofreien Zinsreihe} \\
\texttt{value} & \text{Beobachteter Zinssatz} \\
\texttt{unit} & \text{Einheit des Wertes}
\end{array}

{]}

Auch Makrodaten wie das BIP unterscheiden sich von täglichen
Marktzeitreihen. Sie werden deshalb separat gespeichert:

{[}

\begin{array}{c|c}
\textbf{Feld} & \textbf{Bedeutung} \\
\hline
\texttt{period} & \text{Zeitraum der Beobachtung} \\
\texttt{country} & \text{Land, auf das sich der Wert bezieht} \\
\texttt{indicator} & \text{Makroökonomischer Indikator} \\
\texttt{value} & \text{Beobachteter Wert} \\
\texttt{unit} & \text{Einheit des Wertes} \\
\texttt{available\_from} & \text{Zeitpunkt, ab dem die Information verfügbar war}
\end{array}

{]}

Das zusätzliche Feld available\_from soll verhindern, dass ein
historischer Backtest einen Jahreswert bereits verwendet, bevor dieser
damals verfügbar gewesen wäre. Die genaue Regel dafür wird später bei
der BIP-Strategie festgelegt. Sie folgt dem gleichen Grundprinzip wie
die im Theorieteil behandelte Vermeidung von Look-ahead Bias.\footnote{Hilpisch,
  \emph{Python for Finance}, 319--21.}

5.4 Datenquellen und Umwandlung

5.4.1 Trennung zwischen Datenquelle und Modell

Nachdem das interne Format festgelegt ist, kann untersucht werden, wie
verschiedene Anbieter darauf abgebildet werden. Die Datenquelle soll
dabei nicht direkt von einer Strategie aufgerufen werden. Stattdessen
übernimmt ein Importmodul die Umwandlung.

{[}

\begin{array}{c}
\boxed{\text{Datenquelle}} \\[6pt]
\downarrow \\[6pt]
\boxed{\text{Import und Umwandlung}} \\[6pt]
\downarrow \\[6pt]
\boxed{\text{Einheitliches Datenformat}} \\[6pt]
\downarrow \\[6pt]
\boxed{\text{Simulation}}
\end{array}

{]}

Dadurch bleibt die eigentliche Berechnung gleich, auch wenn später eine
andere Datenquelle verwendet wird. Für meine Untersuchung sind Yahoo
Finance und FRED bereits in der Methodik als mögliche Quellen
vorgesehen.\footnote{Yahoo, {„Yahoo FinanzenDeutschland \textbar{} Börse
  live, Kurse, Nachrichten aus Wirtschafts- und Finanzwelt``}; Federal
  Reserve Bank of St. Louis, {„Federal {Reserve Economic Data}
  \textbar{} {FRED} \textbar{} {St}. {Louis Fed}``}.} Weitere Quellen
können später ergänzt werden, solange sie auf dasselbe interne Format
abgebildet werden können.

5.4.2 Marktdaten

Für Marktzeitreihen soll ein erster Adapter historische Daten eines
Anbieters laden und daraus die benötigten Felder erzeugen. Bei Yahoo
Finance beziehungsweise einer Python-Schnittstelle wie yfinance muss der
Adapter insbesondere festhalten:

\begin{itemize}
\tightlist
\item
  welches Wertpapier geladen wurde,
\item
  welcher Zeitraum verwendet wurde,
\item
  welche Datenfrequenz angefordert wurde,
\item
  welche ursprüngliche Spalte zu performance\_value wird,
\item
  welche Reihe für signal\_value verwendet wird,
\item
  und wie Dividenden beziehungsweise Splits in den Quelldaten behandelt
  werden.
\end{itemize}

Die technische Umbenennung einer Spalte reicht damit nicht aus. Zuerst
muss geklärt sein, was die ursprüngliche Reihe tatsächlich bedeutet. Für
konkrete Parameter einer später verwendeten API soll bei der
Implementierung zusätzlich deren direkte technische Dokumentation
verwendet werden.

5.4.3 Zins- und Makrodaten

FRED eignet sich als mögliche Quelle für Zinsreihen und andere
wirtschaftliche Daten. Für eine risikofreie Reihe ist besonders wichtig,
die Einheit der Quelle zu übernehmen. Ein als jährlicher Prozentsatz
veröffentlichter Zinssatz darf nicht unverändert von einer täglichen
Rendite abgezogen werden. Die Umrechnung muss deshalb erst erfolgen,
nachdem Frequenz und Einheit bekannt sind.

Für die spätere BIP-Gewichtung wird zusätzlich eine Quelle mit
historischen Länderwerten benötigt. Diese Daten werden zunächst in das
allgemeine Makroformat überführt. Welche konkrete Datenbank, welche
Länder und welche Veröffentlichungsverzögerung in der Hauptuntersuchung
verwendet werden, wird erst bei der endgültigen Versuchskonfiguration
festgelegt.

Damit bleibt an dieser Stelle noch offen, welche konkrete Serie später
verwendet wird. Festgelegt ist nur, welche Informationen das Modell
benötigt.

5.4.4 Import eigener CSV-Dateien

Damit das System auch ausserhalb meiner später ausgewählten Anbieter
verwendet werden kann, soll zusätzlich ein allgemeiner CSV-Import
möglich sein. Dabei kann eine fremde Datei über eine kleine Zuordnung
auf das interne Format abgebildet werden.

Beispiel einer Quelldatei:

{[}

\begin{array}{c|c|c}
\textbf{Date} & \textbf{Ticker} & \textbf{TR} \\
\hline
\texttt{2020-01-02} & \texttt{ABC} & 100.0
\end{array}

{]}

Dazu könnte folgende Zuordnung gehören:

{[}

\begin{aligned}
\texttt{Date}   &\;\longrightarrow\; \texttt{date} \\
\texttt{Ticker} &\;\longrightarrow\; \texttt{provider\_symbol} \\
\texttt{TR}     &\;\longrightarrow\; \texttt{performance\_value}
\end{aligned}

{]}

Der Import kann damit unterschiedliche Spaltennamen vereinheitlichen. Er
kann jedoch nicht automatisch beurteilen, ob die Werte fachlich für eine
bestimmte Untersuchung geeignet sind. Diese Prüfung bleibt Teil der
Datenwahl und Dokumentation.

5.5 Datenaufbereitung und Qualitätskontrolle

5.5.1 Rohdaten und verarbeitete Daten

Die ursprünglich geladenen Rohdaten und die später tatsächlich
verwendeten Daten werden getrennt gespeichert:

{[}

\begin{array}{l}
\texttt{data/} \\
\quad \texttt{├── raw/} \\
\quad \texttt{└── processed/}
\end{array}

{]}

Diese Trennung erleichtert die Nachvollziehbarkeit. Wenn sich eine Regel
bei der Aufbereitung ändert, kann sie erneut auf denselben Rohdaten
ausgeführt werden. Die Originaldatei muss dafür nicht manuell verändert
werden.

McKinney unterscheidet beim Arbeiten mit Daten ebenfalls zwischen
Einlesen, Bereinigen, Umformen und anschliessender Analyse.\footnote{McKinney,
  \emph{Python for Data Analysis}, Kap. 6--8.} Die konkrete
Ordnerstruktur ist an mein Projekt angepasst.

5.5.2 Validierung vor der Simulation

Bevor ein Datensatz in die Simulation gelangt, soll er automatisch
geprüft werden. Dabei sind mindestens folgende Punkte relevant:

\begin{itemize}
\tightlist
\item
  die benötigten Spalten sind vorhanden,
\item
  Datumswerte können korrekt eingelesen werden,
\item
  Wertspalten sind numerisch,
\item
  asset\_id ist vorhanden,
\item
  die Kombination aus Datum und Anlage kommt nicht doppelt vor,
\item
  fehlende oder unendliche Werte werden erkannt,
\item
  der verfügbare Zeitraum wird für jede Anlage bestimmt.
\end{itemize}

Fehlende Werte sind bei realen Datensätzen grundsätzlich möglich und
werden von pandas ausdrücklich als eigener Bestandteil der
Datenaufbereitung behandelt.\footnote{McKinney, Kap. 7.} Für meine
Simulation sollen solche Stellen deshalb sichtbar gemacht und nicht
unbemerkt weiterverarbeitet werden.

Das Ergebnis der Prüfung wird zusätzlich in einem kleinen Datenbericht
gespeichert. So lässt sich später nachvollziehen, ob eine Simulation mit
vollständigen Daten oder mit Warnungen durchgeführt wurde.

5.5.3 Umgang mit fehlenden Marktwerten

Fehlende Marktwerte werden in der ersten Version nicht automatisch
linear interpoliert. Eine solche Interpolation würde einen Wert
erzeugen, der in der ursprünglichen Datenquelle nicht beobachtet wurde.
Dadurch könnten auch künstliche Renditen oder Signale entstehen.

Nicht jede fehlende Kalenderbeobachtung ist allerdings ein Datenfehler.
An Wochenenden oder Feiertagen findet beispielsweise kein regulärer
Handel statt. Deshalb soll die Prüfung zwischen einem fehlenden
Handelstag und einer tatsächlichen Lücke in der vorhandenen Reihe
unterscheiden können.

Wenn mehrere Anlagen gemeinsam verwendet werden, muss zudem festgelegt
werden, wie unterschiedliche Handelstage ausgerichtet werden. Die genaue
Regel wird mit den später ausgewählten Datensätzen getestet. Wichtig ist
bereits jetzt, dass kein Wert ohne dokumentierte Regel ergänzt wird.

5.5.4 Datenfrequenz und Resampling

Die Theorie hat bereits gezeigt, dass Finanzzeitreihen in
unterschiedlichen Frequenzen vorliegen können. Für das Downsampling
einer Wertreihe wird der letzte verfügbare Wert eines abgeschlossenen
Intervalls verwendet. Die zeitlich korrekte Zuordnung ist wichtig, damit
keine zukünftige Information in eine frühere Periode gelangt.\footnote{Hilpisch,
  \emph{Python for Finance}, 319--21.}

Dabei muss zwischen Kurs- beziehungsweise Wertreihen und bereits
berechneten Renditen unterschieden werden. Wird beispielsweise aus
Tageswerten eine Monatsreihe erzeugt, kann der letzte Wert des Monats
verwendet werden. Mehrere Tagesrenditen eines Monats dürfen dagegen
nicht einfach durch die letzte Tagesrendite ersetzt werden; sie müssen
über ihre Wachstumsfaktoren miteinander verknüpft werden.

Das vorhandene resample\_dataframe() bildet damit einen Teil der
späteren Verarbeitung ab. Bei der Integration muss geprüft werden, für
welche Datentypen die jeweilige Aggregationsregel zulässig ist.

5.6 Veränderbare Einstellungen des Systems

5.6.1 Allgemeine Einstellungen

Das Modell soll nicht nur für einen einzigen fest eingebauten Versuch
funktionieren. Deshalb werden veränderbare Einstellungen ausserhalb der
Berechnungsfunktionen gespeichert.

Zu den allgemeinen Einstellungen gehören zunächst:

{[}

\begin{array}{l}
\textbf{Startdatum} \\
\textbf{Enddatum} \\
\textbf{Startkapital} \\
\textbf{Basiswährung} \\
\textbf{Datenfrequenz} \\
\textbf{Regel für fehlende Daten}
\end{array}

{]}

Diese Möglichkeiten des Systems sind noch nicht mit den späteren Werten
meiner Untersuchung gleichzusetzen. Das Modell kann beispielsweise
grundsätzlich verschiedene Zeiträume verarbeiten; für die Maturarbeit
wird später trotzdem ein bestimmter Zeitraum festgelegt und begründet.

5.6.2 Strategieabhängige Einstellungen

Auch die Strategien benötigen unterschiedliche Parameter.

Buy-and-Hold

{[}

\begin{array}{l}
\textbf{Verwendete Anlage} \\
\textbf{Startgewicht}
\end{array}

{]}

Rebalancing-Portfolio

{[}

\begin{array}{l}
\textbf{Verwendete Anlagen} \\
\textbf{Zielgewichte} \\
\textbf{Rebalancing-Rhythmus}
\end{array}

{]}

Damit muss die Engine technisch nicht auf genau 60/40 begrenzt sein. Die
Maturarbeit kann später trotzdem 60 \% Aktien und 40 \% Anleihen als
feste Versuchskonfiguration verwenden.

Trendfolge

{[}

\begin{array}{l}
\textbf{Verwendete Anlage} \\
\textbf{Kurzes SMA-Fenster} \\
\textbf{Langes SMA-Fenster} \\
\textbf{Long/Cash-Regel} \\
\textbf{Signalverzögerung} \\
\textbf{Behandlung der Cash-Phasen}
\end{array}

{]}

Die Signalverzögerung ist fachlich bereits durch die Vermeidung von
Look-ahead Bias vorgegeben.\footnote{Hilpisch, 694.} Die konkreten
SMA-Fenster werden dagegen später vor dem Hauptversuch begründet
festgelegt, damit die historisch beste Kombination nicht einfach
nachträglich ausgewählt wird.\footnote{Hilpisch, 697.}

5.6.3 Konfigurationsdatei

Für die erste Version sollen diese Einstellungen in einer JSON-Datei
gespeichert werden. Im Gegensatz zu den historischen Finanzdaten sind
die Einstellungen nicht einfach tabellarisch aufgebaut, sondern
teilweise verschachtelt. So enthält eine Strategie beispielsweise eigene
Parameter und beim Rebalancing wiederum mehrere Zielgewichte. JSON
eignet sich für solche hierarchischen Strukturen besser als CSV.
Zusätzlich kann Python JSON direkt einlesen und das Format lässt sich
später auch von einer Weboberfläche einfach erzeugen.

Ein vereinfachtes Beispiel sieht so aus:

\begin{Shaded}
\begin{Highlighting}[]
\ImportTok{import}\NormalTok{ json}

\NormalTok{beispiel\_config }\OperatorTok{=}\NormalTok{ \{}
    \StringTok{"period"}\NormalTok{: \{}
        \StringTok{"start"}\NormalTok{: }\StringTok{"2020{-}01{-}01"}\NormalTok{,}
        \StringTok{"end"}\NormalTok{: }\StringTok{"2020{-}12{-}31"}
\NormalTok{    \},}
    \StringTok{"start\_capital"}\NormalTok{: }\DecValTok{100000}\NormalTok{,}
    \StringTok{"base\_currency"}\NormalTok{: }\StringTok{"USD"}\NormalTok{,}
    \StringTok{"strategies"}\NormalTok{: \{}
        \StringTok{"buy\_hold"}\NormalTok{: \{}
            \StringTok{"enabled"}\NormalTok{: }\VariableTok{True}\NormalTok{,}
            \StringTok{"asset"}\NormalTok{: }\StringTok{"EQ\_DEMO"}
\NormalTok{        \},}
        \StringTok{"rebalance"}\NormalTok{: \{}
            \StringTok{"enabled"}\NormalTok{: }\VariableTok{True}\NormalTok{,}
            \StringTok{"target\_weights"}\NormalTok{: \{}
                \StringTok{"EQ\_DEMO"}\NormalTok{: }\FloatTok{0.60}\NormalTok{,}
                \StringTok{"BD\_DEMO"}\NormalTok{: }\FloatTok{0.40}
\NormalTok{            \},}
            \StringTok{"rebalance\_frequency"}\NormalTok{: }\StringTok{"annual"}
\NormalTok{        \},}
        \StringTok{"trend"}\NormalTok{: \{}
            \StringTok{"enabled"}\NormalTok{: }\VariableTok{True}\NormalTok{,}
            \StringTok{"asset"}\NormalTok{: }\StringTok{"EQ\_DEMO"}\NormalTok{,}
            \StringTok{"short\_window"}\NormalTok{: }\DecValTok{20}\NormalTok{,}
            \StringTok{"long\_window"}\NormalTok{: }\DecValTok{100}\NormalTok{,}
            \StringTok{"signal\_lag"}\NormalTok{: }\DecValTok{1}
\NormalTok{        \}}
\NormalTok{    \}}
\NormalTok{\}}

\BuiltInTok{print}\NormalTok{(json.dumps(beispiel\_config, indent}\OperatorTok{=}\DecValTok{2}\NormalTok{))}
\end{Highlighting}
\end{Shaded}

\begin{verbatim}
{
  "period": {
    "start": "2020-01-01",
    "end": "2020-12-31"
  },
  "start_capital": 100000,
  "base_currency": "USD",
  "strategies": {
    "buy_hold": {
      "enabled": true,
      "asset": "EQ_DEMO"
    },
    "rebalance": {
      "enabled": true,
      "target_weights": {
        "EQ_DEMO": 0.6,
        "BD_DEMO": 0.4
      },
      "rebalance_frequency": "annual"
    },
    "trend": {
      "enabled": true,
      "asset": "EQ_DEMO",
      "short_window": 20,
      "long_window": 100,
      "signal_lag": 1
    }
  }
}
\end{verbatim}

Die Zahlen in diesem Beispiel dienen nur dazu, die Struktur zu zeigen.
Sie sind noch keine endgültigen Parameter der Maturarbeit.

Der Vorteil der Konfigurationsdatei besteht darin, dass die
Berechnungslogik nicht geändert werden muss, wenn später ein anderer
Zeitraum oder ein anderer Parameter getestet wird. Eine grafische
Oberfläche müsste später im Wesentlichen dieselben Werte erfassen und an
die Engine übergeben.

5.7 Einheitliches Ergebnisformat

5.7.1 Gemeinsamer Output aller Strategien

Damit die Strategien später direkt verglichen werden können, sollen sie
einen gemeinsamen Grundoutput liefern. Für jede Strategie wird deshalb
eine vollständige Zeitreihe gespeichert:

{[}

\begin{array}{l}
\texttt{portfolio\_history.csv} \\[6pt]
\textbf{date} \\
\textbf{strategy} \\
\textbf{portfolio\_value} \\
\textbf{period\_return} \\
\textbf{drawdown}
\end{array}

{]}

Aus diesen Daten lassen sich anschliessend die bereits im Theorieteil
definierten Kennzahlen berechnen. Dazu gehören unter anderem
Gesamtrendite, annualisierte Rendite, Volatilität, Sharpe Ratio und
Maximum Drawdown.

Die Zeitreihe bleibt zusätzlich erhalten, weil ein einzelner Endwert
nicht zeigt, wie eine Strategie während des Untersuchungszeitraums
verlaufen ist. Gerade Drawdowns und unterschiedliche Marktphasen sollen
später Teil der Analyse werden.\footnote{Kommer, \emph{Souverän
  investieren mit Indexfonds und ETFs}, 7. Auflage (Weinheim: Campus
  Verlag, 2025), 143--66.}

5.7.2 Strategieabhängige Resultate

Neben dem gemeinsamen Grundoutput werden zusätzliche Informationen
gespeichert, wenn sie für eine Strategie relevant sind.

Für ein Rebalancing-Portfolio:

{[}

\begin{array}{l}
\texttt{weights\_history.csv} \\[6pt]
\texttt{date} \;|\; \texttt{strategy} \;|\; \texttt{asset\_id} \;|\; \texttt{weight\_before} \;|\; \texttt{target\_weight} \;|\; \texttt{weight\_after}
\end{array}

{]}

{[}

\begin{array}{l}
\texttt{trades.csv} \\[6pt]
\texttt{date} \;|\; \texttt{strategy} \;|\; \texttt{asset\_id} \;|\; \texttt{value\_before} \;|\; \texttt{target\_value} \;|\; \texttt{transaction\_value}
\end{array}

{]}

Für die Trendfolge:

{[}

\begin{array}{l}
\texttt{signals.csv} \\[6pt]
\texttt{date} \;|\; \texttt{strategy} \;|\; \texttt{asset\_id} \;|\; \texttt{signal} \;|\; \texttt{position} \;|\; \texttt{signal\_value} \;|\; \texttt{sma\_short} \;|\; \texttt{sma\_long}
\end{array}

{]}

Diese zusätzlichen Daten helfen später bei der Interpretation. Wenn eine
Strategie in einer bestimmten Phase besonders gut oder schlecht
abschneidet, kann überprüft werden, welche Gewichtung oder welches
Signal zu diesem Zeitpunkt vorlag.

5.7.3 Zusammenfassung und Run-Manifest

Für einen schnellen Vergleich wird zusätzlich eine Zusammenfassung
erzeugt. Sie enthält beispielsweise:

{[}

\begin{array}{l}
\texttt{strategy} \\
\texttt{start\_value} \\
\texttt{end\_value} \\
\texttt{total\_return} \\
\texttt{annualized\_return} \\
\texttt{annualized\_volatility} \\
\texttt{sharpe\_ratio} \\
\texttt{max\_drawdown}
\end{array}

{]}

Zusätzlich soll jeder Simulationslauf ein run\_manifest.json erhalten.
Darin werden die tatsächlich verwendete Konfiguration, der Datenstand,
der Zeitpunkt des Laufs und später auch die verwendete Git-Version
gespeichert.

Damit lässt sich ein Ergebnis auf die Kombination aus Daten, Parametern
und Code zurückführen. Das ist besonders wichtig, wenn später mehrere
Varianten des Modells getestet werden.

5.8 Grundgerüst der Simulationsengine

5.8.1 Aufteilung der Aufgaben

Die zentrale Simulationsfunktion soll nicht alle Berechnungen selbst
enthalten. Stattdessen werden die Aufgaben getrennt:

{[}

\begin{array}{rcl}
\texttt{data} & \longrightarrow & \text{Daten einlesen, vereinheitlichen und prüfen} \\[6pt]
\texttt{strategies} & \longrightarrow & \text{Regeln einer Anlagestrategie anwenden} \\[6pt]
\texttt{analysis} & \longrightarrow & \text{Gemeinsame Kennzahlen berechnen} \\[6pt]
\texttt{export} & \longrightarrow & \text{Resultate speichern}
\end{array}

{]}

Dadurch muss beispielsweise eine neue Datenquelle nicht wissen, wie die
Sharpe Ratio berechnet wird, und eine neue Strategie muss nicht selbst
CSV-Dateien laden.

Diese Trennung folgt vor allem aus dem Aufbau meines eigenen Systems.
Hilpisch zeigt ebenfalls, dass wiederverwendbarer Python-Code aus
interaktiven Beispielen in Funktionen und Module ausgelagert werden
kann.\footnote{Hilpisch, \emph{Python for Finance}.}

5.8.2 Vorgesehene Projektstruktur

{[}

\begin{array}{l}
\texttt{src/} \\
\texttt{├── funktionen.py} \\
\texttt{├── data/} \\
\quad \texttt{├── normalize.py} \\
\quad \texttt{├── validate.py} \\
\quad \texttt{└── adapters/} \\
\texttt{├── engine/} \\
\quad \texttt{├── config.py} \\
\quad \texttt{├── context.py} \\
\quad \texttt{├── result.py} \\
\quad \texttt{└── simulation.py} \\
\texttt{├── strategies/} \\
\quad \texttt{├── buy\_hold.py} \\
\quad \texttt{├── rebalance.py} \\
\quad \texttt{├── trend.py} \\
\quad \texttt{└── country\_weighting.py} \\
\texttt{├── analysis/} \\
\quad \texttt{└── metrics.py} \\
\texttt{└── export/} \\
\quad \texttt{└── results.py}
\end{array}

{]}

Die bereits vorhandene funktionen.py bleibt zunächst bestehen. Sie
enthält die mathematischen Bestandteile, die im Theorieteil entwickelt
wurden. Erst bei der eigentlichen Integration wird entschieden, ob
einzelne Funktionen einen Wrapper benötigen oder sinnvoll in ein anderes
Modul verschoben werden.

5.8.3 Gemeinsame Schnittstelle der Strategien

Damit die Engine unterschiedliche Strategien gleich behandeln kann,
sollen sie dieselbe äussere Schnittstelle erhalten.

Vereinfacht könnte eine Strategie später so aufgerufen werden:

result = strategy.run(context, params)

context enthält die für den Lauf vorbereiteten Markt-, Zins- und
Makrodaten sowie Zeitraum und Startkapital. params enthält nur die
Einstellungen der jeweiligen Strategie.

Das Ergebnis wird in einer gemeinsamen Struktur zurückgegeben. Die
Vermögensentwicklung ist dabei Pflicht; Gewichtungen, Transaktionen oder
Signale sind nur vorhanden, wenn die jeweilige Strategie sie benötigt.

Damit kann die spätere Auswertung alle Strategien über denselben
Grundaufbau vergleichen.

5.8.4 Ablauf eines Simulationslaufs

Ein vollständiger Lauf soll später in folgender Reihenfolge erfolgen:

{[}

\begin{array}{c}
\boxed{\text{Konfiguration laden}} \\[6pt]
\downarrow \\[6pt]
\boxed{\text{Daten laden und prüfen}} \\[6pt]
\downarrow \\[6pt]
\boxed{\text{Gewünschten Zeitraum vorbereiten}} \\[6pt]
\downarrow \\[6pt]
\boxed{\text{Aktivierte Strategien auswählen}} \\[6pt]
\downarrow \\[6pt]
\boxed{\text{Strategien ausführen}} \\[6pt]
\downarrow \\[6pt]
\boxed{\text{Gemeinsame Kennzahlen berechnen}} \\[6pt]
\downarrow \\[6pt]
\boxed{\text{Resultate und Manifest speichern}}
\end{array}

{]}

Während einer Simulation sollen keine neuen Marktdaten automatisch aus
dem Internet geladen werden. Dadurch bleibt klar, mit welchem Datenstand
ein bestimmtes Resultat entstanden ist.

5.9 Vorbereitung der technischen Implementierung

5.9.1 Verwendung der bestehenden Funktionen

Die Funktionen aus dem Theorieteil sollen nicht durch ein zweites,
unabhängiges Berechnungssystem ersetzt werden. Sie bilden die Grundlage
der späteren Engine.

{[}

\begin{array}{l|l}
\textbf{Vorhandene Funktion} & \textbf{Geplante Verwendung} \\
\hline
\texttt{resample\_dataframe()} & \text{Datenfrequenz anpassen} \\
\texttt{prozentuale\_aenderung()} & \text{Periodische Renditen} \\
\texttt{kumulierte\_rendite()} \,/\, \texttt{wachstumsfaktor()} & \text{Kumulierte Entwicklung} \\
\texttt{geometrisches\_mittel()} \,/\, \texttt{annualisierte\_rendite()} & \text{Langfristige Rendite} \\
\texttt{standardabweichung()} \,/\, \texttt{annualisierte\_volatilitaet()} & \text{Volatilität} \\
\texttt{drawdown()} \,/\, \texttt{maximum\_drawdown()} & \text{Verlustphasen} \\
\texttt{sharpe\_ratio()} & \text{Risikobereinigte Rendite} \\
\texttt{korrelationsmatrix()} \,/\, \texttt{portfolio\_risiko()} & \text{Portfolioanalyse} \\
\texttt{neue\_gewichtung()} \,/\, \texttt{rebalancing()} & \text{Rebalancing} \\
\texttt{buy\_and\_hold()} & \text{Passive Benchmark} \\
\texttt{trendfolge()} & \text{SMA-Strategie}
\end{array}

{]}

Vor der Integration muss geprüft werden, welche Datentypen diese
Funktionen erwarten und welche Resultate sie zurückgeben. Falls für die
neue Struktur zusätzliche Hilfsfunktionen fehlen, können sie ergänzt
werden. Die bereits im Theorieteil festgelegten mathematischen
Definitionen sollen dabei nicht still verändert werden.

5.9.2 Übergabe an Codex

Bis zu diesem Punkt sind Datenformat, Verarbeitung, Einstellungen,
Ergebnisstruktur und der grundsätzliche Ablauf des Systems beschrieben.
Die eigentliche technische Zusammenführung dieser Vorgaben mit den
vorhandenen Funktionen erfolgt im nächsten Schritt mit Unterstützung von
Codex.

Codex soll dabei nicht selbst festlegen, welche finanzwirtschaftliche
Methode verwendet wird. Vorgegeben werden:

\begin{itemize}
\tightlist
\item
  der Datenvertrag,
\item
  die Konfigurationsstruktur,
\item
  das Ergebnisformat,
\item
  die geplante Aufteilung des Systems,
\item
  die vorhandenen Funktionen aus funktionen.py,
\item
  sowie die bereits geschriebenen theoretischen Definitionen.
\end{itemize}

Der technische Auftrag besteht anschliessend darin, diese Teile zu einem
lauffähigen System zusammenzuführen und Tests mit künstlichen Daten zu
erstellen. Falls eine vorhandene Funktion angepasst werden muss, soll
die Änderung sichtbar dokumentiert werden.

5.9.3 Stand vor der Implementierung

Vor dem eigentlichen Bau des Systems sind damit die wichtigsten
Schnittstellen festgelegt. Noch nicht endgültig bestimmt sind bewusst
die konkreten Wertpapiere, der Untersuchungszeitraum, die risikofreie
Serie und die strategieabhängigen Hauptparameter. Diese Entscheidungen
werden nach der technischen Integration festgelegt und begründet, bevor
der eigentliche Hauptversuch durchgeführt wird.

Damit endet die Planung des allgemeinen Systems. Der nächste Schritt ist
die Umsetzung und Prüfung der Engine.

5.10 Festlegung der Untersuchungsdaten und Parameter

5.10.1 Von der Engine zur eigentlichen Untersuchung

Die bisher beschriebenen Bestandteile wurden zu einer allgemeinen
Backtesting-Engine zusammengeführt. Der für die Untersuchung vorgesehene
Stand ist Version 1.0.0 mit dem Git-Tag engine-v1.0. Die zugehörigen
technischen Prüfungen sind in der Entwicklungsdokumentation
festgehalten. Damit kann ich die Berechnung von der Auswahl der
konkreten Untersuchungsdaten trennen: Die Engine legt fest, wie ein Lauf
durchgeführt wird. Im folgenden Abschnitt lege ich fest, welche Daten
und Einstellungen ich für meine Arbeit verwenden möchte.

Die Forschungsfrage lautet: «Wie können historische Finanzdaten mithilfe
eines Python-basierten Modells verwendet werden, um langfristige
Investmentstrategien hinsichtlich Risiko und Rendite zu vergleichen?»
Die Untersuchung soll deshalb nicht nur einen möglichst hohen Endwert
finden. Sie soll zeigen, welche Unterschiede ein nachvollziehbarer
Vergleich sichtbar macht und wo seine Grenzen liegen. Die
Risikobetrachtung besteht dabei nicht aus einer einzigen Zahl. Kommer
erläutert, dass verschiedene Risikokennzahlen unterschiedliche Aspekte
erfassen. Für meine Arbeit verwende ich die bereits eingeführten Grössen
Volatilität, Sharpe Ratio und maximaler Drawdown neben der
Rendite.\footnote{Kommer, \emph{Souverän investieren mit Indexfonds und
  ETFs}, 7. Auflage (Weinheim: Campus Verlag, 2025), Abschn. 2.2.}

Ich lege die Hauptparameter vor der Auswertung fest. Andernfalls könnte
ich so lange andere Zeiträume oder SMA-Fenster ausprobieren, bis eine
gewünschte Strategie besonders gut abschneidet. Diese Gefahr wurde im
Theorieteil als Data Snooping beziehungsweise Overfitting
behandelt.\footnote{Kommer, \emph{Souverän investieren mit Indexfonds
  und ETFs}, 7. Auflage (Weinheim: Campus Verlag, 2025), Glossar «Data
  Snooping»; Hilpisch, \emph{Python for Finance}, Kap. 15,
  «Optimization».} Die folgenden Festlegungen sind meine eigenen
Untersuchungsentscheidungen. Die Quellen begründen die verwendeten
Konzepte, schreiben aber nicht genau diese Ticker, Zeiträume oder
Parameter vor.

5.10.2 Auswahl der Wertreihen

Für die praktische Umsetzung verwende ich historische ETF-Reihen statt
einzelner ausgewählter Unternehmen. Dadurch untersuche ich breit
angelegte Markt- beziehungsweise Länderpositionen. Die ETFs sind dabei
Stellvertreter für die jeweiligen Anlagebereiche. Ein ETF ist nicht
dasselbe wie ein kostenfreier theoretischer Index. Diese Unterscheidung
bleibt bei der Interpretation wichtig.

Für den globalen Aktienanteil wähle ich die US-Börsenlinie mit dem
Symbol VT. Vanguard beschreibt sie als weltweit ausgerichteten
Aktienfonds, der entwickelte Märkte und Schwellenländer abdeckt. Die
Anteilsklasse besteht seit Juni 2008.\footnote{\textbf{vanguardVTProdukt2026?}}
Für den Anleihenanteil verwende ich IEF. Dieser ETF bezieht sich auf
US-Staatsanleihen mit Restlaufzeiten zwischen sieben und zehn
Jahren.\footnote{\textbf{isharesIEFProdukt2026?}} Das passt zur in der
Methodik vorgesehenen Staatsanleihenkomponente. Es bedeutet jedoch
nicht, dass dieser Anteil risikofrei ist; insbesondere bleiben
Kursänderungen durch Zinsbewegungen Bestandteil seiner beobachteten
Wertentwicklung.

Für die Ländergewichtung beschränke ich mich auf sieben klar benannte
Länder. Die Auswahl umfasst die USA, Kanada, Japan, das Vereinigte
Königreich, Deutschland, Frankreich und Italien. Sie ist eine begrenzte
Fallstudie, keine vollständige Abbildung der Weltwirtschaft. Die
Datenprüfung muss die geplante historische Abdeckung bestätigen. Eine
Produktseite allein beweist noch nicht, dass der konkrete Download
vollständig und richtig ist.

{\def\LTcaptype{none} % do not increment counter
\begin{longtable}[]{@{}
  >{\raggedright\arraybackslash}p{(\linewidth - 6\tabcolsep) * \real{0.2500}}
  >{\raggedright\arraybackslash}p{(\linewidth - 6\tabcolsep) * \real{0.2500}}
  >{\raggedright\arraybackslash}p{(\linewidth - 6\tabcolsep) * \real{0.2500}}
  >{\raggedright\arraybackslash}p{(\linewidth - 6\tabcolsep) * \real{0.2500}}@{}}
\toprule\noalign{}
\begin{minipage}[b]{\linewidth}\raggedright
Interne Kennung
\end{minipage} & \begin{minipage}[b]{\linewidth}\raggedright
Yahoo-Symbol
\end{minipage} & \begin{minipage}[b]{\linewidth}\raggedright
Verwendung
\end{minipage} & \begin{minipage}[b]{\linewidth}\raggedright
Produktnachweis
\end{minipage} \\
\midrule\noalign{}
\endhead
\bottomrule\noalign{}
\endlastfoot
\emph{\texttt{WORLD\_EQ}} & \emph{\texttt{VT}} & Weltweiter Aktienanteil
& \footnote{\textbf{vanguardVTProdukt2026?}} \\
\emph{\texttt{US\_TREASURY\_7\_10}} & \emph{\texttt{IEF}} &
US-Staatsanleihen, 7--10 Jahre & \footnote{\textbf{isharesIEFProdukt2026?}} \\
\emph{\texttt{US\_EQ}} & \emph{\texttt{VTI}} & US-Aktienproxy und
längere US-Zusatzuntersuchung & \footnote{\textbf{vanguardVTIProdukt2026?}} \\
\emph{\texttt{CA\_EQ}} & \emph{\texttt{EWC}} & Kanada & \footnote{\textbf{isharesEWCProdukt2026?}} \\
\emph{\texttt{JP\_EQ}} & \emph{\texttt{EWJ}} & Japan & \footnote{\textbf{isharesEWJProdukt2026?}} \\
\emph{\texttt{GB\_EQ}} & \emph{\texttt{EWU}} & Vereinigtes Königreich &
\footnote{\textbf{isharesEWUProdukt2026?}} \\
\emph{\texttt{DE\_EQ}} & \emph{\texttt{EWG}} & Deutschland &
\footnote{\textbf{isharesEWGProdukt2026?}} \\
\emph{\texttt{FR\_EQ}} & \emph{\texttt{EWQ}} & Frankreich &
\footnote{\textbf{isharesEWQProdukt2026?}} \\
\emph{\texttt{IT\_EQ}} & \emph{\texttt{EWI}} & Italien & \footnote{\textbf{isharesEWIProdukt2026?}} \\
\end{longtable}
}

Ich verwende jeweils die US-notierte Linie und prüfe deren USD-Notierung
in den Metadaten. Die internen Kennungen unterscheiden die Rolle im
Modell von der Anbieterkennung. Die Länderfonds können unterschiedliche
Indexabgrenzungen und Begrenzungsregeln besitzen. Deshalb behaupte ich
nicht, dass alle Proxies sämtliche Unternehmen ihres Landes vollständig
oder nach genau derselben Indexmethode abbilden. Historische Produkt-
und Benchmarkänderungen werden im Datennachweis berücksichtigt.

5.10.3 Untersuchungsblöcke und Vergleichbarkeit

Der Hauptvergleich umfasst den Zeitraum Januar 2010 bis Dezember 2025.
Die anfängliche Bewertung erfolgt zum Monatsende Dezember 2009; danach
werden 192 Monatsrenditen ausgewertet. Das Ende 2025 schliesst das noch
laufende Jahr 2026 aus. Der Beginn ermöglicht eine längere gemeinsame
Untersuchung nach der Auflegung von VT und einen vorgelagerten
SMA-Anlauf. Dieser Zeitraum ist eine begründete Auswahl, nicht der
Nachweis, dass andere Zeiträume dieselben Resultate liefern würden.

Im Hauptlauf werden vier Strategien auf derselben vorbereiteten
Zeitachse ausgeführt: Buy-and-Hold mit VT, ein jährlich auf 60 \% VT und
40 \% IEF zurückgesetztes Portfolio, die Long/Cash-Trendfolge mit VT
sowie die BIP-Gewichtung der sieben Länderproxies. Damit werden
verschiedene Anlagekonzepte verglichen. Insbesondere besitzt die
BIP-Strategie ein anderes Aktienuniversum als VT. Ein Unterschied
zwischen beiden kann daher nicht allein der Gewichtungsmethode
zugeschrieben werden.

Für eine gezieltere Betrachtung ergänze ich einen zweiten Lauf: Die
BIP-Strategie wird mit einem jährlich gleichgewichteten Portfolio aus
genau denselben sieben Länder-ETFs verglichen. Jedes Land erhält dort
das Zielgewicht 1/7. Das entspricht einer zusätzlichen Konfiguration der
bestehenden Rebalancing-Strategie und keiner neuen Engine-Strategie. Bei
diesem Vergleich bleiben die Bausteine gleich, während sich die
Gewichtungsregel unterscheidet. Die Gleichgewichtung ist allerdings
keine Marktkapitalisierungsgewichtung. Der Zusatzlauf beantwortet
deshalb die Frage «BIP oder gleiche Ländergewichte?», nicht «BIP oder
Marktkapitalisierung?».

Die Literatur zeigt bereits, dass die Bewertung einer Gewichtungsmethode
vom betrachteten Zeitraum abhängen kann. Kommers Gegenüberstellung der
Ländergewichtungen liefert damit eine Grundlage für einen vorsichtigen
Vergleich, nicht die Vorhersage unseres eigenen Resultats.\footnote{Kommer,
  \emph{Souverän investieren mit Indexfonds und ETFs}, 7. Auflage
  (Weinheim: Campus Verlag, 2025), Abschn. 9.4.}

Als längere Ergänzung untersuche ich zudem Buy-and-Hold, 60/40 und
Trendfolge mit VTI beziehungsweise VTI/IEF von Januar 2003 bis Dezember
2025. Ausgangswert ist das Monatsende Dezember 2002. Dieser Block
umfasst 276 Monatsrenditen und geht weiter zurück als der globale
Hauptlauf. Er ist eine eigene US-Stichprobe. Ich verlängere VT nicht
rückwirkend durch eine andere Reihe und füge die beiden
Untersuchungsblöcke nicht zu einer scheinbar durchgehenden
Weltmarktgeschichte zusammen.

{\def\LTcaptype{none} % do not increment counter
\begin{longtable}[]{@{}
  >{\raggedright\arraybackslash}p{(\linewidth - 6\tabcolsep) * \real{0.2500}}
  >{\raggedright\arraybackslash}p{(\linewidth - 6\tabcolsep) * \real{0.2500}}
  >{\raggedright\arraybackslash}p{(\linewidth - 6\tabcolsep) * \real{0.2500}}
  >{\raggedright\arraybackslash}p{(\linewidth - 6\tabcolsep) * \real{0.2500}}@{}}
\toprule\noalign{}
\begin{minipage}[b]{\linewidth}\raggedright
Lauf
\end{minipage} & \begin{minipage}[b]{\linewidth}\raggedright
Zeitraum der Renditen
\end{minipage} & \begin{minipage}[b]{\linewidth}\raggedright
Vergleich
\end{minipage} & \begin{minipage}[b]{\linewidth}\raggedright
Funktion in der Arbeit
\end{minipage} \\
\midrule\noalign{}
\endhead
\bottomrule\noalign{}
\endlastfoot
\textbf{H1} & 2010--2025 & VT Buy-and-Hold; VT/IEF 60/40; VT Trend 3/12;
sieben Länder BIP-gewichtet & Hauptvergleich verschiedener Konzepte \\
\textbf{K1} & 2010--2025 & Gleiche sieben Länder: BIP versus jährliche
Gleichgewichtung & Gewichtungsvergleich bei identischen Bausteinen \\
\textbf{Z1} & 2003--2025 & VTI Buy-and-Hold; VTI/IEF 60/40; VTI Trend
3/12 & Längerer regionaler Zusatzvergleich \\
\textbf{S1 / S2} & 2010--2025 & VT Trend 2/12 beziehungsweise 6/12 &
Vorab festgelegte Sensitivität der Trendregel \\
\end{longtable}
}

Diese fünf Läufe sind der geplante Umfang. Die Sensitivitätsläufe
ersetzen den Hauptparameter nicht nachträglich. Auch ungünstige oder
widersprüchliche Resultate werden dargestellt. Die Läufe sind
historische Vergleiche, keine ausserhalb der Stichprobe durchgeführte
Prognoseprüfung.

5.10.4 Monatsfrequenz, Startkapital und Währung

Die Hauptauswertung verwendet Monatswerte. Das hält den Umfang
überschaubar und passt zum langfristigen Vergleich. Gleichzeitig ist
diese Frequenz eine Einschränkung: Bewegungen innerhalb eines Monats
werden nicht vollständig sichtbar. Der ausgewiesene maximale Drawdown
bezeichnet deshalb den stärksten Rückgang zwischen den verwendeten
Monatsbewertungen, nicht den grössten zwischenzeitlichen Tages- oder
Intradayverlust.

Die Rohdaten werden trotzdem täglich beschafft. Aus ihnen wird für jedes
Monatsende die letzte tatsächlich vorhandene gemeinsame
Handelstagsbewertung ausgewählt. Die Monatsreihe erhält das rechte
Kalender-Monatsendlabel. Hilpisch beschreibt dieses Vorgehen beim
Downsampling und warnt davor, den letzten Wert eines Intervalls nach
links an dessen Anfang zu beschriften.\footnote{Hilpisch, \emph{Python
  for Finance}, Kap. 8, «Resampling» und «Avoiding Foresight Bias».}

Für die Nachvollziehbarkeit speichere ich zusätzlich das tatsächliche
Quelldatum jeder Monatsbewertung. Das Monatslabel ist eine
Periodenbezeichnung, keine Behauptung eines Börsenhandels am Wochenende.
Es wird kein Kurs zwischen zwei Tagen interpoliert und kein fehlender
Monat aufgefüllt. Die Verarbeitung muss prüfen, dass am ausgewählten
Monatsende alle benötigten Reihen tatsächlich dieselbe Schlussbewertung
besitzen. Unvollständige Monate werden vor der Untersuchung geklärt und
nicht still entfernt.

Die Engine verlangt für ihre Jahreskennzahlen eine nachgewiesene
Periodenlogik. Für diese vorbereiteten Monatsdaten werden
period\_frequency=``ME'' und periods\_per\_year=12 gesetzt. Diese
Festlegung betrifft zwölf Monatsintervalle pro Jahr und keine
vermeintlich täglich handelbare Reihe. Die tatsächlichen Quelltage
bleiben im Aufbereitungsprotokoll sichtbar.

Jeder Lauf startet mit 100'000 USD und enthält keine weiteren
Einzahlungen oder Entnahmen. Der Betrag macht die Ergebnisse
anschaulich; bei den hier verwendeten proportionalen Portfolios ohne
feste Handelskosten beeinflusst er nicht die prozentuale
Strategierendite. Alle Ergebnisse werden nominal in USD betrachtet. Es
findet weder eine Inflationsbereinigung noch eine Umrechnung in
Schweizer Franken statt. USD-Notierung allein bedeutet zudem keine
Absicherung der im Fonds enthaltenen Fremdwährungsrisiken. Die
Untersuchung stellt somit nicht die persönliche Nettorendite eines
Schweizer Anlegers dar.

5.10.5 Dividenden, Kosten und Performance-Reihe

Für die Performance benötige ich eine Reihe, die Ausschüttungen
berücksichtigt. Ein reiner Kursindex genügt nicht, um den Gesamtertrag
einer Aktienanlage darzustellen.\footnote{Kommer, \emph{Souverän
  investieren mit Indexfonds und ETFs}, 7. Auflage (Weinheim: Campus
  Verlag, 2025), Abschnitt zu Wertpapierindizes.} Als Ausgangspunkt
verwende ich die Spalte Adj Close von Yahoo Finance. Yahoo beschreibt
diese als um Splits und Dividenden bereinigten Schlusskurs.\footnote{\textbf{yahooAdjustedClose2026?}}

Die umbenannte Spalte performance\_value wird deshalb als Anbieter-Proxy
für eine ausschüttungsbereinigte Entwicklung verwendet. Sie wird nicht
ohne Kontrolle als exakt identische Kopie einer offiziellen
Fonds-Total-Return-Reihe bezeichnet. Bei der Datenabnahme werden
Ausschüttungs- und Splitereignisse sowie verfügbare Jahresrenditen des
Anbieters stichprobenweise verglichen. Marktpreisrendite und NAV-Rendite
dürfen dabei nicht verwechselt werden. Dividenden werden nicht ein
zweites Mal auf die bereits bereinigte Reihe addiert.

Bei tatsächlichen ETF-Reihen können bereits fondsinterne Kosten in der
beobachteten Entwicklung enthalten sein. iShares erläutert dies
beispielsweise bei seinen Performanceangaben.\footnote{\textbf{isharesIEFProdukt2026?}}
Die bisherige Methodik ist deshalb zu präzisieren: Zusätzliche
Transaktionskosten, persönliche Steuern und weitere separat modellierte
Gebühren werden nicht berechnet. Bereits in der ETF-Reihe enthaltene
Kosten oder Quellensteuereffekte werden nicht rückgängig gemacht. Diese
Abgrenzung vermeidet die unzutreffende Behauptung einer vollständig
kosten- und steuerfreien Indexsimulation.

\bookmarksetup{startatroot}

\chapter{Fazit}\label{fazit}

\cleardoublepage
\phantomsection
\addcontentsline{toc}{part}{Anhang}
\appendix

\chapter{Dokumentation der
KI-Nutzung}\label{dokumentation-der-ki-nutzung}

Dieses Kapitel dokumentiert transparent, wie Codex bei der Erstellung
der Maturarbeit eingesetzt wurde. Die ausführlichen Arbeitsnachweise
werden projektintern im KI-Nutzungslog geführt. Hier wird der für die
Buchversion relevante aktuelle Stand zusammengefasst.

Die KI-Nutzungsdokumentation wurde von Codex formuliert und fortgeführt.
In den jeweiligen Prompts wurde Codex ausdrücklich gebeten, die
ausgeführten Arbeiten und die Art der KI-Nutzung direkt zu
dokumentieren. Der Autor hat diese Dokumentation gelesen und inhaltlich
überprüft, sie jedoch nicht selbst geschrieben.

\section{Grundsätze}\label{grundsuxe4tze}

Codex wurde vor allem für technische Einrichtung, Strukturierung,
Übertragung bereits vorhandener Inhalte und sprachliche Unterstützung
eingesetzt. KI-Ausgaben ersetzen keine fachliche Prüfung. Fachliche
Aussagen, Quellenqualität, mathematische Korrektheit und die
Interpretation der Resultate müssen vom Autor kontrolliert und
verantwortet werden.

Bei sprachlichen Korrekturen durfte Codex Rechtschreibung, Grammatik und
Satzbau verbessern, den Inhalt und den persönlichen Schreibton jedoch
grundsätzlich nicht verändern. Bei Überarbeitungen zur besseren
Lesbarkeit durften Formulierungen und Satzstrukturen vereinfacht werden.
Wichtige Fachbegriffe sollten dabei möglichst erhalten bleiben.

\section{Bisherige Einsätze}\label{bisherige-einsuxe4tze}

\subsection{Technische Einrichtung des
Projekts}\label{technische-einrichtung-des-projekts}

Am 11. Mai 2026 wurde Codex für die initiale technische Einrichtung des
Quarto-Book-Projekts verwendet. Dazu gehörten die Projektstruktur,
Kapitelgerüste, Bibliografie-Einbindung, Styles, Beispielnotebooks und
die Dokumentation der KI-Nutzung. Diese Unterstützung betraf das
technische Grundgerüst und nicht die endgültige fachliche Argumentation.

\subsection{Planung der Maturarbeit}\label{planung-der-maturarbeit}

Am 6. August 2026 übertrug und strukturierte Codex den bereits in der
README vorhandenen Zeitplan in der Todo-Anwendung. Die Planung wurde
anschliessend manuell geprüft.

\subsection{Übertragung der
Überblick-Inhalte}\label{uxfcbertragung-der-uxfcberblick-inhalte}

Am 7. August 2026 übertrug Codex alle bereits vorhandenen Inhalte aus
\texttt{README.md} in die Überblick-Seite des Quarto-Books. Codex
verfasste dabei weder die Daten noch die fachlichen Texte und ergänzte
keine neuen fachlichen Inhalte. Zusätzlich wurde ein wiederverwendbarer
Ablauf für spätere Aktualisierungen dokumentiert.

\subsection{Korrektur des Titels des
Theoriekapitels}\label{korrektur-des-titels-des-theoriekapitels}

Am 8. August 2026 untersuchte und korrigierte Codex die fehlerhafte
Anzeige des Titels von Kapitel 3. Dafür wurden die Quarto-Konfiguration
und der YAML-Frontmatter des Jupyter Notebooks angepasst. Der Fachtext,
die Berechnungen und die Quellen des Theoriekapitels wurden nicht
verändert.

\subsection{Ordnung der
Python-Dateien}\label{ordnung-der-python-dateien}

Am 8. August 2026 wurde Codex verwendet, um die lose Datei
\texttt{funktionen.py} sinnvoll im Repository einzuordnen. Sie wurde als
wiederverwendbarer Analysecode nach \texttt{src/funktionen.py}
verschoben. Der bereits vorhandene Funktionscode wurde bei dieser
organisatorischen Änderung nicht fachlich neu verfasst.

\subsection{Bündelung der
Analysefunktionen}\label{buxfcndelung-der-analysefunktionen}

Am 8. August 2026 übertrug Codex die nach \texttt{maximum\_drawdown()}
folgenden eigenen Funktionen aus dem Theoriekapitel in
\texttt{src/funktionen.py}. Beispielaufrufe und erklärende
Notebook-Kommentare wurden nicht übernommen. Ziel war eine einheitliche
Funktionssammlung für die spätere Analyse und den Vergleich der
Anlagestrategien. Die Vollständigkeit sowie grundlegende
Funktionsaufrufe wurden technisch geprüft.

\subsection{Rechtschreibung, Grammatik und
Satzbau}\label{rechtschreibung-grammatik-und-satzbau}

Codex wurde für die sprachliche Korrektur der Texte im Projekt
eingesetzt. Dabei wurden Rechtschreibung, Grammatik und Satzbau
kontrolliert. Inhalt, fachliche Aussagen und persönlicher Schreibton
durften grundsätzlich nicht verändert werden. Starke sprachliche
Eingriffe waren ausdrücklich nicht erlaubt. Die endgültige Kontrolle
bleibt beim Autor.

\subsection{Verbesserung der
Lesbarkeit}\label{verbesserung-der-lesbarkeit}

Codex wurde ausserdem verwendet, um Formulierungen in den Texten des
Repositorys allgemein verständlicher zu gestalten. Die Überarbeitung
orientierte sich an der Wiener Sachtextformel, dem Flesch-Index und dem
Gunning-Fog-Index. Komplexe Sätze und Formulierungen wurden vereinfacht,
wichtige Fachbegriffe jedoch möglichst beibehalten, damit Präzision und
Professionalität erhalten bleiben. Die fachlichen Aussagen wurden nicht
von Codex entwickelt.

\subsection{Erstellung der
KI-Nutzungsdokumentation}\label{erstellung-der-ki-nutzungsdokumentation}

Codex wurde in den jeweiligen Prompts ausdrücklich angewiesen, den
eigenen Einsatz und die ausgeführten Arbeiten zu dokumentieren. Codex
formulierte und aktualisierte deshalb das projektinterne KI-Nutzungslog,
dieses Kapitel der Buchversion und die wiederverwendbaren
Arbeitsanweisungen. Der Autor überprüfte die Dokumentation, schrieb sie
jedoch nicht selbst. Die Formulierungen der KI-Dokumentation stammen
somit von Codex; der Autor erteilte die Aufträge, kontrollierte die
Ergebnisse und übernahm die geprüfte Fassung.

\subsection{Einheitliche Formatierung}\label{einheitliche-formatierung}

Codex wurde mehrfach verwendet, um Dateien und Inhalte im gesamten
Repository einheitlich, übersichtlich und technisch korrekt zu
formatieren. Dies betraf unter anderem HTML-, Markdown-, QMD-, Quarto-
und Python-Strukturen sowie Überschriften, Listen, Tabellen, Codeblöcke,
Abstände und die Darstellung der Buchversion. Ziel war eine konsistente
und gut lesbare Gestaltung. Die Formatierung durfte fachliche Aussagen
und den persönlichen Schreibton nicht verändern. Die endgültige
Darstellung und die Einhaltung der formalen Vorgaben müssen vom Autor
kontrolliert werden.

\subsection{Einbindung des KI-Nachweises in die
Buchversion}\label{einbindung-des-ki-nachweises-in-die-buchversion}

Am 8. August 2026 wurde Codex verwendet, um den bisherigen
KI-Nutzungsnachweis als Anhang in das Quarto-Book einzubinden. Der
Inhalt dieses Kapitels ist eine für die Arbeit aufbereitete
Zusammenfassung des bestehenden KI-Nutzungslogs. Zusätzlich wurde ein
wiederverwendbarer Ablauf eingerichtet, mit dem zukünftige Logeinträge
in diese Buchversion übertragen werden können.

\section{Manuelle Verantwortung}\label{manuelle-verantwortung}

Der Autor muss sämtliche übernommenen Änderungen kontrollieren. Dazu
gehören insbesondere die unveränderte fachliche Bedeutung, die korrekte
Wiedergabe der Quellen, die mathematische Richtigkeit, die
Nachvollziehbarkeit des Codes und der Erhalt des persönlichen
Schreibtons. Die Verantwortung für die endgültige Fassung der Arbeit
bleibt beim Autor.

\protect\phantomsection\label{refs}
\begin{CSLReferences}{1}{1}
\bibitem[\citeproctext]{ref-admiralsWieSieMit2026}
Admirals. {„Wie Sie mit der Trendfolgestrategie bestehende Trends
nutzen``}. Admirals, 4. Juni 2026.
\url{https://admiralmarkets.com/de/wissen/articles/forex-strategy/trendfolgestrategie}.

\bibitem[\citeproctext]{ref-aroussiRanaroussiQuantstats2026}
Aroussi, Ran. \emph{Ranaroussi/Quantstats}. 1. Mai 2019, released 5.
August 2026. \url{https://github.com/ranaroussi/quantstats}.

\bibitem[\citeproctext]{ref-buchhaltung-einfach-sicherVolatilitatEinfachErklart2026}
Buchhaltung-Einfach-Sicher. {„Volatilität Einfach Erklärt \textbar{}
{Definition}, {Berechnung} Und {Nutzen}``}. Volatilität einfach erklärt
\textbar{} Definition, Berechnung und Nutzen, 7. Juni 2026.
\url{https://www.buchhaltung-einfach-sicher.de/finanzen/volatilitaet}.

\bibitem[\citeproctext]{ref-federalreservebankofst.louisFederalReserveEconomic2026}
Federal Reserve Bank of St. Louis. {„Federal {Reserve Economic Data}
\textbar{} {FRED} \textbar{} {St}. {Louis Fed}``}. Federal Reserve
Economic Data \textbar{} FRED \textbar{} St. Louis Fed, 7. Juni 2026.
\url{https://fred.stlouisfed.org/}.

\bibitem[\citeproctext]{ref-fidelityMaximumDrawdown2026}
Fidelity. {„Maximum {Drawdown}``}. FIL Limited 2026, 7. Juni 2026.
\url{https://www.fidelity.de/wissen/tipps-and-strategien/risiko-kennziffern/maximum-drawdown/}.

\bibitem[\citeproctext]{ref-fidelitySharpeRatio2026}
---------. {„Sharpe-{Ratio}``}. Sharpe-Ratio; FIL Limited 2026, 7. Juni
2026.
\url{https://www.fidelity.de/wissen/tipps-and-strategien/risiko-kennziffern/sharpe-ratio/}.

\bibitem[\citeproctext]{ref-gerdkommerInvestieren2026}
Gerd Kommer. {„Investieren``}. Gerd Kommer, 4. Juni 2026.
\url{https://gerd-kommer.de/investieren/}.

\bibitem[\citeproctext]{ref-hilpischPythonFinanceMastering2019}
Hilpisch, Yves. \emph{Python for Finance: Mastering Data-Driven
Finance}. Second edition. Beijing Boston Farnham Sebastopol Tokyo:
O'Reilly, 2019.

\bibitem[\citeproctext]{ref-kommerLeichteEinstiegWelt2025}
Kommer, Gerd. \emph{Der leichte Einstieg in die Welt der ETFs:
unkompliziert vorsorgen -- ein Starterbuch für Finanzanfänger}. 2.
vollständig überarbeitete und aktualisierte Neuausgabe. München: FBV,
FinanzBuch Verlag, 2025.

\bibitem[\citeproctext]{ref-kommerSouveraenInvestierenMit2025}
---------. \emph{Souverän investieren mit Indexfonds und ETFs: Ein
Investmentbuch für fortgeschrittene Privatanleger}. 7. Auflage.
Weinheim: Campus Verlag, 2025.

\bibitem[\citeproctext]{ref-markowitzPortfolioSelection1952}
Markowitz, Harry. {„Portfolio {Selection}``}. \emph{The Journal of
Finance} 7, Nr. 1 (1952): 77--91.
\url{https://doi.org/10.1111/j.1540-6261.1952.tb01525.x}.

\bibitem[\citeproctext]{ref-mckinneyPythonDataAnalysis2022}
McKinney, Wes. \emph{Python for Data Analysis: Data Wrangling with
Pandas, {NumPy}, and {Jupyter}}. Third edition. Beijing Boston Farnham
Sebastopol Tokyo: O'Reilly, 2022.

\bibitem[\citeproctext]{ref-PandasSeriespct_changePandas305}
{„Pandas.{Series}.Pct\_change --- Pandas 3.0.5 Documentation``}.
Zugegriffen 4. August 2026.
\url{https://pandas.pydata.org/pandas-docs/stable/reference/api/pandas.Series.pct_change.html\#pandas.Series.pct_change}.

\bibitem[\citeproctext]{ref-PandasSeriesstdPandas305}
{„Pandas.{Series}.Std --- Pandas 3.0.5 Documentation``}. Zugegriffen 7.
August 2026.
\url{https://pandas.pydata.org/docs/reference/api/pandas.Series.std.html}.

\bibitem[\citeproctext]{ref-phamPerformance60402025}
Pham, Nga, Bei Cui, und Ummul Ruthbah. \emph{The {Performance} of the
60/40 {Portfolio}: {A Historical Perspective}}. 18. Februar 2025.
\url{https://doi.org/10.56227/25.1.12}.

\bibitem[\citeproctext]{ref-Preise}
{„Preise``}. Zugegriffen 4. August 2026.
\url{https://www.bfs.admin.ch/content/bfs/de/home/statistiken/preise.html}.

\bibitem[\citeproctext]{ref-probstGleitendeDurchschnittBestimmen2025}
Probst, Wendelin. {„Der Gleitende Durchschnitt - so bestimmen Sie die
Trendrichtung``}. Online Broker LYNX, 21. November 2025.
\url{https://www.lynxbroker.de/boerse/trading/technische-analyse/technische-indikatoren/der-moving-average-so-bestimmen-sie-die-richtung-des-trends/}.

\bibitem[\citeproctext]{ref-ratgeberRelativeStrengthIndex2022}
Ratgeber, Redaktion. {„Relative Strength Index - so funktioniert der
RSI! - finanzen.ch``}. Ratgeber: Wissenswertes rund um die Geldanlage
\textbar{} finanzen.ch, 26. Oktober 2022.
\url{https://www.finanzen.ch/ratgeber/relative-strength-index/}.

\bibitem[\citeproctext]{ref-sharpeSharpeRatioFall2021}
Sharpe, William F. {„The {Sharpe Ratio} ({Fall} 1994)``}. In
\emph{Streetwise}, 169--78. Princeton University Press, 2021.
\url{https://www.degruyterbrill.com/document/doi/10.1515/9781400829408-022/html?lang=en}.

\bibitem[\citeproctext]{ref-vanguardStartseiteVanguardSwitzerland2026}
Vanguard. {„Startseite \textbar{} Vanguard Switzerland Private
Anleger``}. Startseite \textbar{} Vanguard Switzerland Private Anleger,
7. Juni 2026. \url{https://www.ch.vanguard/de/private-anleger}.

\bibitem[\citeproctext]{ref-yahooYahooFinanzenDeutschlandBoerse2026}
Yahoo. {„Yahoo FinanzenDeutschland \textbar{} Börse live, Kurse,
Nachrichten aus Wirtschafts- und Finanzwelt``}. Yahoo Finance, 7. Juni
2026. \url{https://de.finance.yahoo.com/}.

\end{CSLReferences}




\end{document}
